\documentclass[12pt,a4paper]{article}
\usepackage{graphicx}
\usepackage{slashed}
\usepackage{tikz}
\usepackage{colortbl}
\usepackage{calc}


% Macros-v2_188: re-based current paper identities and A--G family; no change to geometry.
% Macros-v2_184: active Paper III-E v1.200/III-F v1.267 binding for framework v3.031; mathematical macros unchanged.
% Macros-v2_183: updated shared cross-reference bindings for Paper1B/III-E/III-F v3.030; mathematical macros unchanged.
% Macros-v2_182: active Paper II-F/V-A cross-references for framework v3.029; definitions unchanged.
% Paper-series v2.381: new I-B torsion/modular results, bibliography and current source roots.
% Paper-series v2.380: new exceptional bibliography and synchronized paper references.
% Paper-series v2.379 current 31-paper macro/aux registry.
% Paper-series v2.378: current 31-paper AUX/bibliography registry updated for Paper I-B v1.2 literature additions.
% Framework v2.371: current-source reference registry refreshed; old theorem labels retained.
% Framework v2.370: current 30-paper source/AUX root rebinding; native mathematical macros unchanged.
% Framework v2.368: synchronized merged 30-root registry; mathematical notation unchanged.
% Framework v2.358: active 30-paper source registry synchronized; existing mathematical macros unchanged.
% Framework v2.357: selective new source-owned certificates; active reference roots updated for five revised papers.
% Framework v2.356: active 30-root registry synchronized to the 30 papers in this release.
% Framework v2.275: exact positive dense-twin word cubature, finite repeatability and eighth-moment classification, nested sampling frames, joint action and POVM-jet exports; preserves v2.274.
% Framework v2.274: filename and root-source synchronization.
% Framework v2.273: restored missing Problem 3 theorems from v2.268; exact exceptional frames, twin geometry, convolution repeatability and projector/occurrence firewalls; merged on v2.272 without overwriting independent branches.
% Framework v2.272: merged v2.270 covariant native-readout mathematics with v2.271 Problem 4 finite-action Hopf hypercharge/holonomy corrections; no universal bulk-gauge uniqueness claim.
% Independent branches: v2.270 readout/guard exports and v2.270 K16 action-consumed charge closure; all prior scoped assumptions retained.
% Framework v2.269: synchronized 30 active AUX roots; mathematics and bibliography preserved from Macros-v2_161.
% Framework v2.266: Paper II-C converse-derived outcome resources; II-E, III-C, IV-A/B, V-A and VI-B downstream mathematical closures; 30-root registry synchronized.
% Framework v2.265: II-C coordinate-stage constitution and exact first-symbol host inheritance; prior mathematical domains retained.
% Framework v2.264: paper II-B reverse/audit and IV-A observable-minimal exchange; synchronized external references to 30 active roots.
% Framework v2.263: Paper I-C adds bridge-controlled root-carry parity and source-relative old-cap cancellation; other root mathematical bodies retained; active 30-root references synchronized.
% Framework v2.262: integrates Paper II-F radius, Casimir and readout results; active 30-root references synchronized.
% Framework v2.261: merged Paper II-B GP2-level/geometric repairs; preserves v2.260 II-D repairs; active references synchronized.
% Framework v2.260: Paper II-D signed RS, finite-symbol typing, and conditional curvature/Jordan repairs; references synchronized.
% Framework v2.259: conditional native SU(2) atlas and explicit later discharge; v2.258 mathematical bodies retained elsewhere.
% Framework v2.258: preserves parallel II-A/dynamics update and adds complete readout/frequency closures.
% Framework v2.257: exact minimal complete position readout and arbitrary-preparation frequency law.
% Framework v2.256: exact native inertia and sharp bounded ordered-source closure; preserves concurrent v2.255 work.
% Framework v2.255: preserves parallel detector integration and adds exact finite readout closures.
% Framework v2.254: Pasterski bibliography and synchronized 30-paper external-reference registry.
% Framework v2.253: typed raw coordinate defect in I-A and synchronized VI-A import; active references synchronized.
% Framework v2.252: synchronized source registry; optional III-G shell-frame calculation.
% Framework v2.251: forward common-action execution, joint source jets and scoped coupled control.
% Framework v2.251: exact full-heat variations, fermion source jets and explicit physical completion frontier.
% Framework v2.251: certified interacting source jets, pole stability and scoped completion proofs.
% Framework v2.251: synchronized active source registry; mathematical notation unchanged.
% Framework v2.251: synchronized active source registry; mathematical notation unchanged.
% Framework v2.251: synchronized active source registry; mathematical notation unchanged.
% Framework v2.251: synchronized source registry; no mathematical notation changes.
% Framework v2.244: synchronized 30-paper AUX registry and rooted phase-source action-order tests.
% Framework v2.242: signed rooted phase lift, complete transverse return and exact family cubic.
% Framework v2.241: proved rooted conormal geometry, conditional cap realization and scoped action tests.
% Framework v2.239: all 30 active source roots synchronized; no mathematical macro changes.
% Framework v2.238: Paper IV A/B/C split and 30-root reference registry.
% Framework v2.237: Paper VII-D technical split and 28-root reference registry.
% Framework v2.236: Paper VI-C technical-volume reference added; all current root bindings advanced.
% Framework v2.235: 26-root AUX and bibliography bindings synchronized; mathematical macros unchanged.
% BEGIN INLINE SOURCE: Paper00-TeXBuildCompatibility-v1_0.tex
% Rendering-only repairs for inherited long appendices and bitmap text fonts.
% Integer counters through 26 retain the original labels; later values use AA,
% AB, ... . This file changes no action coefficient, state or theorem premise.
\RequirePackage{alphalph}
\makeatletter
\let\BHPrimitiveAlph\@Alph
\newalphalph{\BHExtendedAlph}{\BHPrimitiveAlph}{26}
\let\@Alph\BHExtendedAlph
\AtBeginDocument{%
  \@ifpackageloaded{microtype}{\microtypesetup{expansion=false}}{}%
}
\makeatother

% END INLINE SOURCE: Paper00-TeXBuildCompatibility-v1_0.tex

% Framework v2.164: current reference/bibliography profile; mathematical definitions retained.
% Framework v2.162: finite metric, functional Hilbert and scoped Ward integration.
% Framework v2.161: inherits Macros-v2_57.tex exactly except the current external-reference binding.
% Macros v2.57 / framework v2.156: current reference registry only; mathematical definitions unchanged.
% Framework v2.143 / Macros v2.55: current external-reference bindings; mathematical notation unchanged.
% Framework v2.110 / Macros v2.54: current external-reference registry and bibliography; mathematical notation inherited.
% Framework v2.67 / Macros-v2_53: adds Born-action-cell and joint-EH readout symbols; absolute unit realization is metrological; MATCH_{J/scal}^{EH} remains the sole G24 normalization bridge.
% Framework v2.35 / Macros-v2_52: grouped paper numbering, active AUX registry, and Born-v6_20 bibliography binding synchronized; notation unchanged.
% Framework v2.26: active-source AUX and bibliography bindings synchronized; notation unchanged.
% Framework v2.09: reviewed early-reciprocity integration; scoped physical imports.
% Framework v2.08: Born--Euler mathematics and conditional complete-action potential criterion.
% Framework v2.07: active AUX registry includes Paper I-B v1.2; shared bibliography Born-v6_18.
% Macros-v2_46 / framework v2.04: mathematical atlas and downstream realization correction.
% Macros-v2_46 / framework v2.03: scoped structure integration and current dependency bindings.
% Macros-v2_46 / framework v2.01: shared source bindings; prior revision comments are historical.
% v2.42 / framework v1.99: Synchronized RG/AS owner and bibliography bindings.
% v2.41 / framework v1.99: active manuscript AUX bindings and geometric-closure bibliography synchronized.
% Release-v1.80 centralized source synchronization for scalar-circle/comparison-band integration.
% Release v1.79: Hurwitz hierarchy and kinetic-normalized action response integration; physical mass matching remains conditional.
% Release v1.78: approved M8 closure-batch integration; physical and analytic domains retained.
% Current v1.78 registry: all active TeX roots are synchronized below.
% Earlier release comments are historical scope records, not current status.
% Release-v1.77 synchronization: native same-owner Noether closure, generator--flow/history typing, analytic-certification firewalls, scalaron/K16, LF/QCD, quantum-Ward, and RAR boundary-owner transport are registered; shared mathematical notation is unchanged.
% Release-v1.76 synchronization: registry points to the v1.76 active paper roots.
% Release-v1.76 synchronization: registers the local A+Phi+g Noether/constraint closure and dynamics-frontier propagation; shared mathematical notation is unchanged.
% Release-v1.75 synchronization: registry points to the v1.75 active paper roots.
% Release-v1.75 synchronization: integrates the joint-Dirac Einstein bulk real form and the local odd-scalar heat-coefficient action on explicit analytic first-jet/interior domains; v1.74 Yang--Mills closure is retained. Fermionic reality, pure-metric curvature-squared/exact-nonlocal terms, noncanonical boundaries, physical vacuum selection, and actual solution-history tokens remain open.
% Release-v1.74 synchronization: registry points to the v1.74 active paper roots and retains the v1.73 theorem labels.
% Release-v1.74 synchronization: canonical Macros-v2_34.tex; generated local nonlinear Yang--Mills action, constrained connection-field phase space, source-free Gauss momentum map, and Legendre closure are integrated on their proved domains. The remaining non-gauge ACT_cl sectors and actual solution-history selection remain open.
% Release-v1.71 synchronization: closes the former JOINT_x/JOINT_5
% connection-occurrence gap by canonical structure-group extension of the
% nonflat Born connection, faithful associated representation, and the local
% cotangent--Clifford operator constructor.  Actual fluctuation members,
% actions, vacua, solutions, calibrations, and analytic QFT domains remain
% separately typed.  The active-AUX registry is advanced once.
% Release-v1.71 canonical macro and active-AUX registry.
% ======================================================================
% Macros-v2_39.tex -- canonical publication macro file
% Release v1.71 retains the v1.70 null-ray/QCD boundaries and additionally
% removes the duplicate structural connection actuality premise.  It does not
% claim global bundle topology, select a gauge/Higgs member, or promote the
% finite construction to an interacting continuum theory.
% Release v1.68 constructs the finite light-front total-fraction operators and
% parton number/momentum operator measures, proves their simplex and Hausdorff-limit
% theorems, and separates positive cut measures from scheme-renormalized
% light-ray distributions.  It also separates finite real Schur poles, stable
% norm-resolvent/Riesz pole-residue limits, and continued resonance
% pole/residue limits.  The Paper-VII-C scattering package is claim-indexed:
% local QFT and spectral claims import no S-matrix or infrared premise, while
% asymptotic amplitudes, LSZ, S-matrix poles, completeness, and cross sections
% import only their named massive or massless/long-range clauses.
% Release v1.67 factorizes BND48_H into a same-action Hessian/Schur
% descendant and still-open typed boundary-trace, critical-lift, and
% refinement-limit conditions.  It proves the finite-collar Dirichlet-to-
% Neumann law sqrt(y)tanh(L sqrt(y)), its nonuniform critical limit, and the
% finite-return ultra-infrared turnover.  The Paper-I-A Hardy constant mode is
% an exact geometric near-hit, not yet the physical metric boundary mode.
% The current local f(R) Hessian remains massive and therefore still refutes
% a RAR prediction on the minimal action branch.
% Release v1.66 derives the transverse depth-six normalizer bridge from the
% already generated C6 Reynolds image and Paper-I-A binary-dihedral normalizer.
% The bridge is the unique least 2D3-stable scalar-containing realization of
% the faithful quotient 2D3/C6=C2, has trace defect six and hence K=48, but is
% transverse to the nested active carrier.  Exact K=44 and K=52 contraction,
% semigroup, and energy-midpoint identities are registered as structural
% readouts rather than projector events.  The active E48 carrier gap remains
% unchanged.  A same-action boundary-symbol/Calderon intertwiner is isolated
% as the necessary and still unproved owner bridge to the static DtN sort;
% representation equivalence alone neither derives RAR nor creates a K=48
% action junction, pole, occupation, or cosmological transition.
% Release v1.65 derives the native E_48-centered carrier mass gap from the
% GP1--GP3-generated active spectrum, distinguishes its carrier-free interval
% from a particle located at E_48, and withdraws the unsupported promotion to
% a positive continuum Yang--Mills mass gap.  The latter remains a valid
% conditional transfer theorem on an enlarged continuum branch, but it is not
% an active framework prediction or an unfinished native-gap proof.
% Release v1.64 synchronizes Derivation Batches 14--22.  It registers the
% GP1--GP3 heat-profile orbit, same-Hessian electroweak response and pole
% ownership, K16/K48 action-response corrections, perturbation state and
% non-Gaussianity typing, finite-vacuum/gap/QCD limit repairs, factorized
% position-area ownership, the actual static metric response, and the ordered
% Wilson/string-tension limits.  Strict continuum convergence, positive
% physical gap, surviving higher cumulants, and positive tension remain open.
% Release v1.61 closes TIMEWIDTH16 as an independent direct-production datum.
% JUNC_{2|7} continues to own the finite intrinsic microscopic profile/width,
% while the actual same-action background/production history determines its
% physical-time pullback, support, and duration.  The current Paper-I-A--VII-C AUX
% registry is advanced once; shared mathematical notation is unchanged.
% Release v1.60 derives the final-octet primitive-remnant depletion, canonical
% Hopf-current representative, exact Hopf descent to dOmega/(4pi), and
% closed-FLRW a^{-3} neutral-current shape.  It retypes K=48 as the onset
% of the final one-unit geometric transfer and records cDM/RAR as distinct
% physical projections of that common neutral parent while preserving
% occurrence, Hilbert-owner, stress-matching, and static-kernel premises.
% Release v1.59 transports the v1.58 exceptional-shell/K16 closure through
% Papers III-F and VI-B--VII-C.  It replaces Paper VI-B's obsolete mixed-Hessian
% output typing by the third-variation pair vertex and Bogoliubov occupation
% support, derives the post-seed Higgs-pair feedback vertex and its exact
% alpha--delta kernel factors, propagates the unique K16/K48 action-junction
% split, identifies the QCD delta count with the upstream 2O shell atom only
% as a common-scalar numerical compatibility, and narrows THERM16 to dynamical
% approach/beta selection because the finite Gibbs/KMS target is already
% unique on a supplied simple observable algebra once its Hamiltonian and
% inverse temperature are supplied; direct-sum central weights remain typed.
% Release v1.58 registers the exceptional nearest-shell closure, the derived
% transverse K=16 action event, the unique interior action junction, the
% type-correct ordered pair vertex, the thin-layer squeeze law, and the
% radiation-like/matter-like kinematic split.  REG_0, transfer magnitude,
% physical layer width, and thermalization remain explicitly conditional.
% Release v1.57 registers the active-depth-six obstruction, the native
% K=48 carrier desert, the typed 7|9 owner reduction, and the dynamical
% reheating handoff across Papers III-B--VI-B.
% Release v1.56 registers the least-support n_2 inflation dressing of the
% persistent n_7 Laurent response, the analytic global-convexity and positive
% endpoint-measure proofs, the derived inflation core and K=16 coefficient
% freeze, and the no-particle dark-sector split into static RAR plus the
% conditional homogeneous effective-curvature response at reflected K=48.
% There is no n_6 action, scalaron-DM mass, or second Einstein--Hilbert term.
% Release v1.55 registers the protected-octet resolved/residual reflection
% q_sp+Z_W=8, its j<->8-j and K<->64-K scale covariance, and the derived
% rank-one-exit quadratic grade C_1^2/U=C_2 in Papers III-C, III-E, and VIII.
% The reflection is a typed trace/scale theorem, not a time reversal or a
% physical event-occurrence theorem; reheating transfer remains conditional.
% Release v1.54 registers the exact positive-part Laurent indices, the
% active-domain/terminal-endpoint action split, the no-second-EH firewall, and
% the q_DM=6/K_DM=48 scalaron readout typing across Papers III-D/E/F, VI,
% VIII, and IX.
% Release v1.53 registers the common nonnegative Laurent obstacle projection
% for n_1^+ and n_7^+, discharges the three former exit/freeze assumptions,
% and advances the Paper-III-D/E/F, VIII, and IX consumers.  Shared notation
% and bibliography ownership are unchanged.
% Release v1.51 registers the PQTI derivation batch and its updated
% Paper-III-A/C/D/E/F, VIII, and IX consumers.  Shared notation and
% bibliography ownership are unchanged.
% Release v1.50 advances the Paper-III-B AUX target to v1.40 after the
% rigorous audit repair; shared notation and bibliography ownership are unchanged.
% Release v1.49 advances the Paper-II-E/F AUX target to v3.14 after the
% minimal audit repair; shared notation and bibliography ownership are unchanged.
% Release v1.46 synchronizes Paper II-A's minimal frozen constructor-origin
% theorem and its five dependent AUX targets.  Shared notation is unchanged.
% Release v1.45 updates the Paper-II-B AUX target to v2.23 after the
% GP3-rigidity and appendix-placement repair; shared notation is unchanged.
% Release v1.44 advances the centralized Paper-I-A--VII-C AUX registry after the
% Paper-I-A v5.16 protected-completion scope repair.  Every current
% paper imports this one file; no paper-level opt-in or \externaldocument list
% is permitted.  The bibliography remains centralized here through Born-v6_13.
% Shared bibliography, theorem infrastructure, colors, notation alignment,
% and current Paper-I-A--VII-C external-reference registration.
% This file centralizes common imports, bibliography sorting, colors, theorem
% infrastructure, and the notation shared by the complete series.
% v2.9 (release v1.43) made the Paper-I-A--VII-C AUX registry automatic and
% centralized the bibliography and all external-reference ownership.
% v2.8 release-registry refresh (v1.33) registers the rooted generated-parent
% cascade and the POL/REG/free-law/11+4/Born-group closure revisions.  No
% shared mathematical notation or Paper-II-A fundamental principle is changed.
% v2.8 release-registry refresh (v1.32) registers the universal generated-
% continuation revisions in Papers II-A, III-A, III-C, IV, and VIII.  No
% shared mathematical notation or Paper-II-A fundamental principle is changed.
% v2.8 release-registry refresh (v1.30) registers the CRP4/projective-state
% source revisions and their direct registry-only consumers.  No shared
% mathematical notation or Paper-II-A axiom is changed.
% v2.7 release-registry refresh (v1.27) updates only the current external
% Paper-I-A--VII-C basenames after the placement/assumption synchronization; no
% shared mathematical notation or Paper-II-A axiom is changed.
% v2.7 synchronizes the geometric defect-repair revision: Paper-II-A v3.42,
% Paper-III-C v1.26, Paper-III-D v1.27, Paper-III-E v1.32, and Paper-III-F
% v1.34.  It registers the common trace-defect/spectral-filtration theorem,
% its owner-typed response pullback, and the reduced GRED and TERM interfaces.
% No shared mathematical notation is changed.
% v2.6 synchronizes Paper-III-E v1.31's cone-type repair, minimal-response
% classification, and explicit classical action-role boundary, together with
% Paper-III-F v1.32's downstream response-atlas synchronization.  No shared
% mathematical notation is changed.
% v2.5 synchronizes the Paper-III-D v1.26 faithful-trace-duality proof,
% exact ordered-owner joint locus, and typed alternative-presentation
% bookkeeping repair, together with its Paper-III-E and closure-frontier
% references.  No shared mathematical notation is changed.
% v2.4 synchronizes the Paper-III-C v1.25 orbit/process/owner repair, its
% downstream Papers III-B--F, V, VI, VIII, and IX revisions, and Born-v6_13.
% v2.3 synchronizes the Paper-II-A v3.41 witness/support/principal inheritance
% layer, the Paper-III-A v1.24 joint-realizability/readout firewall, all dependent
% host-rate revisions, and the current Paper-I-A--VII-B source registry.
% v2.2 synchronizes the generated active-realization, determinant-depth,
% dimension-weighted hierarchy, and typed cosmological-event interfaces.
% v2.1 synchronizes the exceptional-support interface revisions: Paper I-A
% v5.9, Paper II-A v3.30, Paper II-E/F v3.7, and Paper IV v4.16.
% v2.0 introduces the stage-safe G0 occurrence notation, synchronizes the
% active Paper-I-A--IV source registry (current rigorous Paper-III closure:
% III-A v1.17, III-B v1.22, III-C--E v1.17, III-F v1.14, IV v4.14), and makes
% external-reference self-skip depend on logical paper role rather than a
% versioned job name.
% ======================================================================

% ======================================================================
% macros.tex — Shared macros for the Born Geometry series
% ======================================================================

% ======================================================================
% 1. PACKAGES, BIBLIOGRAPHY, COLORS, AND UNIT DEFINITIONS
% ======================================================================
% Core packages shared across the complete Paper-I-A--VII-C series.  Loading them here makes
% the macro file a genuine series preamble while remaining compatible with existing
% paper preambles that load the same packages before \input.
\usepackage{amsmath}
\usepackage{amssymb}
\usepackage{amsthm}
\usepackage{mathrsfs}
\usepackage{array}
\usepackage{booktabs}
\usepackage{xcolor}
\makeatletter
\@ifpackageloaded{geometry}{}{\usepackage{geometry}\geometry{margin=1in}}
\makeatother
\usepackage{mdframed}

% Shared support packages required by the macros below and by most papers.
\usepackage{ragged2e}
\usepackage{longtable}
\usepackage{tabularx}
\usepackage{enumitem}
\usepackage[most]{tcolorbox}

% Series-wide numeric bibliography convention using the portable BibTeX path.
% All Papers I-A--VII-C use the ordinary \cite interface.
\usepackage[numbers,sort&compress]{natbib}
\providecommand{\printbibliography}{%
  \bibliographystyle{plainnat}%
  \bibliography{Born-v6_59}%
}

% Shared colors.  \providecolor avoids clobbering a paper-specific override.
\providecolor{postulatecolor}{RGB}{230,242,255}
\providecolor{compare}{RGB}{230,245,230}
\providecolor{definitioncolor}{RGB}{230,242,255}

\providecolor{openriskcolor}{RGB}{255,235,235}
\providecolor{bridgecolor}{RGB}{240,240,255}
\providecolor{gold}{RGB}{255,215,0}

% ======================================================================
% 2. CUSTOM ENVIRONMENTS
% ======================================================================
% Common table column types used in all papers.  Guarded so that a paper may
% define its own versions before loading this file.
\makeatletter
\@ifundefined{NC@find@L}{%
  \newcolumntype{L}[1]{>{\RaggedRight\arraybackslash}p{#1}}%
}{}
\@ifundefined{NC@find@C}{%
  \newcolumntype{C}[1]{>{\Centering\arraybackslash}p{#1}}%
}{}
\@ifundefined{NC@find@Y}{%
  \newcolumntype{Y}{>{\RaggedRight\arraybackslash}X}%
}{}
\makeatother

% ======================================================================
% 3. THEOREM STYLES AND ENVIRONMENTS
% ======================================================================
% Custom underlined theorem style (useful for remarks or notes)
\newtheoremstyle{underlined}
  {3pt}   % Space above
  {3pt}   % Space below
  {}      % Body font (empty means normal/roman text, typical for remarks)
  {}      % Indent amount
  {}      % Theorem head font (can be \itshape, \bfseries, etc.)
  {.}     % Punctuation after theorem head
  { }     % Space after theorem head
  {\underline{\thmname{#1}\thmnumber{ #2}}\thmnote{ (#3)}} % Theorem head formatting

% Standard plain theorem style environments.  Each environment is guarded so
% the shared macro file can be loaded after a paper-specific theorem setup.
\theoremstyle{plain}
\makeatletter
\@ifundefined{theorem}{\newtheorem{theorem}{Theorem}[section]}{}
\@ifundefined{lemma}{\newtheorem{lemma}{Lemma}[section]}{}
\@ifundefined{proposition}{\newtheorem{proposition}{Proposition}[section]}{}
\@ifundefined{corollary}{\newtheorem{corollary}{Corollary}[section]}{}
\makeatother

% Definition style environments (typically non-italicized body)
\theoremstyle{definition}
\makeatletter
\@ifundefined{definition}{\newtheorem{definition}{Definition}[section]}{}
\@ifundefined{postulate}{\newtheorem{postulate}{Postulate}[section]}{}
\@ifundefined{closure}{\newtheorem{closure}{Closure}[section]}{}
\makeatother

% Applying the custom underlined style
\theoremstyle{underlined}
\makeatletter
\@ifundefined{remark}{\newtheorem{remark}{Remark}[section]}{}
\@ifundefined{convention}{\newtheorem{convention}{Convention}[section]}{}
\@ifundefined{principle}{\newtheorem{principle}{Principle}[section]}{}
\makeatother

% Framed proof environment.
% Every \begin{proof} ... \end{proof} is automatically framed/shaded.
% This avoids mdframed/trivlist nesting, which can drop the top border after
% theorem-like environments.
\makeatletter
\renewenvironment{proof}[1][\proofname]{%
  \par\addvspace{6pt}%
  \noindent
  \begin{tcolorbox}[%
    enhanced,
    breakable,
    colback=black!4,
    colframe=black!35,
    boxrule=0.35pt,
    arc=1pt,
    left=6pt,
    right=6pt,
    top=4pt,
    bottom=4pt,
    boxsep=0pt,
    before skip=0pt,
    after skip=6pt,
    parbox=false
  ]%
  \pushQED{\qed}%
  \normalfont\itshape #1\@addpunct{.}\enspace
  \normalfont\ignorespaces
}{%
  \unskip\nobreak\hfill\popQED
  \end{tcolorbox}%
  \par\addvspace{2pt}%
}
\makeatother

\mdfdefinestyle{postulatebox}{
  backgroundcolor=yellow!10,
  linecolor=black,
  linewidth=0.4pt,
  skipabove=\baselineskip,
  skipbelow=\baselineskip,
  nobreak=true
}

% ======================================================================
% PROOF-READABILITY HELPERS
% ======================================================================

\newcommand{\ProofPlan}[1]{%
  \par\smallskip
  \noindent\textbf{Proof plan. }#1%
  \par\smallskip
}

\newlist{proofsteps}{enumerate}{2}
\setlist[proofsteps,1]{
  label=\textbf{(\alph*)},
  ref=\alph*,
  leftmargin=2.2em,
  itemsep=0.45\baselineskip,
  topsep=0.35\baselineskip
}
\setlist[proofsteps,2]{
  label=\textbf{(\roman*)},
  ref=\roman*,
  leftmargin=2.0em,
  itemsep=0.25\baselineskip,
  topsep=0.2\baselineskip
}


% ======================================================================
% 4. NUMBER SYSTEMS AND ALPHABETS
% ======================================================================
\providecommand{\RR}{\mathbb{R}} % Real numbers
\providecommand{\PP}{\mathbb{P}} % Projective space / Probability
\providecommand{\CC}{\mathbb{C}} % Complex numbers
\providecommand{\ZZ}{\mathbb{Z}} % Integers
\providecommand{\NN}{\mathbb{N}} % Natural numbers
\providecommand{\HH}{\mathbb{H}} % Quaternions / Hyperbolic space
\providecommand{\LFRayGeom}{\mathsf{LF\!\mbox{-}RAY\!\mbox{-}GEOM}}
\providecommand{\LFRep}{\mathsf{LF\!\mbox{-}REP}}
\providecommand{\LFState}{\mathsf{LF\!\mbox{-}STATE}}
\providecommand{\LFLim}{\mathsf{LF\!\mbox{-}LIM}}
\providecommand{\LFScat}{\mathsf{LF\!\mbox{-}SCAT}}
% Number systems / projective spaces
\providecommand{\CP}{\mathbb{CP}}
\providecommand{\RP}{\mathbb{RP}}

% ======================================================================
% 5. GROUPS AND LIE THEORY
% ======================================================================
\providecommand{\SO}{\ensuremath{\mathrm{SO}}}     % Special Orthogonal group
\providecommand{\SU}{\ensuremath{\mathrm{SU}}}     % Special Unitary group
\providecommand{\Spin}{\ensuremath{\mathrm{Spin}}} % Spin group
\providecommand{\CS}{\ensuremath{\mathrm{CS}}}     % Chern-Simons (capitalized)
\providecommand{\cs}{\ensuremath{\mathrm{cs}}}     % Chern-Simons (lowercase)

% ======================================================================
% 6. COMMON MANIFOLDS AND TOPOLOGY
% ======================================================================
\providecommand{\Sone}{S^1}                 % Standard circle

% ======================================================================
% 7. MATH OPERATORS (Declared via \DeclareMathOperator)
% ======================================================================
\DeclareMathOperator{\End}{End}   % Endomorphisms
\DeclareMathOperator{\Aut}{Aut}   % Automorphisms
\DeclareMathOperator{\spec}{spec} % Lowercase spectrum
\DeclareMathOperator{\tr}{tr}     % Lowercase trace (if you already have \Tr)
\DeclareMathOperator{\supp}{supp} % Support
\DeclareMathOperator{\TrA}{Tr_{A}} % Partial trace over subsystem A
\DeclareMathOperator{\TrB}{Tr_{B}} % Partial trace over subsystem B
% Lie/algebra operators
\DeclareMathOperator{\Ad}{Ad}
\DeclareMathOperator{\Sym}{Sym}
\DeclareMathOperator{\su}{su}
% ======================================================================
% 8. STANDARD OPERATOR-LIKE SYMBOLS (Safe definitions via \providecommand)
% ======================================================================
\providecommand{\vol}{\operatorname{vol}}   % Volume form
\providecommand{\Vol}{\operatorname{Vol}}   % Total volume
\providecommand{\Ric}{\operatorname{Ric}}   % Ricci tensor
\providecommand{\Rm}{\operatorname{Rm}}     % Riemann tensor
\providecommand{\Tr}{\operatorname{Tr}}     % Uppercase trace
\providecommand{\Spec}{\operatorname{Spec}} % Uppercase spectrum
\providecommand{\diag}{\operatorname{diag}} % Diagonal matrix
\providecommand{\rank}{\operatorname{rank}} % Rank
\providecommand{\scal}{\mathcal R} % Scalar curvature
\providecommand{\Cl}{\operatorname{Cl}}     % Clifford algebra
\providecommand{\diam}{\operatorname{diam}} % Diameter

% ======================================================================
% 9. FUNCTIONS, TRIGONOMETRY, AND CALCULUS
% ======================================================================

% Codifferential naming conventions
% (Avoid \ddag: commonly defined by physics.sty and symbol packages)
\providecommand{\codiff}{\mathrm{d}^{\dagger}} % Codifferential operator

% Quantum mechanics bra-ket notation (needed for basis tables)
\providecommand{\ket}[1]{\lvert #1 \rangle}
\providecommand{\bra}[1]{\langle #1 \rvert}
\providecommand{\braket}[2]{\langle #1 \mid #2 \rangle}

% ======================================================================
% 10. PHYSICAL SCALES AND CONSTANTS (Paper 1c / Born Geometry)
% ======================================================================
\providecommand{\Rp}{\ensuremath{R_{\mathrm P}}}     % matched common EH curvature scale (not defined from observed G)
\providecommand{\hBorn}{\ensuremath{h_{\mathrm B}}} % physical representative 2\pi\hbar of one primitive Born scalar action cell
\providecommand{\UEHjoint}{\ensuremath{U_{J}}} % joint-Dirac EH curvature-coordinate scale 8 kappa_J
\providecommand{\Gread}{\ensuremath{G_{\rm scal}}} % scalar-action coordinate c_op^3/(h_B U); physical matching explicit
\providecommand{\GreadJoint}{\ensuremath{G_{J,\mathrm{raw}}}} % raw ordinary-trace coordinate = G_scal/96
\providecommand{\Rds}{R_{\mathrm{dS}}}    % de Sitter Curvature
\providecommand{\Mpl}{M_{\mathrm{Pl}}}    % Planck mass
\providecommand{\lp}{\ell_{\mathrm{Pl}}}  % Planck length
\providecommand{\Pmax}{P_{\mathrm{max}}}  % Maximum momentum
\providecommand{\Mplbar}{\bar{M}_{\mathrm{Pl}}} % Reduced Planck mass
\providecommand{\eff}{\mathrm{eff}}       % Subscript for effective quantities
\providecommand{\Cosmo}{\ensuremath{\Lambda}}                % Cosmological constant (use \Cosmo in source)
\providecommand{\LCDM}{\ensuremath{\Cosmo\mathrm{CDM}}}  % \LambdaCDM (typesets with \Cosmo)
\providecommand{\Lcut}{\ensuremath{\Lambda_{\mathrm{cut}}}}  % Cutoff scale (brace when used as script)
\providecommand{\Rcs}{\ensuremath{R_{\mathrm{cs}}}}      % Configuration-space radius
\providecommand{\CosmoEff}{\ensuremath{\Cosmo_{\eff}}}   % Effective cosmological constant

% Disambiguation: \Rmom is the momentum-space (S^3) radius. 
% Reserve plain R for the 4D curvature scalar in f(R).
\providecommand{\Rmom}{\ensuremath{R_{\mathrm{mom}}}} 
\providecommand{\Punit}{P_{0}} % Momentum unit set by \Rmom (= \hbar/\Rmom)
\providecommand{\Rvec}{\ensuremath{R_{\mathrm{vec}}}} % invariant vector-space radius in the reduced round sector
% Specific numeric values and expressions
\providecommand{\Vlogcap}{L_{\mathrm{Hopf}}^{(1)}}
\providecommand{\Hopf}{\operatorname{Hopf}}
% Scalar phase (S^1) sector + universal cover notation (Paper I-A)
\providecommand{\tper}{T_\star}      % Time period/scale
\providecommand{\Eper}{E_\star}      % Energy period/scale
\providecommand{\ttilde}{\tilde t}   % Universal cover time
\providecommand{\Etilde}{\tilde E}   % Universal cover energy

% ======================================================================
% 10A. SERIES-WIDE LOGICAL / NOTATIONAL DISAMBIGUATIONS
% ======================================================================
% Use these symbols to avoid collisions among: (i) the Paper--III
% no-insertion/geometric-closure class, (ii) Paper--I-A trimer diagonals,
% (iii) Paper--V-A/V-C family-space diagonals, (iv) scalar-index running,
% and (v) protected shell-clock labels.
\providecommand{\Rprim}{\ensuremath{\mathsf R_{\mathrm{prim}}^{\Sigma}}} % primitive finite structural readouts
\providecommand{\Rloop}{\ensuremath{\mathsf R_{\circ}^{\Sigma}}} % closed-loop readout role
\providecommand{\Rbdry}{\ensuremath{\mathsf R_{\partial}^{\Sigma}}} % boundary-conditioned readout role
\providecommand{\Rgen}{\ensuremath{\mathsf R_{\mathrm{fin}}^{\Sigma}}} % finite generated closure of primitive roles
\providecommand{\PhysOcc}{\operatorname{PhysOcc}}
\providecommand{\Actual}{\operatorname{Actual}}
\providecommand{\GeomLic}{\operatorname{GeomLic}}
\providecommand{\InvDist}{\operatorname{InvDist}}
\providecommand{\Dcl}{\ensuremath{\mathsf{D}_1}}      % unmarked positive Dirac-even closed spectral-readout form class
\providecommand{\Dheat}{\ensuremath{\mathsf D_{\mathrm{heat}}}} % one-parameter heat-semigroup subclass
\providecommand{\DtriI}{\ensuremath{D^{\mathrm{tri}}_1}} % Paper--I-A cyclic-trimer diagonal
\providecommand{\DtriII}{\ensuremath{D^{\mathrm{tri}}_2}}
\providecommand{\DfamI}{\ensuremath{D^{\mathrm{fam}}_1}} % Family-space diagonal selector
\providecommand{\DfamII}{\ensuremath{D^{\mathrm{fam}}_2}}
\providecommand{\alpharun}{\ensuremath{\alpha_{\mathrm{run}}}} % scalar spectral-index running
\providecommand{\Ksh}{\ensuremath{K_{\mathrm{sh}}}} % protected shell-clock coordinate
\providecommand{\ddef}{\ensuremath{\delta_{\mathrm{def}}}} % universal defect scalar after Paper--III realization

% ======================================================================
% 10B. SERIES-WIDE LOGICAL STATUS TAGS
% ======================================================================
% Upright small-caps wrapper for audit/status labels.  It prevents LaTeX from
% trying to use unavailable small-caps italic font shapes inside theorem bodies
% such as TU/lmr/m/scit.
\providecommand{\AuditLabel}[1]{}
\renewcommand{\AuditLabel}[1]{{\normalfont\scshape #1}}

% These tags are bookkeeping labels for consistency audits. They do not
% create new mathematical hypotheses; they record how a datum is being used
% in the proof chain.
\providecommand{\StatusPrimitive}{}
\renewcommand{\StatusPrimitive}{\AuditLabel{Primitive}}
\providecommand{\StatusDerived}{}
\renewcommand{\StatusDerived}{\AuditLabel{Derived}}
\providecommand{\StatusImported}{}
\renewcommand{\StatusImported}{\AuditLabel{Imported}}
\providecommand{\StatusConditional}{}
\renewcommand{\StatusConditional}{\AuditLabel{Conditional}}
\providecommand{\StatusBridge}{}
\renewcommand{\StatusBridge}{\AuditLabel{Bridge-derived}}
\providecommand{\StatusComparison}{}
\renewcommand{\StatusComparison}{\AuditLabel{Comparison-only}}

% ======================================================================
% 10C. SERIES-WIDE CLAIM-WORDING AUDIT LABELS
% ======================================================================
% These are textual audit labels used to keep modal language consistent
% across the series.  They introduce no new mathematical assumptions.
\providecommand{\ClaimForced}{}
\renewcommand{\ClaimForced}{\AuditLabel{forced}}
\providecommand{\ClaimUnique}{}
\renewcommand{\ClaimUnique}{\AuditLabel{unique}}
\providecommand{\ClaimCanonical}{}
\renewcommand{\ClaimCanonical}{\AuditLabel{canonical}}
\providecommand{\ClaimMinimal}{}
\renewcommand{\ClaimMinimal}{\AuditLabel{minimal}}
\providecommand{\ClaimClosed}{}
\renewcommand{\ClaimClosed}{\AuditLabel{closed}}
\providecommand{\ClaimPrediction}{}
\renewcommand{\ClaimPrediction}{\AuditLabel{prediction}}

% ======================================================================
% 10D. SERIES-WIDE THEOREM-CHAIN AUDIT LABELS
% ======================================================================
% These labels mark the direction of load-bearing theorem chains. They are
% audit annotations only; they introduce no new mathematical assumptions.
\providecommand{\ChainInput}{}
\renewcommand{\ChainInput}{\AuditLabel{Chain input}}
\providecommand{\ChainOutput}{}
\renewcommand{\ChainOutput}{\AuditLabel{Chain output}}
\providecommand{\ChainGuard}{}
\renewcommand{\ChainGuard}{\AuditLabel{Chain guard}}
\providecommand{\ChainNoBackflow}{}
\renewcommand{\ChainNoBackflow}{\AuditLabel{No backflow}}

% ======================================================================
% 10E. SERIES-WIDE LOCAL THEOREM-AUDIT LABELS
% ======================================================================
% These labels mark a statement-level dependency audit around individual
% theorem/definition nodes. They are bookkeeping labels only.
\providecommand{\LocalAudit}{}
\renewcommand{\LocalAudit}{\AuditLabel{Local audit}}
\providecommand{\UsesOnly}{}
\renewcommand{\UsesOnly}{\AuditLabel{Uses only}}
\providecommand{\Exports}{}
\renewcommand{\Exports}{\AuditLabel{Exports}}
\providecommand{\DoesNotUse}{}
\renewcommand{\DoesNotUse}{\AuditLabel{Does not use}}

% ======================================================================
% 10F. SERIES-WIDE NUMERICAL-ROW AUDIT LABELS
% ======================================================================
% These labels classify individual rows in numerical, benchmark, status,
% and comparison tables.  They are audit annotations only and do not change
% any theorem, value, formula, or comparison number.
\providecommand{\RowAudit}{}
\renewcommand{\RowAudit}{\AuditLabel{Numerical-row audit}}
\providecommand{\RowTheorem}{}
\renewcommand{\RowTheorem}{\AuditLabel{Theorem-fixed row}}
\providecommand{\RowBridge}{}
\renewcommand{\RowBridge}{\AuditLabel{Bridge-fixed row}}
\providecommand{\RowRepresentative}{}
\renewcommand{\RowRepresentative}{\AuditLabel{Representative row}}
\providecommand{\RowComparison}{}
\renewcommand{\RowComparison}{\AuditLabel{Comparison row}}
\providecommand{\RowCalibrated}{}
\renewcommand{\RowCalibrated}{\AuditLabel{Calibrated row}}
\providecommand{\RowEstimate}{}
\renewcommand{\RowEstimate}{\AuditLabel{Estimate row}}
\providecommand{\RowOpen}{}
\renewcommand{\RowOpen}{\AuditLabel{Open row}}
\providecommand{\RowNoSelector}{}
\renewcommand{\RowNoSelector}{\AuditLabel{No selector}}

% ======================================================================
% 10G. TABLE HELPERS
% ======================================================================
% Table helpers used by Paper V-C.  The unused readout-label preset has been
% removed; only active table helpers remain.
\providecommand{\gpsection}[1]{%
  \addlinespace[0.55em]
  \multicolumn{3}{@{}l}{\textsc{#1}}\\[-0.25em]
  \midrule
}
\providecommand{\mathcell}[1]{%
  \(\begin{gathered}#1\end{gathered}\)%
}

\newcommand{\DerBy}[1]{%
  \begin{tabular}[t]{@{}l@{}}#1\end{tabular}%
}

\newcommand{\TableNotes}[1]{%
  \par\medskip
  \begin{minipage}{0.94\linewidth}
  \footnotesize\emph{Notes.} #1
  \end{minipage}\par
}
% ======================================================================
% 10H. PAPER VI-A / QCD NOTATION
% ======================================================================
% Shared symbols used by Paper VI-A and downstream QCD/Yang--Mills handoffs.
\providecommand{\Pqcd}{\mathfrak P_{\rm VIII}^{\rm QCD}}
\providecommand{\Iqcd}{\mathfrak I_{\rm QCD}^{(7)}}
\providecommand{\muQCD}{\mu_{\rm QCD}^{\rm int}}
\providecommand{\LambdaQCD}{\Lambda_{\rm QCD}}
\providecommand{\LambdaBH}{\Lambda_{\rm BH}^{\rm QCD}}
\providecommand{\HQCD}{\mathcal H_{\rm QCD}^{\rm BH}}
\providecommand{\HQCDfin}{H_{\rm QCD}^{\rm fin}}
\providecommand{\Hraw}{\mathcal H_{\rm raw}}
\providecommand{\HphysQCD}{\mathcal H_{\rm phys}^{\rm QCD}}
\providecommand{\Vcol}{V_C}
\providecommand{\Ps}{P_{\rm sing}}
\providecommand{\PiOct}{\Pi_8}
\providecommand{\MSbar}{\overline{\rm MS}}
\providecommand{\YM}{\mathrm{YM}}
\providecommand{\Ktri}{K_{\triangle}}
\providecommand{\KYM}{K_{\YM}^{\rm BH}}
\providecommand{\aBH}{a_{\rm BH}}
\providecommand{\epsBH}{\varepsilon_{\rm BH}}
\newcommand{\RowConstructive}{\textnormal{\textsc{Constructive}}}
\newcommand{\RowObligation}{\textnormal{\textsc{Bridge obligation}}}
\newcommand{\RowTarget}{\textnormal{\textsc{Target row}}}
\newcommand{\RowValidation}{\textnormal{\textsc{Validation row}}}
\providecommand{\slashed}[1]{\not{#1}}

% ======================================================================
% 10I. PAPER VI-B / CONTINUOUS-QFT NOTATION
% ======================================================================
% Shared symbols used by Paper VI-B and any downstream continuous-QFT or
% Yang--Mills handoff paper.  All definitions are guarded so individual
% papers may override them locally before loading this file if needed.
\providecommand{\A}{\mathfrak A}
\providecommand{\Hphys}{\mathcal H_{\rm phys}}
\providecommand{\QCD}{\mathrm{QCD}}
\providecommand{\BH}{\mathrm{BH}}
\providecommand{\loc}{\mathrm{loc}}

\providecommand{\RBH}{\mathcal R_{\rm BH}}
\providecommand{\DG}{\mathfrak D_G}
\providecommand{\PW}{\mathcal{P\!W}} % Peter--Weyl spectral-symbol subspace
\providecommand{\Toe}{\operatorname{Toe}} % Toeplitz compression/representative
\providecommand{\Hol}{\operatorname{Hol}} % Holonomy

% ======================================================================
% 11. LATE-LOADED PACKAGES (Must be at the end of the preamble)
% ======================================================================
% Load late packages only if a paper has not already loaded them, avoiding
% option clashes in paper-specific preambles.
\makeatletter
\@ifpackageloaded{hyperref}{}{\usepackage{hyperref}}
\makeatother
% --- Series-wide disambiguations (Pass 3) ---
\providecommand{\diracdelta}{\delta_{\rm D}} % Dirac delta distribution (delta reserved for hierarchy)
\providecommand{\Hg}{H_{\mathrm g}} % Hodge harmonic projector (keep H for Hubble)
\providecommand{\Gr}{G_{\mathrm r}} % Green operator (keep G for Newton)
\providecommand{\FD}{\mathrm{D}} % functional/variational derivative symbol (delta reserved)
\providecommand{\Id}{\mathbf{1}} % identity/Kronecker tensor (avoid \delta_{ij})


% BEGIN INLINE SOURCE: Paper00-ActiveExternalReferences-v2_130.tex
% Framework v2.191: current reference binding; existing definitions retained.
% Framework v2.190: current22-root registry with native clock-evolution proofs.
% Framework v2.189: exact current22-root external registry, including retained v2.188 hardening owners.
% Framework v2.187: current registry synchronization only; source release v2.186 preserved.
% Framework v2.186: authoritative current 22-paper registry.
% Framework v2.182: all current 22 logical paper AUX bindings.
% Framework v2.170: current I-B v1.43 external AUX binding.
% Framework v2.167: exact current 22-root registry after CP--BL2 integration.
% Framework v2.166: exact current 22-root compact-release AUX registry.
% Framework v2.165: authoritative current root bindings.
% Framework v2.164: exact current 22 manuscript AUX bindings.
% Framework v2.162: finite metric, functional Hilbert and scoped Ward integration.
% Framework v2.161: current 22 manuscript AUX bindings; source only.
% Framework v2.156 / registry v2.113: synchronized current 22 manuscript roots; source-only registry.
% Framework v2.143 / registry v2.111: active manuscript AUX basenames; no generated AUX or PDF included.
% ======================================================================
% Central external-reference imports, numbers only, no PDF jumping
% ======================================================================

% xr-hyper is loaded centrally so that papers may retain their existing
% \usepackage{xr-hyper} lines without redefining \externaldocument after this
% registry is installed.  If plain xr was loaded earlier, retain it.
\makeatletter
\@ifpackageloaded{xr-hyper}{}{%
  \@ifpackageloaded{xr}{}{\usepackage{xr-hyper}}%
}
\makeatother

% Make external-document registration idempotent.  The registry identifies a
% logical paper by its Paper<number><optional letter> prefix.  Prefix-qualified
% imports remain independent.  Current paper roots contain no local
% \externaldocument calls; this guard also makes accidental duplicates harmless.
\ExplSyntaxOn
\prop_new:N \g_bh_external_refs_prop
\cs_new_eq:NN \BHOriginalExternalDocument \externaldocument

\cs_new_protected:Npn \bh_external_document:nnnn #1#2#3#4
 {
  \tl_set:Nn \l_tmpa_tl {#3}
  \regex_replace_once:nnN
   { \A (?: xref/ )? (Paper[0-9]+[A-G]?) .* \Z }
   { \1 }
   \l_tmpa_tl
  \tl_set:Nx \l_tmpb_tl { #1 | \tl_use:N \l_tmpa_tl }
  \prop_if_in:NVTF \g_bh_external_refs_prop \l_tmpb_tl
   {
    \PackageInfo{BornHopf}
     {Skipping~duplicate~external~refs~for~#3}
   }
   {
    \prop_gput:NVn \g_bh_external_refs_prop \l_tmpb_tl {#3}
    \IfNoValueTF{#4}
     { \BHOriginalExternalDocument[#1][#2]{#3} }
     { \BHOriginalExternalDocument[#1][#2]{#3}[#4] }
   }
 }

\RenewDocumentCommand{\externaldocument}{O{} O{} m o}
 { \bh_external_document:nnnn {#1}{#2}{#3}{#4} }

\NewDocumentCommand{\BHExternalRefs}{m}
 {
  % Do not import the aux file of the logical paper currently being compiled,
  % even when the job and registry use different version suffixes.
  \tl_set:Nx \l_tmpa_tl { \c_sys_jobname_str }
  \regex_replace_once:nnN
   { \A (Paper[0-9]+[A-G]?) .* \Z }
   { \1 }
   \l_tmpa_tl
  \tl_set:Nn \l_tmpb_tl { #1 }
  \regex_replace_once:nnN
   { \A (Paper[0-9]+[A-G]?) .* \Z }
   { \1 }
   \l_tmpb_tl
  \str_if_eq:eeTF { \tl_use:N \l_tmpa_tl } { \tl_use:N \l_tmpb_tl }
   {
    \PackageInfo{BornHopf}
     {Skipping~external~refs~for~#1~because~this~is~the~current~document}
   }
   {
    \IfFileExists{xref/#1.aux}
     { \externaldocument[][nocite]{xref/#1}[#1.pdf] }
     {
      \IfFileExists{#1.aux}
       { \externaldocument[][nocite]{#1} }
       {
        \PackageInfo{BornHopf}
         {External~aux~file~#1.aux~not~found;~no~labels~imported}
       }
     }
   }
 }
\ExplSyntaxOff

% External references are automatic.  This file is the sole owner of both the
% bibliography and the current Paper-I-A--VII-D AUX registry (30 roots including I-C).  Missing auxiliaries
% are informational on the first pass and are populated by the source-only
% two-pass series build.  Basenames deliberately omit temporary download/copy
% suffixes such as ``(1)'' or ``-1''.
  % Paper I-A: round S3 compendium.
  \BHExternalRefs{Paper1A-Geometry-v5_180}

  % Paper I-B: complete exceptional geometry and finite spectral structures.
  \BHExternalRefs{Paper1B-ExceptionalGeometry-v1_22}
  % Paper I-D: technical companion of I-A through I-C.
  \BHExternalRefs{Paper1D-TechnicalMaterial-v1_69}
  % Paper I-C: reciprocal M8 compendium.
  \BHExternalRefs{Paper1C-M8Geometry-v1_144}

\ifdefined\BHMathOnly
  % Papers I-A and I-B import only the mathematical atlas labels.
\else
  % Paper II: five physical-mathematical main papers plus one technical proof volume.
  \BHExternalRefs{Paper2A-GeometricConstitution-v3_227}
  \BHExternalRefs{Paper2B-MomentumCarrier-v2_147}
  \BHExternalRefs{Paper2C-ReciprocalLocalization-v2_173}
  \BHExternalRefs{Paper2D-SymmetrySpectra-v2_144}
  % The former II-F export labels remain in II-E; current II-F is the technical volume.
  \BHExternalRefs{Paper2E-KinematicExports-v3_139}
  % Paper II-F: complete technical proof and calculation owner for II-A--II-E.
  \BHExternalRefs{Paper2F-TechnicalMaterial-v1_70}

  % Paper III: Dynamics and Intrinsic Geometric Constitution (six components).
  \BHExternalRefs{Paper3A-GeneratedDynamics-v1_180}
  \BHExternalRefs{Paper3B-SpectralAction-v1_191}
  \BHExternalRefs{Paper3C-ProtectedHierarchy-v1_193}
  \BHExternalRefs{Paper3D-RunningResponse-v1_167}
  \BHExternalRefs{Paper3E-GravityAction-v1_203}
  \BHExternalRefs{Paper3F-CompleteDynamics-v1_276}
  % Paper III-G: complete preserved technical proof owner.
  \BHExternalRefs{Paper3G-TechnicalMaterial-v1_57}

  % Paper IV-A: generated finite quantum mechanics.
  \BHExternalRefs{Paper4A-FiniteQuantumMechanics-v4_216}
  % Paper IV-B: complete finite quantum processes and physical response.
  \BHExternalRefs{Paper4B-QuantumResponse-v1_54}
  % Paper IV-C: extended proofs and finite technical calculations.
  \BHExternalRefs{Paper4C-TechnicalMaterial-v1_50}

  % Paper V-A: The Standard Model.
  \BHExternalRefs{Paper5A-FiniteMatterGeometry-v3_218}

  % Paper V-B: Finite Spectrum and Flavor Closure.
  \BHExternalRefs{Paper5B-TerminalSpectrum-v5_202}

  % Paper V-C: Finite Interaction Boundaries of the Standard Model.
  \BHExternalRefs{Paper5C-InteractionResponse-v2_182}
  % Paper V-D: preserved technical proof and calculation owner.
  \BHExternalRefs{Paper5D-TechnicalMaterial-v1_58}

  % Paper VI-A: Gravity and Closed Cosmology.
  \BHExternalRefs{Paper6A-GravityCosmology-v2_267}

  % Paper VI-B: Geometric Dark Sector and Cosmological Readouts.
  \BHExternalRefs{Paper6B-CosmologyReadouts-v1_214}

  % Paper VI-C: complete technical proof and comparison owner.
  \BHExternalRefs{Paper6C-TechnicalMaterial-v1_53}

  % Paper VII-A: Finite QCD Kernel and Nonperturbative Closure.
\BHExternalRefs{Paper7A-QCDKernel-v1_180}

  % Paper VII-B: Hadronic, Thermal, and Nuclear QCD Readouts.
  \BHExternalRefs{Paper7B-QCDReadouts-v1_165}

  % Paper VII-C: Continuous Quantum Field Theory.
  \BHExternalRefs{Paper7C-ContinuousQFT-v2_208}

  % Paper VII-D: complete technical material for VII-A through VII-C.
  \BHExternalRefs{Paper7D-TechnicalMaterial-v1_60}

\fi

% END INLINE SOURCE: Paper00-ActiveExternalReferences-v2_130.tex


\hypersetup{hidelinks}

\usetikzlibrary{positioning, calc} 


\title{\textbf{Born--Hopf Geometry II-B: The Generated Momentum Carrier}\\[0.5em]
\large Compact SU(2), Spin(4), finite level one, and recovery}

\author{Peter Kugelmann\thanks{\texttt{pkugelmann@proton.me}}}
\date{Version 2.147 --- 20 November 2026}

\setlength{\emergencystretch}{3em}
\begin{document}

\maketitle

\begin{abstract}
Relative to the constituted round momentum parent
$M_p\cong SU(2)\cong S^3(\Rmom)$ and the frozen
Paper~II-A momentum-seed frontier, this paper derives the complete
GP2-generated momentum type-support.  The connected isometry group
$SO(4)$ and its spin lift $\Spin(4)$ act on the single parent; they
are not additional momentum carriers.  The intrinsic scalar unit,
lowest chiral spinors, first scalar harmonics and linearized
gradient-bracket Killing fields generate the six irreducible types
$(0,0)$, $(\frac12,0)$, $(0,\frac12)$, $(\frac12,\frac12)$,
$(1,0)$ and $(0,1)$.  Their explicit one-copy model is
$\mathcal H_1$ of complex dimension $15$.

The intrinsic representation word length of that GP2 support gives
$m_p=1$ without a tensor-tower selector or a universal GP3
multiplicity-one premise.  A separate complete whole-candidate
comparison certificate is required to identify every physical GP3
endpoint with the one-copy model; that certificate remains open.
The canonical smooth presentation, finite-state and operator-algebra
capacities, product-closure obstructions, continuum recovery and
inverse diagnostics are developed at their stated mathematical
and conditional scopes.  Multiplicity-blind native tests, including
finite-ancilla protocols, cannot themselves discharge the GP3
comparison obligation.  Abstract carriers, state spaces and
conditional finite-support comparisons do not by themselves establish
actual preparations, interactions, rates or measurements.
\end{abstract}

\tableofcontents
\listoftables
\newpage

\section{Momentum-role input and constructor dependencies}
\label{paper2b:sec:organized-scope}

\subsection{Introduction}
\label{p2b:sec:introduction}

\paragraph{Full technical derivations.}
The active auxiliary classifications, exact calculations, and negative
results supporting this paper are collected without abridgment in
Paper~II-F, Section~\ref{p2ftech:sec:momentum}.
That companion is a proof owner, not an additional constructor,
physical occurrence, state, or calibration assumption.

Paper~II-B imports the Paper~II-A momentum parent
\[
 M_p\cong SU(2)\cong S^3(\Rmom)
\]
and the constitutional semantics
\[
 \begin{aligned}
  \mathrm{G0}&:\text{ actualized/distinguishable},&
  \mathrm{GP1}&:\text{ finite/bounded},\\
  \mathrm{GP2}&:\text{ exact/host-faithful},&
  \mathrm{GP3}&:\text{ irreducible/least}.
 \end{aligned}
\]
Paper~II-A owns the parent-role allocation.  This paper owns the next
constructor-relative step: it derives the connected isometry completion,
classifies the \(\Spin(4)\) sectors, identifies the complete generated
momentum-seed carrier, and only afterward attaches the intrinsic word-length
label \(m_p=1\).  Thus no tower membership, spectral exhaustivity, component
ceiling, or numerical level is used to select the initial momentum carrier.

The established type-support and level chain is
\[
 S^3(\Rmom)
 \Longrightarrow \Spin(4)
 \Longrightarrow E_0,S_+,S_-,E_1
 \Longrightarrow (\mathcal K_{\rm Kill})_{\mathbb C}
 \Longrightarrow \mathcal G^{\rm car}_{\Sigma_p}\cong\mathcal H_1
 \Longrightarrow m_p=1.
\]
The stronger arrow from this common generated content to the carrier of
\emph{every} GP3-irreducible structured endpoint requires the whole-candidate
comparison certificate of Theorem~\ref{p2b:thm:momentum-carrier-rigidity}.
The generic word tower and its component/shell bounds remain useful
classification and recovery structures, especially for later roles, but they
are not the constitutional origin of the momentum seed.

\paragraph{Co-realization of the generated momentum constituents.}
The momentum seed, \(E_0\), \(S_+\), \(S_-\), \(E_1\),
\(\mathcal K_{\rm Kill}\), its common generated support, and the associated one-copy smooth model lie
over one rooted constructor component.  Their
universal bridges are constituent ancestry, inclusions, and the licensed
section-to-operation map.  Different dimensions and sorts prevent a
pairwise Born interpretation.  A Killing section and its induced operation
remain different typed nodes but, when the derivation constructor is
licensed, require only the one momentum-root occurrence.

\subsubsection{Logical scope and import ledger}
\label{p2b:subsec:import-ledger}
The following table distinguishes constitutional inputs from structures
derived in this paper.
\begingroup
\small
\setlength{\LTpre}{6pt}
\setlength{\LTpost}{6pt}
\renewcommand*{\theHtable}{importledger.\arabic{table}}%
\begin{longtable}{L{0.20\linewidth}L{0.35\linewidth}L{0.37\linewidth}}
\toprule
\textbf{Source} & \textbf{Imported datum} & \textbf{Not imported} \\
\midrule
\endfirsthead
\toprule
\textbf{Source} & \textbf{Imported datum} & \textbf{Not imported} \\
\midrule
\endhead
Paper~II-A
& the least parent-role skeleton
\(M_p\simeq SU(2)\), its round negative-Killing geometry up to \(\Rmom\),
the stage-safe G0 geometric-actualization certificate and its
occurrence/readout boundary,
GP1 finite/bounded realization, GP2 exact/host-faithful constructor-relative
closure including finite generated-image inheritance, canonical-image
objectification, and the same-role no-splitting lemma, the frozen momentum-seed
dependency frontier, generated-orbit and unresolved-family inheritance,
no-representative-selection semantics, and the universal GP3 host-faithful
reduction semantics
& no finite momentum carrier, no momentum level, no arbitrary tensor closure,
no component ceiling, no selected coordinate-completion point, and no
core-support occurrence merely from the parent theorem \\
Paper~I-A
& the round-sphere scalar and spinor harmonic package and the intrinsic
\(\Spin(4)\) representation dictionary
& no physical selection of a finite carrier or level \\
Paper~I-B
& the finite representation identities, the canonical bracket/Killing
comparison and the retained one-copy smooth model
\(J_R:\mathcal H_1\simeq\mathcal L_R\)
& no independent GP2 carrier-generation or universal GP3 endpoint proof \\
Paper~I-B: \(\mathsf C_{\rm PI}\)

& only where explicitly invoked, normalized finite coordinate intertwiners,
the protected Euclidean Clifford vector, and its independent algebraic scale
\(\rho_{\rm prot}\)
& one representative of
\(\mathfrak M_{\rm coord}(\mathcal H_1)/\!\sim\cong\mathbb R_{\ge0}^3\),
with residual coordinates \((r_{1,0},r_{0,1},\varrho_{\rm sp})\), not a point
selected by \(\mathcal H_1\), and no identification
\(\rho_{\rm prot}=\Rmom\) without a separate convention \\
Paper~I-B: \(\mathsf{Supp}_{2T}\)
& only where explicitly invoked, the exceptional-core atlas/constructor and
its finite octet/septet and oriented trimer-pair supports with their
constituted transitive actions
& no implication from generic point resolution alone to \([2T]\): an
occurring point-resolution support gives \(\mathsf{Supp}_{\varepsilon}(X)\),
and \(\mathsf{Supp}_{2T}\) co-realizes only after the additional
candidate-independent core constructor plus the Paper-II-A domain bridge, or
a direct core-constructor occurrence; no attachment from \(m_p=1\), no
labelled representative, and no physical interpretation of the weights \\
Optional tower grammar
& the mathematical family \(\mathcal H_m\), word stages, component ceilings,
shells, and Peter--Weyl recovery structures developed below
& no premise for the initial momentum-seed theorem; TWR/SPEX/FIT and
\(\mathsf P_{\rm sli}\) are optional tower conditions rather than
momentum-selection axioms \\
\bottomrule
\end{longtable}
\endgroup
\addtocounter{table}{-1}%
Everything else proved below is native to II-B.  In particular, later
Paper-II dynamics, reciprocal-position growth, normalization, and freezing
conditions are neither premises nor proof sources for the momentum seed.

\paragraph{Physical occurrence boundary.}
All physical interfaces below are role-specific specializations of
Paper~II-A's stage-safe G0 geometric-actualization certificate.  Whenever an
interface is asserted to occur, its certificate must exhibit a contingent
token of a geometrically licensed, claim-relative invariant distinction with
result-independent provenance.  The interface tag alone asserts no occurrence
and supplies no independent actualization principle.  No instance reopens the
frozen upstream geometry; a new role enters only through a separately declared
downstream dependency frontier.  Analytic, regularity, continuation, and
reconstruction hypotheses remain separately displayed.  The transfer parent
and the host theorem below are generated descendants of the already occurring
operation role, so neither needs a second occurrence witness.

\subparagraph{Role-local rank boundary.}
All operational rank-one statements in this paper concern the single primitive
momentum role.  They neither remove the independently retained scalar
\(U(1)\) role nor exclude higher-rank products of separately constituted
reciprocal roles.

\section{Round SU(2) momentum geometry and Spin(4) lift}
\label{paper2b:sec:organized-round}

\subsection{Round momentum geometry and spinorial seeds}
\label{p2:sec:classical-momentum-geometry-operation-first}

Theorem~\ref{p2:thm:dependency-complete-parent-allocation} allocates the
generated connected noncommutative transfer role, and
Definition~\ref{p2:def:momentum-branch-specialization} records the derived
relation-exact \(SU(2)\) parent used by the remaining Paper-II
construction.  Its sphere topology and rank-one geometry are the II-A results
in Theorems~\ref{p2:thm:rank-one-topology-spin} and
\ref{p2:thm:rank-one-space-form-einstein}.  From
this point onward, \emph{momentum space} means
\[
        M_p\cong S^3_A(\Rmom)\cong SU(2).
\]
The circle \(S^1_A\cong U(1)\) is exported as the compact scalar succession
parent but is not part of the Paper-II momentum construction.  The round
metric, its \(SO(4)\) isometry group, and the lifted \(\Spin(4)\) symmetry are
consequences of the derived \(SU(2)\) branch.
The compatible Hopf fibrations form a conjugacy class with no preferred
member; one left/common-phase Hopf/contact representative and its normalization
are now fixed from the Paper-I-A atlas.  This convention does not constitute the
independent scalar parent.  The finite momentum seed is derived below from the constructor-relative host
signature and finite generated-image inheritance.  The generic \(\Spin(4)\)-equivariant
tower is retained only as secondary mathematical infrastructure.  Coordinate
actions form the unresolved family \(\mathfrak M_{\rm coord}\), while
\(\mathsf C_{\rm PI}\) is the normalized Paper-I-A representative used when
its numerical formulas are quoted.
The center-blind \(SO(3)\) alternative has the same local geometry but is excluded
at the preceding GP2 exact/host-faithful globalization step.

\begin{theorem}[Derived isometry completion of the single \(SU(2)\) parent]
\label{p2:thm:single-su2-derived-spin4-isometries}
Let
\[
 G_{\rm op}=M_p\cong SU(2)\cong S^3(\Rmom)
\]
be the single nonabelian operation parent selected above, and equip it with
the canonical automorphism-invariant negative-Killing metric of
Theorem~\ref{p2:thm:master-compact-role-allocation}.  Then
this metric is round up to its radius scale, and
\[
 \operatorname{Isom}_0(M_p)
 \cong
 \frac{SU(2)_L\times SU(2)_R}{\Delta\ZZ_2}
 \cong SO(4).
\]
The universal spin cover of this connected isometry group is
\[
 \Spin(4)\cong SU(2)_L\times SU(2)_R,
\]
acting by
\[
 (a,b)\cdot q=a q b^{-1}.
\]
The point stabilizer is \(SU(2)_{\rm diag}\), and therefore
\[
 M_p\cong
 \frac{\Spin(4)}{SU(2)_{\rm diag}}
 \cong S^3.
\]
In particular,
\[
 \operatorname{Lie}(G_{\rm op})=\mathfrak{su}(2),
 \qquad
 \operatorname{Lie}(\operatorname{Isom}_0(M_p))
 =\mathfrak{so}(4)
 \cong\mathfrak{su}(2)_L\oplus\mathfrak{su}(2)_R.
\]
The second algebra is the Killing algebra of the derived metric, not the Lie
algebra of two momentum-role carriers.
\end{theorem}

\begin{proof}
For the simple compact Lie algebra \(\mathfrak{su}(2)\), every positive
\(\operatorname{Ad}\)-invariant inner product is a positive multiple of the
negative Killing form.  Hence the derived bi-invariant metric is the round
metric after fixing \(\Rmom\).  Left and right translations are isometries,
so the displayed action defines a homomorphism
\[
 SU(2)_L\times SU(2)_R\longrightarrow\operatorname{Isom}_0(M_p).
\]
If \(a q b^{-1}=q\) for every \(q\), evaluation at \(q=1\) gives \(a=b\),
and then \(a\) must be central.  The kernel is therefore
\(\{(1,1),(-1,-1)\}=\Delta\ZZ_2\).  The isometry group of the round \(S^3\)
is \(O(4)\), whose identity component is \(SO(4)\); the containment
\(\operatorname{Isom}_0\subseteq SO(4)\) used next follows because an isometry
of a connected Riemannian manifold is determined by its value and differential
at one point, and the displayed action is already transitive on oriented
orthonormal frames of \(S^3\).  Since the domain is compact, its image is
a compact and hence closed subgroup of \(SO(4)\).  The quotient by the discrete kernel has dimension six, equal to
\(\dim SO(4)\); its image therefore has full Lie algebra and is open in the
connected group \(SO(4)\), so it is all of \(SO(4)\).  The stabilizer of
\(q=1\) is the diagonal subgroup, proving the homogeneous presentation and
the Lie algebra statements.
\end{proof}

\begin{remark}[Carrier, symmetry, and finite-module boundary]
\label{p2:rem:carrier-symmetry-module-boundary}
The identities
\(S^3\cong SU(2)\) and
\(\Spin(4)\cong SU(2)_L\times SU(2)_R\cong S^3\times S^3\) refer to
different objects.  The former is the three-dimensional carrier.  The latter
is the six-dimensional lifted isometry group acting on that carrier.  A
\(\Spin(4)\)-module used below is a representation of the lifted symmetry
group; it is not a second geometric carrier.
\end{remark}

\begin{proposition}[Finite sector grammar and the exact descent boundary]
\label{p2:prop:algebra-sector-grammar}
Let a GP1-finite, GP2-host-faithful complex Hilbert realization of the momentum
host carry the derived connected isometry symmetry of
Theorem~\ref{p2:thm:single-su2-derived-spin4-isometries}.  The resulting
finite-dimensional unitary representation integrates to the universal lift
\(\Spin(4)\), and its irreducible complex sectors are, up to unitary
equivalence,
\[
 V_{j_L}\otimes V_{j_R},
 \qquad
 (j_L,j_R)\in
 \tfrac12\ZZ_{\ge0}\times\tfrac12\ZZ_{\ge0}.
\]
Every such Lie-algebra representation integrates uniquely to the simply
connected derived symmetry group
\(\Spin(4)=SU(2)_L\times SU(2)_R\).  This integration introduces neither a
second momentum parent nor a choice of deck data on \(M_p\).

Let \(z=(-1,-1)\) be the kernel of
\(\Spin(4)\to SO(4)\).  On \(V_{j_L}\otimes V_{j_R}\),
\[
 \pi_{j_L,j_R}(z)
 =
 (-1)^{2(j_L+j_R)}\mathbf1.
\]
Consequently conjugation on the endomorphisms of any single irreducible sector
descends to \(SO(4)\).  On a direct sum of sectors, conjugation on an
inter-sector block descends precisely when its source and target have the same
central character.  The generated operator doctrine used in this paper is
within that descent domain: Killing generators and equivariant projectors are
block diagonal, while the coordinate type
\((\tfrac12,\tfrac12)\) changes \(j_L+j_R\) by an integer and therefore
preserves the central character.  No descent statement is made for arbitrary
endomorphisms mixing the two central-character classes.
\end{proposition}

\begin{proof}
Finite-dimensional unitary representations of a compact semisimple Lie
algebra are completely reducible
\cite{fulton1991representation,Hall2015lie}.  The irreducibles of a direct sum
are outer tensor products of irreducibles of its two factors, and the
irreducibles of \(\mathfrak{su}(2)\) are indexed by
\(j\in\tfrac12\ZZ_{\ge0}\) \cite{fulton1991representation}.  Lie's
integration theorem \cite{Hall2015lie} gives a unique
representation of the simply connected group \(\Spin(4)\).

The central element \(-1\in SU(2)\) acts on \(V_j\) by
\((-1)^{2j}\), giving the displayed central character.  On
\(\operatorname{Hom}(V_\lambda,V_\mu)\), conjugation by \(z\) is the quotient
of the two central characters, so it is trivial exactly when they agree.
Finally, a coordinate channel has
\((\Delta j_L,\Delta j_R)=(\pm\tfrac12,\pm\tfrac12)\), hence changes
\(j_L+j_R\) by \(0\) or \(\pm1\).  This proves the stated descent boundary.
\end{proof}

\begin{remark}[Constitution status of the rotation--transvection sextet]
\label{p2:rem:identity-anchored-split}
The rotation--transvection split
\(\mathfrak{so}(4)=\mathfrak{so}(3)_J\oplus\mathfrak p\) is anchored at the
identity of the operation parent: \(\mathfrak{so}(3)_J\) is the isotropy
algebra at \(e\), and \(\mathfrak p\) is its invariant Killing complement.
The identity and the operation law are already constituted, so this split
does not select an external basepoint or a second operation family.  Inner
automorphisms \(q\mapsto aqa^{-1}\) are derived covariance transformations,
as follows directly from the constituted group law.
The finite-dimensional unitary carrier action in
Proposition~\ref{p2:prop:algebra-sector-grammar} is the realized form of the
derived host symmetry required by GP1 finiteness and GP2 host faithfulness.
The additional faithful six-operator realization assumed in
Theorem~\ref{p2:thm:compact_selection} remains an explicit hypothesis of that
algebraic classification theorem.
\end{remark}

\subsubsection{Quantized compact momentum algebra}
\label{p2:subsec:quantized-algebra}

Once the parent-allocation and round-geometry theorems have identified the
classical momentum manifold with the round \(3\)-sphere,
Theorem~\ref{p2:thm:single-su2-derived-spin4-isometries} identifies its
six-dimensional Killing algebra.  The following theorem classifies that
isometry algebra in the rotation--transvection basis; it does not replace the
three-dimensional operation-parent algebra by two parent algebras.

\begin{remark}[Compact branch used in the positive chain]
\label{p2:rem:compact-branch}
The round parent already fixes the compact Killing algebra, so the positive
construction uses
\[
 \varkappa=\frac{1}{\Rmom^2}>0,
 \qquad
 \mathfrak g\cong\mathfrak{so}(4)
 \cong\mathfrak{su}(2)_L\oplus\mathfrak{su}(2)_R.
\]
The flat and Lorentzian algebraic alternatives, together with the general
two-parameter covariant bracket reduction, are not needed in the positive
proof and are deferred to Paper~II-F, Section~\ref{p2b:app:algebraic-branch-diagnostics}.
The carrier group itself remains \(SU(2)\), with Lie algebra
\(\mathfrak{su}(2)\).
\end{remark}

Defining dimensionless generators $L_i^\pm$ by
\[
\frac{\hat J_i}{\hbar}=L_i^+ + L_i^-,
\qquad
\hat P_i = \frac{\hbar}{\Rmom}\bigl(L_i^+ - L_i^-\bigr),
\]
the commutation relations become
\begin{equation}
\label{p2:eq:pp-correct}
[\hat P_i,\hat P_j] = \frac{i\hbar}{\Rmom^2}\,\epsilon_{ijk}\hat J_k,
\qquad
[\hat J_i,\hat P_j] = i\hbar\,\epsilon_{ijk}\hat P_k.
\end{equation}

\subsubsection{Bounded momentum components and the finite equivariant module package}
\label{p2:sec:fuzzy-truncation}

Paper~I-A~\cite{Kugelmann2026a} introduces a mathematical finite-dimensional
\(\Spin(4)\)-equivariant family, termed the ``finite Toeplitz \(S^3\)''
regulator.  Proposition~\ref{p2:prop:algebra-sector-grammar} supplies its
sector grammar.  The structural momentum-role carrier type-support is fixed by
Theorem~\ref{p2:thm:momentum-host-H1}, and conditional carrier rigidity of GP3-irreducible
endpoints is supplied by Theorem~\ref{p2b:thm:momentum-carrier-rigidity}; the optional
\(\mathsf P_{\rm sli}\) package merely organizes a larger lower-ideal tower,
and \(\mathsf C_{\rm PI}\) separately equips chosen members with normalized
Toeplitz coordinate maps.  Only diagonal subspaces are scalar Peter--Weyl
cutoffs; non-diagonal sectors have the module/bundle interpretation stated
below.

\begin{proposition}[Component Spectrum]
\label{p2:prop:component_spectrum}
Let $(j_L,j_R)$ be an irreducible $\Spin(4)$ sector, i.e., a tensor product $V_{j_L}\otimes V_{j_R}$ with $j_L,j_R\in\frac12\ZZ_{\ge0}$. Then:
\begin{enumerate}[label=(\roman*), leftmargin=*, itemsep=6pt]
    \item \textbf{Eigenvalues of $\hat P_3$:}
    \begin{equation}
    \begin{aligned}
    \Spec(\hat P_3)\big|_{(j_L,j_R)}
    =\Bigl\{&\frac{\hbar}{\Rmom}(m_L-m_R)\;\Big|\\
             &m_L\in\{-j_L,\ldots,j_L\},\quad
              m_R\in\{-j_R,\ldots,j_R\}\Bigr\}.
    \end{aligned}
    \end{equation}
    \item \textbf{Operator norm:}
    \begin{equation}
    \|\hat P_3\|_{(j_L,j_R)}=
    \frac{\hbar}{\Rmom}(j_L+j_R)=\Punit(j_L+j_R).
    \end{equation}
    \item \textbf{Rotational symmetry:}
    \begin{equation}
    \|\hat P_i\|_{(j_L,j_R)}=\Punit(j_L+j_R),\qquad i=1,2,3.
    \end{equation}
    \item \textbf{\(\Spin(4)\)-Casimir / kinematic kinetic eigenvalue:}
    \begin{equation}
    \begin{aligned}
    \mathscr C_{\rm kin}
    &:=\frac{2}{\Rmom^2}\sum_{i=1}^3
       \bigl[(L_i^+)^2+(L_i^-)^2\bigr],\\
    \lambda_{j_L,j_R}
    &=\frac{2}{\Rmom^2}\bigl[j_L(j_L+1)+j_R(j_R+1)\bigr].
    \end{aligned}
    \end{equation}
\end{enumerate}
\end{proposition}

\begin{proof}
The weight basis diagonalizes $L_3^\pm$, and
$\hat P_3=\Punit(L_3^+-L_3^-)$ gives the spectrum directly.  The extremal
weights give the norm, and rotational symmetry implies the same norm for
each component.  The displayed kinetic eigenvalue belongs to the
\(\Spin(4)\)-Casimir \(\mathscr C_{\rm kin}\).  Since the two spin triples
commute and \(\hat J_i=\hbar(L_i^++L_i^-)\),
\[
 \sum_{i=1}^3\hat P_i^2
 =\hbar^2\mathscr C_{\rm kin}
  -\frac{1}{\Rmom^2}\sum_{i=1}^3\hat J_i^2.
\]
Thus the Casimir must be distinguished from the transvection-square operator
on both diagonal and non-diagonal sectors.  On scalar diagonal sectors
\((j,j)\), \(\mathscr C_{\rm kin}\) agrees with the positive round-$S^3$
scalar Laplacian; this does not remove the displayed angular-momentum term.
\end{proof}

\begin{remark}[Scalar Laplacian versus full module Casimir]
\label{p2:rem:scalar-vs-full-carrier-casimir}
The scalar Hilbert space \(L^2(S^3)\) contains only diagonal Peter--Weyl
sectors \((j,j)\).  With the integer level convention used here, its finite
subspace is
\[
 \mathcal H_m^{\rm sc}
 =
 \bigoplus_{\substack{j\in\frac12\ZZ_{\ge0}\\2j\le m}}
 V_j\otimes V_j^*.
\]
Since every \(SU(2)\) irreducible is self-dual, \(V_j\otimes V_j^*\) is the
\((j,j)\) sector \(V_j\otimes V_j\) of the labelling used elsewhere in this
paper; the dual notation records the Peter--Weyl origin.
The integer coordinate-stable module family
\(\mathcal H_m=\bigoplus_{j_L+j_R\le m}V_{j_L}\otimes V_{j_R}\)
also keeps non-diagonal \(\Spin(4)\) sectors.  Those sectors are not scalar
functions generated by \(S^3\); they require specified bundle-valued field
types or remain abstract symmetry modules.  In particular,
\[
 \dim\mathcal H_1^{\rm sc}=5,
 \qquad
 \dim\mathcal H_1=15.
\]
On a non-diagonal sector the displayed second-order quantity is the chosen
\(\Spin(4)\)-Casimir kinetic normalization, not a scalar Laplacian eigenvalue.
\end{remark}

\begin{remark}[Constitution versus scale and kinetic interpretation]
\label{p2b:rem:carrier-scale-kinetic-firewall}
The derived GP2 structural type-support \(\mathcal G^{\rm car}_{\Sigma_p}\cong\mathcal H_1\) and \(m_p=1\)
do not determine a numerical value of \(\Rmom\), do not identify the protected
coordinate scale with the parent radius, do not constitute a physical kinetic
Hamiltonian, and do not actualize a measurement or calibration token.  In
particular,
\[
 \boxed{
 \mathcal G^{\rm car}_{\Sigma_p},\ m_p
 \ \not\Rightarrow\
 \Rmom^{\rm numerical},\
 \rho_{\rm prot}=\Rmom,\
 H_{\rm kin},\
 \text{measurement/calibration token}.}
\]
The component spectrum and the displayed Casimir are parameterized kinematic
mathematics until a later role supplies the corresponding scale, dynamical
operator, or token occurrence.
\end{remark}

\subsubsection{Intrinsic spinor seed, canonical host images, and intrinsic word length}
\label{p2:subsec:constructor-grammar-word-length}

Let \(S_+\) and \(S_-\) denote the two intrinsic \(k=0\) chiral
Dirac/Killing-spinor spaces of the round host, and put
\[
 F_L:=S_+,\qquad F_R:=S_-,\qquad S:=S_+\oplus S_-.
\]

\begin{theorem}[Canonical lowest spinor seed and weight classification]
\label{p2b:thm:seed-atomicity}
For the round momentum parent and its derived universal lift \(\Spin(4)\), the
intrinsic lowest two chiral Dirac/Killing-spinor spaces supplied by the round
spin geometry satisfy
\[
 \boxed{S_+\cong V_{\frac12,0},\qquad S_-\cong V_{0,\frac12}.}
\]
Their direct sum \(S=S_+\oplus S_-\) is faithful and the unordered pair is
stable under factor interchange.  After this geometric generation, their
highest weights are the two fundamental weights \(\omega_L,\omega_R\), which
are the atoms of the mathematical dominant-weight monoid
\[
 P^+(\Spin(4))\cong\mathbb Z_{\ge0}\omega_L\oplus
 \mathbb Z_{\ge0}\omega_R.
\]
This last statement classifies the already generated spinor seed; it does not
identify the physical construction monoid with the whole dominant-weight
monoid and is not the selector of \(S_\pm\).
\end{theorem}

\begin{proof}
The \(k=0\) cases of the Paper-I-A spinor harmonic theorem give the two intrinsic
lowest chiral spaces and their \(\Spin(4)\)-types.  The center of one factor
acts nontrivially on one chiral module and trivially on the other, so the
kernels of the two factor actions intersect trivially and \(S\) is faithful.
Factor interchange exchanges the two chiral spaces.  The final dominant-weight
statement is the standard classification of irreducible
\(SU(2)_L\times SU(2)_R\)-modules and is used here only after the geometric
objects have been generated.
\end{proof}

\begin{remark}[Word monoid versus physical constructor monoid]
\label{p2b:rem:seed-semantic-revision}
For representation-theoretic bookkeeping define the free word monoid
\[
 M_{\rm word}:=\mathbb N[F_L]\oplus\mathbb N[F_R].
\]
It records tensor-word length built from the already generated spinor seed.
It is not asserted to equal the GP2 physical construction monoid of the
momentum host.  In particular, the former non-atomic seed countermodels remain
useful diagnostics of weaker generation-based definitions, but they do not
compete with the geometrically fixed lowest spinor spaces above.  Their
detailed comparison is deferred to Paper~II-F, Section~\ref{p2b:app:legacy-selector}.
\end{remark}

\begin{lemma}[Minimal tensor-word length]
\label{p2b:lem:minimal-word-length}
For an irreducible \(V_{j_L,j_R}\), define the representation-theoretic word
length
\[
 \ell(V_{j_L,j_R})
 :=\min\{r\ge0:V_{j_L,j_R}\prec(F_L\oplus F_R)^{\otimes r}\}.
\]
Then
\[
 \boxed{\ell(V_{j_L,j_R})=2j_L+2j_R.}
\]
\end{lemma}

\begin{proof}
A word with \(r_L\) left and \(r_R\) right fundamentals has every constituent
satisfying \(j_L\le r_L/2\) and \(j_R\le r_R/2\), so
\(r\ge2j_L+2j_R\).  Conversely, \(V_{j_L}\) occurs as the highest-spin
constituent of \((V_{1/2})^{\otimes2j_L}\), and similarly on the right.
\end{proof}

\begin{corollary}[Intrinsic word length and half-level normalization]
\label{p2b:cor:word-degree-derived}
The half-level convention is the derived representation label
\[
 m(V_{j_L,j_R})=\frac12\ell(V_{j_L,j_R})=j_L+j_R.
\]
At fixed integer word length \(N\), the multiplicity-free shell and cumulative
prefix have dimensions
\[
 \binom{N+3}{3},\qquad \binom{N+4}{4},
\]
respectively.  This is classification data; it does not make unrestricted
tensor powers physical GP2-generated content.
\end{corollary}

\begin{proof}
The length identity is Lemma~\ref{p2b:lem:minimal-word-length}.  At fixed
\(N=a+b\), sum \((a+1)(b+1)\) over \(a+b=N\), then sum the shells.
\end{proof}

\section{Generated six-type momentum host and GP3 comparison}
\label{paper2b:sec:generated-six-type-host}
\label{p2b:subsec:unconditional-momentum-host}
\subsection{Intrinsic bracket-image constructor and typed closure}

Paper~II-A freezes the initial momentum-seed dependency frontier in
Definition~\ref{p2:def:momentum-seed-frontier}.  This paper verifies its
round-host provenance.  Let \(E_k\) denote the scalar harmonic level \(k\) of
the round parent.  By the Paper-I-A representation-harmonic package,
\[
 E_0\cong(0,0),\qquad E_1\cong(\tfrac12,\tfrac12),
\]
and the intrinsic lowest two spinor objects are
\(S_+\cong(\tfrac12,0)\) and \(S_-\cong(0,\tfrac12)\).  The round gradient and
the vector-field Lie bracket are host-natural constructors available before
any finite representation level is chosen.  By Paper~II-A's parent-geometric
object typing, \(\Gamma(TS^3)_{\rm sec,obj}\) is the tangent-section object
sort of the same momentum parent, whereas the induced derivation action is a
distinct operation-typed output co-realized by the same licensed constructor
word.  No tensor-word boundary,
compression map, component ceiling, or pointwise scalar multiplication as a
finite carrier constructor belongs to this seed frontier.

\begin{lemma}[Canonical bracket linearization and bivector--Killing bridge]
\label{p2b:lem:canonical-bracket-image}
Identify the real first harmonic space \(E_1^{\mathbb R}\) with the ambient
vector space \(\mathbb R^4\) by
\[
 u\longmapsto f_u,\qquad f_u(x):=\langle u,x\rangle,
 \qquad x\in S^3(\Rmom).
\]
Then
\[
 \nabla f_u=u-\frac{f_u}{\Rmom^2}x
\]
and
\[
 [\nabla f_u,\nabla f_v]
 =\frac1{\Rmom^2}\bigl(f_u v-f_v u\bigr).
\]
Thus the constituted bilinear constructor
\[
 \beta:E_1^{\mathbb R}\times E_1^{\mathbb R}
 \longrightarrow\Gamma(TS^3)_{\rm sec,obj},\qquad
 \beta(f_u,f_v):=[\nabla f_u,\nabla f_v],
\]
has a canonical real linear hull.  Because \(\beta\) is alternating, its
unique mathematical linearization is
\[
 \bar\beta:\Lambda^2E_1^{\mathbb R}
 \longrightarrow\Gamma(TS^3)_{\rm sec,obj},
 \qquad
 \bar\beta(u\wedge v):=[\nabla f_u,\nabla f_v].
\]
It satisfies
\[
 Z_\beta:=\operatorname{span}_{\mathbb R}\operatorname{im}\beta
 =\operatorname{im}\bar\beta
 =\mathcal K_{\rm Kill}.
\]
Equivalently,
\[
 \Theta:=\Rmom^2\bar\beta:
 \Lambda^2E_1^{\mathbb R}\longrightarrow\mathcal K_{\rm Kill},
 \qquad
 \Theta(u\wedge v):=\Rmom^2[\nabla f_u,\nabla f_v]
\]
is a canonical \(SO(4)\)-equivariant isomorphism onto the six-dimensional
Killing-field module.  After complexification,
\[
 (\mathcal K_{\rm Kill})_{\mathbb C}
 \cong\Lambda^2_{\mathbb C}
 \cong(1,0)\oplus(0,1).
\]
\end{lemma}

\begin{proof}
The displayed gradient is the orthogonal projection of the constant ambient
vector \(u\) to \(T_xS^3\).  Differentiating it in the ambient space and
antisymmetrizing gives the bracket formula.  The vector field
\[
 x\longmapsto\langle u,x\rangle v-\langle v,x\rangle u
\]
is the infinitesimal rotation associated with the ambient bivector
\(u\wedge v\).  The standard map \(\Lambda^2\mathbb R^4\to\mathfrak{so}(4)\)
is an isomorphism, so \(\Theta\), and hence \(\bar\beta\), is injective and
has image all of \(\mathcal K_{\rm Kill}\).  Since the image of the induced
linear map is exactly the real span of the outputs of the original bilinear
map, this proves the displayed identity for \(Z_\beta\).  Naturality of the
ambient inner product makes \(\Theta\) equivariant.  The Paper-I-A
quadratic-shell theorem gives the same linear-image identity and the displayed
\(\Spin(4)\)-type.
\end{proof}


\begin{lemma}[Canonical normalized first-harmonic gradient port]
\label{p2b:lem:gradient-port}
On the round host, the gradient is a natural, \(\Spin(4)\)-equivariant
isomorphism of object modules
\[
 \nabla:E_1\xrightarrow{\ \sim\ }\nabla E_1
 \subset\Gamma(TS^3)_{\rm sec,obj}.
\]
For the round \(L^2\) scalar and vector-field inner products,
\[
 \langle\nabla f,\nabla g\rangle_{L^2}
 =\frac3{\Rmom^2}\langle f,g\rangle_{L^2},\qquad
 U_\nabla:=\frac{\Rmom}{\sqrt3}\nabla,
 \qquad U_\nabla^{-1}=-\frac{\Rmom}{\sqrt3}\operatorname{div}
 \quad\text{on }\nabla E_1.
\]
This unitary object-port map identifies the two isomorphic four-dimensional
representations inside a one-copy presentation.  Their source/target roles
remain distinct; it does not identify a section object with an induced
operation occurrence or add a physical gradient-evaluation constructor.
\end{lemma}
\begin{proof}
With the positive Laplacian convention
\(\Delta=-\operatorname{div}\nabla\), every \(f\in E_1\) satisfies
\(\Delta f=3f/\Rmom^2\).  Integration by parts on the compact round
sphere gives the inner-product identity and the displayed inverse.
The metric and Laplacian commute with the derived \(\Spin(4)\) action,
so the map is equivariant and independent of a chosen basis.
\end{proof}

\begin{proposition}[Conformal-field Lie closure, in the smooth host]
\label{p2b:prop:conformal-host-closure}
For \(u,v\in\mathbb R^4\) put \(G_u:=\nabla f_u\), and for
\(A\in\mathfrak{so}(4)\) put \(K_A(x):=Ax\).  Their ordinary
vector-field brackets obey
\[
 [G_u,G_v]=K_{(vu^\top-uv^\top)/\Rmom^2},\quad
 [K_A,G_u]=-G_{Au},\quad
 [K_A,K_B]=K_{BA-AB}.
\]
Consequently
\(\mathfrak c:=\operatorname{span}_{\mathbb R}\{G_u,K_A\}\)
is a ten-dimensional Lie algebra isomorphic to \(\mathfrak{so}(4,1)\).
As a representation of the \emph{compact} derived \(\Spin(4)\),
\[
 \mathfrak c_{\mathbb C}
 \cong(\tfrac12,\tfrac12)\oplus(1,0)\oplus(0,1)
 \cong\operatorname{Sym}^2(S_+\oplus S_-),
 \qquad \mathcal H_1\cong\mathbf1\oplus S\oplus\mathfrak c_{\mathbb C}.
\]
Thus \(15=1+4+10\) is an intrinsic host-classification of the
already generated six types.  Neither \(\mathfrak{so}(4,1)\) nor
its noncompact finite-dimensional action is introduced as a new physical
momentum group or a faithful unitary finite-carrier operation constructor.
\end{proposition}
\begin{proof}
The first bracket is Lemma~\ref{p2b:lem:canonical-bracket-image}.
Direct differentiation of \(K_A\) and \(G_u\) gives the other two.
The injective map
\[
 K_A\longmapsto\begin{pmatrix}-A&0\\0&0\end{pmatrix},\qquad
 G_u\longmapsto\begin{pmatrix}0&u/\Rmom\\u^\top/\Rmom&0\end{pmatrix}
\]
preserves the three brackets, and its image is the standard
\(\mathfrak{so}(4,1)\) block presentation.  The four gradients have
\((\tfrac12,\tfrac12)\)-type; the six Killing fields decompose as
\((1,0)\oplus(0,1)\).  Expanding the symmetric square of the two
chiral defining modules proves the final representation identity.
\end{proof}

\begin{remark}[Coordinate-algebra blocks versus geometric type blocks]
\label{p2b:rem:conformal-block-compatibility}
The geometric decomposition
\(\mathbf1\oplus\mathfrak c_{\mathbb C}\oplus S\)
organizes \(\mathcal H_1\) into dimensions \(11+4\).
Where the separately constituted intrinsic coordinate tuple and its
completion hypotheses produce
\(\mathcal A_D\cong M_{11}(\mathbb C)\oplus M_4(\mathbb C)\),
those block sizes agree with the geometric organization.  Matrix
simplicity and the dimension \(137=11^2+4^2\) depend on the separate
coordinate-intertwiner/word calculation, not on the decomposition
\(11+4\) alone, and do not select a physical operation or state.
\end{remark}

\begin{proposition}[Maximal exact unital scalar derivation subalgebra]
\label{p2b:prop:maximal-finite-scalar-action}
Set \(\mathcal P_1:=E_0\oplus E_1\subset C^\infty(S^3(\Rmom))\).
For a smooth-host field \(X\in\mathfrak c\), let
\(\delta_X(h)=X(h)\) be its \emph{ordinary} scalar derivation.  Then
\[
 \delta_X(\mathcal P_1)\subseteq\mathcal P_1
 \quad\Longleftrightarrow\quad X\in\mathcal K_{\rm Kill}.
\]
This is the maximal exact action \emph{among fields in}
\(\mathfrak c\), conditional only on taking \(\mathcal P_1\) as the
specified scalar-operation target; the proposition does not assert that
the physical operation grammar has selected that target.
\end{proposition}
\begin{proof}
Write \(X=K_A+G_u\).  A Killing field preserves \(E_0\) and
\(E_1\), since \(K_A(f_v)=-f_{Av}\).
On the other hand,
\[
 G_u(f_v)=u\cdot v-\frac{f_uf_v}{\Rmom^2}.
\]
If \(u\ne0\), take \(v=u\).  The traceless quadratic harmonic part of
\(f_u^2\) is nonzero; the Killing term has degree one and cannot
cancel it.  Thus \(X\) fails to preserve \(\mathcal P_1\).
\end{proof}

\begin{proposition}[Conditional infinite tower under retained gradient evaluation]
\label{p2b:prop:gradient-evaluation-no-go}
Consider a \emph{different, enlarged} physical grammar that (i) licenses
ordinary scalar evaluation by \(G=\nabla f\), \(f=x_1\),
(ii) retains each output in a scalar carrier sort, and (iii) licenses
repeated evaluation on those outputs.  Its scalar carrier closure is
infinite-dimensional.
\end{proposition}
\begin{proof}
Let \(F_0=f\), \(F_{n+1}=G(F_n)\).  On functions of \(f\),
\(G=(1-f^2/\Rmom^2)d/df\).  Inductively \(F_n\) has degree
\(n+1\) with leading coefficient
\((-1)^n n!/\Rmom^{2n}\ne0\).  Therefore the
\(F_n\) are linearly independent as functions on the sphere.  This
is an obstruction only to the explicitly enlarged retention/reuse grammar,
not to the frozen Paper-II-A momentum seed.
\end{proof}

\begin{remark}[Typed constructor and evaluation boundary]
\label{p2b:rem:typed-momentum-operation-grammar}
The frozen seed licenses the \emph{object} words
\(E_1\xrightarrow{\nabla}\Gamma(TS^3)_{\rm sec,obj}\) and
\(E_1\times E_1\xrightarrow{\beta}\Gamma(TS^3)_{\rm sec,obj}\),
with section-object linear-span retention for \(Z_\beta\).  An induced
operation \(K\mapsto\delta_K\) remains a distinct port and is physically
co-realized only on the action target and evaluations explicitly licensed by
the parent constructor signature.  Neither ambient pointwise multiplication,
nor arbitrary \(G(h)\), nor reuse of a scalar evaluation output as a new
finite-carrier seed follows from these object words.  The host Lie closure
of Proposition~\ref{p2b:prop:conformal-host-closure} is a mathematical
classification, not a recursive enlargement of the physical constructor
signature.  A conformal-density Lie derivative would change the typed
operation and is not silently substituted for \(\delta_X\).
\end{remark}

\begin{remark}[Literal bilinear image versus generated linear image]
\label{p2b:rem:beta-set-image-linear-hull}
The literal set image \(\operatorname{im}\beta\) consists of fields
corresponding to decomposable bivectors \(u\wedge v\) and is a proper subset
of \(\mathcal K_{\rm Kill}\).  For example,
\(e_1\wedge e_2+e_3\wedge e_4\) is not decomposable because its exterior
square is nonzero.  No equality
\(\operatorname{im}\beta=\mathcal K_{\rm Kill}\) is claimed.  The
six-dimensional object used below is the generated real linear subspace
\(Z_\beta=\operatorname{span}_{\mathbb R}\operatorname{im}\beta\).
\end{remark}

\begin{lemma}[Sharp literal two-bracket realization of every round Killing section]
\label{p2b:lem:v264-two-bracket-surjection}
On the declared round momentum host, with the \emph{already constituted}
addition in the real vector-space section-object target, the composite
\[
 \mathfrak B:(E_1^{\mathbb R})^4\longrightarrow
 \Gamma(TS^3)_{\rm sec,obj},\qquad
 \mathfrak B(f,g,h,k)=[\nabla f,\nabla g]+[\nabla h,\nabla k]
\]
has literal set image precisely $\mathcal K_{\rm Kill}$.  Under the
isomorphism $\Theta:\Lambda^2E_1^{\mathbb R}\to\mathcal K_{\rm Kill}$
of Lemma~\ref{p2b:lem:canonical-bracket-image}, a Killing section
$\Theta(A)$ belongs to the \emph{single}-bracket set image exactly when
\[
 A\wedge A=0,
 \quad\text{equivalently}\quad
 \|A_+\|=\|A_-\|
\]
for the self-dual/anti-self-dual decomposition of real two-forms.
The bound of two brackets is sharp.  This uses only finite composition
and addition in the frozen target object sort; $\Lambda^2E_1$ is a
\emph{classification} device, not a newly physicalized constructor.
\end{lemma}
\begin{proof}
Every real alternating $4\times4$ matrix admits an orthonormal skew
normal form $A=\lambda e_1\wedge e_2+\mu e_3\wedge e_4$.
Since $u\mapsto f_u$ is real linear, each coefficient can be absorbed
into one harmonic argument, so $\Theta(A)$ is the sum of two bracket
outputs. Conversely each such sum is Killing by the preceding lemma.
A nonzero real bivector in dimension four is decomposable if and only if
$A\wedge A=0$ (the vanishing Pfaffian criterion). As
$A\wedge A=(\|A_+\|^2-\|A_-\|^2)\operatorname{vol}$ with the same
Hodge normalization on the two halves, the norm criterion is equivalent.
The form $e_1\wedge e_2+e_3\wedge e_4$ is not decomposable, proving
that one bracket cannot always suffice.
\end{proof}
\begin{remark}[Finite-image inheritance without a span/image ambiguity]
\label{p2b:rem:v264-literal-inheritance}
The preceding literal surjection makes the premise of
Lemma~\ref{p2:lem:gp2-finite-generated-image-inheritance} directly
applicable to $\operatorname{im}\mathfrak B$ whenever the previously
frozen vector-space section-object addition is consumed. It strengthens
Corollary~\ref{p2b:cor:bracket-image-same-carrier}; it does not infer
that addition is physical merely from a computed image and does not
supply the distinct GP3 whole-candidate comparison certificate.
\end{remark}

\begin{corollary}[Inherited same-role typing of the canonical bracket linear image]
\label{p2b:cor:bracket-image-same-carrier}
The finite-dimensional real section-object subspace
\[
 Z_\beta
 =\operatorname{span}_{\mathbb R}\operatorname{im}\beta
 =\operatorname{im}\bar\beta
 =\mathcal K_{\rm Kill}
\]
belongs to the momentum role's generated object closure.  Its complexification
\((Z_\beta)_{\mathbb C}=(\mathcal K_{\rm Kill})_{\mathbb C}\) belongs to the
momentum carrier.  The corresponding infinitesimal Killing action remains
a distinct operation-typed node and is co-realized, not independently
actualized.
\end{corollary}

\begin{proof}
The constituted composite
\[
 \beta:E_1^{\mathbb R}\times E_1^{\mathbb R}
 \longrightarrow\Gamma(TS^3)_{\rm sec,obj}
\]
lands in the same-role parent geometric object sort fixed by Paper~II-A.
That target sort is a real vector-space object sort, so its already
constituted addition and scalar multiplication force every generated object
subspace containing all \(\beta\)-outputs to contain their real linear hull
\(Z_\beta\).  This is closure within the target's object typing; it neither
constitutes \(\Lambda^2E_1^{\mathbb R}\) as a new physical constructor nor
adds arbitrary tensor or polynomial closure.

By Lemma~\ref{p2b:lem:canonical-bracket-image}, the resulting linear image
\(Z_\beta=\operatorname{im}\bar\beta=\mathcal K_{\rm Kill}\) is
finite-dimensional, candidate-independent, intrinsically characterized,
nonredundant, invariant, and carries the induced \(\Spin(4)\) transformation
law.  Paper~II-A's constructor-relative generated closure and finite
generated-image inheritance
Lemma~\ref{p2:lem:gp2-finite-generated-image-inheritance} therefore preserve
the section-object typing of \(Z_\beta\); retaining only its induced
derivation/Killing-operation occurrence would violate GP2 typed host
faithfulness and no-omission.

The frozen frontier generates no autonomous role label, separate carrier,
cross-role variance, or separating test that would split this inherited
linear object image from the momentum role.  Paper~II-A's same-role
no-splitting Lemma~\ref{p2:lem:gp2-same-role-no-splitting} therefore places
its finite-dimensional object copy, after complexification, in the momentum
carrier.  GP2 sort faithfulness keeps the induced operation occurrence
separate.  The conclusion objectifies only this canonical generated linear
subspace, not the ambient infinite-dimensional section space
\(\Gamma(TS^3)\).
\end{proof}

\subsection{Exact GP2 type-support and multiplicity-blind observables}

\begin{theorem}[Exact momentum carrier type-support]
\label{p2:thm:momentum-host-H1}
Relative to the frozen momentum-seed frontier of Paper~II-A, the GP2 generated
carrier type-support is exactly
\[
 \begin{aligned}
 \mathcal G^{\rm car}_{\Sigma_p}
 &=E_0\oplus S_+\oplus S_-\oplus E_1
   \oplus(\mathcal K_{\rm Kill})_{\mathbb C}\\
 &\cong
 (0,0)\oplus(\tfrac12,0)\oplus(0,\tfrac12)
 \oplus(\tfrac12,\tfrac12)\oplus(1,0)\oplus(0,1)\\
 &=\boxed{\mathcal H_1}.
 \end{aligned}
\]
Equivalently,
\[
 \mathcal H_1\cong\mathbf1\oplus S\oplus\operatorname{Sym}^2S,
 \qquad \dim_{\mathbb C}\mathcal H_1=15,
\]
where the symmetric-square identity classifies the already generated content;
\(\operatorname{Sym}^2\) is not invoked as a physical constructor.

Thus every GP2-exact/host-faithful finite momentum realization has precisely
this irreducible carrier type-support, up to unmarked implementation
multiplicity.  GP2 no-omission excludes proper supports such as the former
\(\mathcal E_9\) countermodel, while constructor-relative no-insertion
excludes higher irreducible types not produced by the frozen seed frontier.
\end{theorem}

\begin{proof}
The scalar, spinor, and first-coordinate sectors are the verified frontier
objects above.  Corollary~\ref{p2b:cor:bracket-image-same-carrier} supplies the
complexified generated linear object sector
\((Z_\beta)_{\mathbb C}=(\mathcal K_{\rm Kill})_{\mathbb C}
\cong(1,0)\oplus(0,1)\).  Assembling
these independently generated invariant sectors is the exact type-support
required by the host signature, not a new direct-sum constructor.
Paper~I-A Proposition~\ref{p1:prop:h1-symmetric-exterior-form}
identifies the resulting module with \(\mathcal H_1\) and gives its dimension.

No other carrier-generating constructor belongs to the frozen momentum-seed
frontier of Definition~\ref{p2:def:momentum-seed-frontier}.  In particular, arbitrary tensor powers and ordinary pointwise
multiplication of coordinate harmonics remain parent mathematics rather than
finite carrier generation at this stage.  Hence their possible outputs, such
as the symmetric \((1,1)\) quadratic sector, are not inserted.  Conversely,
omitting any displayed generated sector violates GP2 exactness or sort
faithfulness.  In particular, the former \(\mathcal E_9\) branch is not a
complete GP2 realization: it retains the induced Killing operations but omits
the independently generated finite-dimensional linear section-object subspace required by
Corollary~\ref{p2b:cor:bracket-image-same-carrier}.
\end{proof}

\begin{corollary}[Exact multiplicity detectors for the frozen six-type momentum support]
\label{p2b:cor:v279-six-type-multiplicity-detectors}
Let \(\mathcal H\) be any finite-dimensional complex \(\Spin(4)\)-module whose
irreducible \emph{type-support} is exactly the six-type set \(T_p\) of
Theorem~\ref{p2:thm:momentum-host-H1}.  Write
\[
 \mathcal H=\bigoplus_{\lambda\in T_p}V_\lambda\otimes\mathbb C^{m_\lambda},
 \qquad m_\lambda\in\mathbb Z_{\ge1},\qquad
 \mathcal H_1=\bigoplus_{\lambda\in T_p}V_\lambda.
\]
With the six type dimensions \(d_\lambda=1,2,2,4,3,3\), respectively,
the exact carrier-dimension excess is
\[
 \boxed{\Delta_{\rm dim}:=\dim_{\mathbb C}\mathcal H-15
   =\sum_{\lambda\in T_p}d_\lambda(m_\lambda-1)\ge0.}
\]
The equivariant-commutant excess and normalized character mismatch are
\[
 \boxed{\begin{aligned}
 \Delta_{\rm com}&:=\dim_{\mathbb C}\End_{\Spin(4)}(\mathcal H)-6
                   =\sum_{\lambda\in T_p}(m_\lambda^2-1)\ge0,\\
 \Delta_{\chi}&:=\int_{\Spin(4)}
                |\chi_{\mathcal H}(g)-\chi_{\mathcal H_1}(g)|^2\,dg
                   =\sum_{\lambda\in T_p}(m_\lambda-1)^2\ge0.
 \end{aligned}}
\]
Here \(dg\) is normalized Haar measure.  The three conditions
\(\Delta_{\rm dim}=0\), \(\Delta_{\rm com}=0\), and
\(\Delta_\chi=0\) are equivalent to \(m_\lambda=1\) for every
\(\lambda\in T_p\).  Thus an \emph{independently established} exact value
\(\dim\mathcal H=15\), commutant dimension six, or vanishing character
mismatch certifies multiplicity one for that particular carrier without
requiring the universal whole-candidate GP3 comparison of
Theorem~\ref{p2b:thm:momentum-carrier-rigidity}.

These are conditional mathematical diagnostics of an already specified
\(\Spin(4)\)-module, \emph{not} new GP2 constructors, an automatically
realized character measurement, nor a proof that the physical momentum
endpoint universally attains any of the three equality values.
\end{corollary}
\begin{proof}
The six representation dimensions in Theorem~\ref{p2:thm:momentum-host-H1} give the first identity.  The
\(\Spin(4)\) commutant decomposition by Schur's lemma is
\(\End_{\Spin(4)}(\mathcal H)\cong\bigoplus_\lambda M_{m_\lambda}(\mathbb C)\),
which gives the second.  The irreducible characters are orthonormal in
\(L^2(\Spin(4),dg)\), hence
\(\chi_{\mathcal H}-\chi_{\mathcal H_1}
 =\sum_\lambda(m_\lambda-1)\chi_\lambda\) proves the third.
Each summand in each expression is nonnegative, and all coefficients
\(d_\lambda\) are strictly positive; therefore any one zero defect
forces every \(m_\lambda-1=0\).  No implication invokes GP3 leastness or
constructs a physical trace/measurement merely from this identity.
\end{proof}

\begin{proposition}[Unmarked multiplicity does not supply a canonical copy selector]
\label{p2b:prop:v279-no-canonical-copy-selector}
Let \(V\ne0\) be an irreducible \(\Spin(4)\)-module and let
\(M\cong\mathbb C^m\) with \(m>1\). Suppose the frozen structure on
\(V\otimes M\) carries no distinguished line, subspace, or multiplicity
label, so every \(I_V\otimes U\) for \(U\in U(M)\) is a symmetry of the
\emph{multiplicity-slot data}.  Then no one-dimensional complex subspace
of \(M\) is invariant under all these symmetries, and there is no
nonzero \(U(M)\)-invariant linear functional \(M\to\mathbb C\) (with
trivial target action).  In particular a projector choosing a single
copy \(V\otimes\mathbb C\subset V\otimes M\), or a nonzero linear
retraction to such an unmarked copy, is not canonically generated
from the multiplicity-slot data alone.

This obstruction concerns \emph{copy selection/compression}. It does not
exclude a separately proved generated whole-structured GP3 comparison
between distinct admissible candidates; that is exactly the independent
certificate of Theorem~\ref{p2b:thm:momentum-carrier-rigidity}.
\end{proposition}
\begin{proof}
The unitary group \(U(m)\) acts transitively on complex lines in
\(M\), so no proper nonzero line is invariant under every unitary.
If \(L:M\to\mathbb C\) were invariant, applying a central unitary
\(e^{it}I_M\) would give
\(L(e^{it}v)=L(v)=e^{it}L(v)\) for every \(t\), hence \(L=0\).
Any equivariant reduction obtained by selecting a copy would require
such an invariant line or a nonzero invariant functional; neither
is supplied by an unmarked multiplicity slot.  The argument does not
address the separate, non-compression notion of a generated comparison
of complete candidate presentations.
\end{proof}

\begin{theorem}[Canonical observable quotient without choosing a multiplicity copy]
\label{p2b:thm:v293-operational-multiplicity-quotient}
Let $K=\Spin(4)$ and let a finite unitary carrier have the isotypic
presentation
\[
 \mathcal H=\bigoplus_{\lambda\in T}V_\lambda\otimes M_\lambda,
 \qquad 1\le\dim M_\lambda=m_\lambda<\infty,
\]
where the $V_\lambda$ are pairwise inequivalent. Assume the specified
\emph{observable test algebra} is precisely the multiplicity-blind
unital $C^*$-algebra
\[
 \mathcal A_{\rm bl}
 =\bigoplus_{\lambda\in T}
   \bigl(\operatorname{End}(V_\lambda)\otimes I_{M_\lambda}\bigr)
 \subset\operatorname{End}(\mathcal H).
\]
Let $\mathcal H_0=\bigoplus_{\lambda\in T}V_\lambda$ and
$\mathcal A_0=\bigoplus_\lambda\operatorname{End}(V_\lambda)$,
represented block diagonally on $\mathcal H_0$. Then the following
statements hold, canonically relative to the isotypic presentation (and
independently of every multiplicity-space basis):
\begin{enumerate}[label=(\roman*),leftmargin=2em]
\item The map $j:\mathcal A_0\to\mathcal A_{\rm bl}$,
$j((a_\lambda))=\bigoplus_\lambda a_\lambda\otimes I_{M_\lambda}$,
is a unital $*$-isomorphism.
\item Every density matrix $\rho$ on $\mathcal H$ induces the
block-diagonal density
\[
 Q(\rho)=\bigoplus_\lambda
 \operatorname{Tr}_{M_\lambda}(P_\lambda\rho P_\lambda)
 \quad\text{on }\mathcal H_0,
\]
with $P_\lambda$ the isotypic projector, and for all
$a\in\mathcal A_0$,
\[
 \operatorname{Tr}_{\mathcal H}(\rho j(a))
 =\operatorname{Tr}_{\mathcal H_0}(Q(\rho)a).
\]
The map $Q$ is an affine surjection onto the block-diagonal density
matrices. It has the canonical mixed-state right inverse
\[
 S((\sigma_\lambda)_\lambda)
 =\bigoplus_\lambda\sigma_\lambda\otimes\frac{I_{M_\lambda}}{m_\lambda},
 \qquad Q\circ S=\mathrm{id}.
\]
This is a \emph{state/effect quotient}, not a projector choosing a Hilbert
subspace isomorphic to $\mathcal H_0$.
\item A Hermitian effect candidate $E\in\operatorname{End}(\mathcal H)$
passes through $Q$ for \emph{every} density matrix, in the sense that
$\operatorname{Tr}(\rho E)$ depends only on $Q(\rho)$, if and only if
$E\in\mathcal A_{\rm bl}$. Thus the complete multiplicity-blind
observable content has an exact converse test.
\item If a specified exact process has a unital completely positive
Heisenberg map $\Phi^*$ on $\operatorname{End}(\mathcal H)$, its state
transition descends to the above observable quotient for all states
and all effects in $\mathcal A_{\rm bl}$ if and only if
$\Phi^*(\mathcal A_{\rm bl})\subseteq\mathcal A_{\rm bl}$. When this
holds, the induced map
$\widetilde\Phi^*=j^{-1}\circ\Phi^*\circ j$
is uniquely defined and unital completely positive, with
$Q\circ\Phi=\widetilde\Phi\circ Q$ on states.
\end{enumerate}
In particular, observational multiplicity redundancy has an exact
basis-free quotient even when
Proposition~\ref{p2b:prop:v279-no-canonical-copy-selector}
forbids an invariant one-copy \emph{subspace} selector. This theorem does
\emph{not} establish that the quotient is an occurring generated GP2
construction or a proper GP3 comparison of complete candidates: those
claims still require all frozen objects, source labels, transformations,
relations, tests, and occurrence ports to be preserved.
\end{theorem}
\begin{proof}
The direct-sum block map $j$ preserves products, adjoints, the unit,
and operator norms and has inverse obtained by reading one coefficient
from each multiplicity-blind block; hence it is a $*$-isomorphism.
The partial-trace identity gives the stated equality of expectations.
Each $Q(\rho)$ is positive, block diagonal, and has trace one. For any
such $\sigma$ the displayed $S(\sigma)$ is positive of trace one and
$Q(S(\sigma))=\sigma$. The identity on each $M_\lambda$ makes $S$
independent of a multiplicity basis.

For (iii), if $E=j(a)$, the trace identity proves the claim. Conversely,
let $\mathsf P_{\mathcal A}$ be the orthogonal projection onto
$\mathcal A_{\rm bl}$ for the Hilbert--Schmidt real inner product on
Hermitian matrices, and suppose $E\notin\mathcal A_{\rm bl}$. Put
$X=E-\mathsf P_{\mathcal A}E$. Then $X=X^*$, $X\ne0$, and $X$ is
orthogonal to all Hermitian elements of $\mathcal A_{\rm bl}$, including
$I$; thus $\operatorname{Tr}X=0$. For sufficiently small $\varepsilon>0$,
\[
 \rho_\pm=\frac{I_\mathcal H}{\dim\mathcal H}
                 \pm\varepsilon X
\]
are density matrices with identical $Q$-image, while
$\operatorname{Tr}[(\rho_+-\rho_-)E]
=2\varepsilon\operatorname{Tr}(X^2)>0$.

For (iv), invariance of $\mathcal A_{\rm bl}$ defines the displayed
conjugated unital completely positive map; complete positivity is
preserved by the $*$-isomorphism. Duality against every $j(a)$ gives
the state intertwining identity. Conversely, if a descended process
exists for all states and multiplicity-blind readouts, then each
$\rho\mapsto\operatorname{Tr}(\Phi(\rho)j(a))$
depends only on $Q(\rho)$. Part (iii) forces
$\Phi^*(j(a))\in\mathcal A_{\rm bl}$ for every $a$, proving invariance.
No step selects a line in any $M_\lambda$ or supplies an unproved
GP3 whole-candidate comparison.
\end{proof}


\subsection{Derived momentum level and structured GP3 comparison}

\begin{corollary}[Level one from GP2 exact type-support, independently of GP3]
\label{p2:thm:m1-component}
Let \(\mathfrak M\) be any GP1-admissible, GP2-exact finite structured
realization of the frozen momentum-seed frontier, with the six irreducible
types of Theorem~\ref{p2:thm:momentum-host-H1} and arbitrary positive
implementation multiplicities.  Then
\[
 \max_{V_{j_L,j_R}\subset\mathcal H_{\mathfrak M}}
 \ell(V_{j_L,j_R})=2,\qquad
 \boxed{m_p:=\tfrac12\max\ell=1}.
\]
This level is a type-support descriptor and does not determine the total
carrier dimension, multiplicities, states, measurements, or a level selector.
\end{corollary}
\begin{proof}
The six types have word lengths \(0,1,1,2,2,2\) by
Lemma~\ref{p2b:lem:minimal-word-length}.  Repeating a representation
changes neither its highest weight nor its minimal tensor-word length.
\end{proof}

\begin{proposition}[Explicit one-copy geometric module, without GP3]
\label{p2b:prop:one-copy-geometric-module}
There is a fifteen-dimensional \(\Spin(4)\)-equivariant one-copy
\emph{mathematical carrier presentation}
\[
 \mathcal H_1=E_0\oplus S_+\oplus S_-\oplus E_1
 \oplus(\mathcal K_{\rm Kill})_{\mathbb C}
\]
carrying the parent-generated object ports \(\nabla\), \(\beta\), and
their canonical linear-image incidence.  Its abstract representation
and these object diagrams require no GP3 conclusion.  Certification that
this presentation is a \emph{complete GP2-exact structured endpoint}
for every applicable operation, relation and test remains a separate
hypothesis in Theorem~\ref{p2b:thm:momentum-carrier-rigidity}.
\end{proposition}
\begin{proof}
Use the six inequivalent irreducibles supplied by
Theorem~\ref{p2:thm:momentum-host-H1}, each once.  The scalar and section
ports are identified equivariantly by Lemma~\ref{p2b:lem:gradient-port};
the bilinear bracket is fixed by Lemma~\ref{p2b:lem:canonical-bracket-image}.
These maps are independent of implementation multiplicity.  They prove
existence of the displayed mathematical realization, without presupposing
the additional universal structured-comparison property.
\end{proof}

The explicit native coordinate presentation contains auxiliary pair
creation/annihilation and degree-preserving terms. Paper~II-F, Section~\ref{p2b:app:v2174-geometric-pair-presentation}
evaluates that presentation and its scoped linear-port comparison; it
supplies no physical occupation boundary or quantum field constructor.


\begin{remark}[Momentum package and geometric host]
\label{p2b:rem:m8-momentum-package-scope}
The derived \(S_1^3\) denotes the enriched physical momentum package
on smooth support \(S_p^3\).  The latter is its support in the
completed reciprocal host, not a second sphere containing a
smaller momentum sphere.  The fifteen-dimensional \(\mathcal H_1\)
carrier is not replaced by functions on the twenty-four Hurwitz
points.  Finite support atlases and the intrinsic coordinate class
are separately typed descendants; II-C transports the intrinsic
class and its joint comparisons in those precise roles.
\end{remark}


\begin{theorem}[GP3 carrier rigidity under a complete structured-comparison certificate]
\label{p2b:thm:momentum-carrier-rigidity}
Let \(\mathfrak M\) range over GP1-admissible, GP2-exact/host-faithful
structured realizations of the frozen momentum-seed signature, with
\[
 \mathcal H_{\mathfrak M}=\bigoplus_{\lambda\in T_p}
 V_\lambda\otimes M_\lambda,
 \qquad T_p=\{(0,0),(\tfrac12,0),(0,\tfrac12),
 (\tfrac12,\tfrac12),(1,0),(0,1)\}.
\]
Assume the following \emph{additional whole-structured certificates}:
\begin{enumerate}[label=(\roman*),leftmargin=2em]
\item The mathematical one-copy model
\(\mathfrak M_1\) of
Proposition~\ref{p2b:prop:one-copy-geometric-module} has a
complete GP1-admissible, GP2-exact and host-faithful presentation of
\emph{all} frozen object, operation, constructor, relation and test diagrams.
\item For every such \(\mathfrak M\) there exists a \emph{generated}
role-faithful whole-candidate comparison
\(\mathfrak M_1\preceq_{\Sigma_p}\mathfrak M\) satisfying
Paper~II-A Postulate~\ref{p2:post:GP3}.  The witness preserves all
frozen diagrams, including typed section-to-operation incidence,
and does not use a non-generated projector, compression or
candidate-specific label as a selection mark.
\item The comparison is proper whenever
\(\dim M_\lambda>1\) for at least one \(\lambda\).
\end{enumerate}
Then every GP3-irreducible endpoint satisfies
\[
 \boxed{\dim M_\lambda=1\ (\lambda\in T_p),\qquad
 \mathcal H_{\mathfrak M}\cong\mathcal H_1,\quad
 \dim_{\mathbb C}\mathcal H_{\mathfrak M}=15.}
\]
This theorem does not assert uniqueness of the entire structured endpoint.
Its three comparison hypotheses are \emph{not} consequences of type-support
or the absence of frozen copy labels alone.  Until a universal certificate
is supplied, the unconditional conclusion is the common six-type GP2
support and \(m_p=1\), not universal endpoint multiplicity-one.
\end{theorem}
\begin{proof}
GP2 no-omission gives \(\dim M_\lambda\ge1\) for every
\(\lambda\in T_p\).  Under (i)--(ii), the one-copy model is a complete
admissible competitor with a generated comparison to every source candidate.
Under (iii), any source candidate with a surplus multiplicity admits a
proper role-faithful comparison and is not GP3-irreducible.  Hence its
multiplicities are all at most one, and the displayed equality follows.
No source candidate is compressed to an ungenerated copy.  If a candidate
has a multiplicity-sensitive, role-essential operation or test, its
compatibility with (ii) must be proved separately; it cannot be discarded
by this proof or by the commutant theorem alone.
\end{proof}

\begin{remark}[The remaining universal GP3 comparison obligation]
\label{p2b:rem:gp3-structured-certificate-obligation}
The decisive outstanding calculation is the source-candidate-independent
incidence witness for hypothesis (ii), with the strictness test (iii),
\emph{after} proving candidate-internal exactness (i).  The common generated
six-type support, canonical gradient port and bracket-image objectification
are necessary diagram components, but do not by themselves certify an
arbitrary candidate's other frozen operations or tests.  Paper~II-F,
Proposition~\ref{p2:prop:gp3-multiplicity-bound}, gives the conditional
unmarked-multiplicity reduction when all generated transformations and
separating tests are trivial on that multiplicity factor; the hypothesis
must not be presumed for every admissible candidate.
\end{remark}

\begin{theorem}[Generated-carrier identification with the retained smooth model]
\label{p2b:thm:canonical-retained-smooth-realization}
Paper~I-B, Theorem~\ref{p1:thm:canonical-retained-smooth-model}, defines the
canonical equivariant unitary class as follows.  Let
\(V_R:=\operatorname{Vol}(S^3(\Rmom))\), let
\(x_a\), \(a=1,\ldots,4\), be the round embedding-coordinate module, and put
\[
 e_0:=V_R^{-1/2},\qquad
 u_a:=\frac{2}{\sqrt{3V_R}}\,dx_a,\qquad
 w_{ab}:=\sqrt{\frac{2}{V_R}}\,dx_a\wedge dx_b
 \quad(a<b).
\]
With \(S_+\) and \(S_-\) the two lowest Dirac eigenspaces already belonging
to the generated seed, define
\[
 \mathcal L_R
 :=\operatorname{span}\{e_0\}
 \oplus\operatorname{span}\{u_a\}_{a=1}^4
 \oplus\operatorname{span}\{w_{ab}\}_{a<b}
 \oplus S_+\oplus S_-.
\]
Then the displayed modes are orthonormal, \(\dim_{\mathbb C}\mathcal L_R=15\),
and the explicit one-copy model of
Proposition~\ref{p2b:prop:one-copy-geometric-module} admits the natural unitary
\(\Spin(4)\)-equivariant identification
\[
 J_R:\mathcal H_1\xrightarrow{\ \simeq\ }\mathcal L_R.
\]
It sends \(E_0\) to the constant line, \(E_1\) to the exact first-one-form
module, and \((\mathcal K_{\rm Kill})_{\mathbb C}\) to its Hodge-dual
two-form module.  It is unique up to the already unmarked phases on the six
multiplicity-one irreducible sectors.  Reversing orientation exchanges the
self-dual and anti-self-dual bivector labels but does not change the unitary
class.
\end{theorem}

\begin{proof}
Paper~I-B, Theorem~\ref{p1:thm:canonical-retained-smooth-model}, proves the
orthonormality, representation types, orientation covariance, and
phase-qualified uniqueness.  Theorem~\ref{p2:thm:momentum-host-H1} fixes the six generated
types, and Proposition~\ref{p2b:prop:one-copy-geometric-module}
constructs their one-copy module without a GP3 premise.  The metric
musical map followed by Hodge star identifies the
generated Killing-field module of
Lemma~\ref{p2b:lem:canonical-bracket-image} with Paper~I-A's retained two-form
summand.  Hence the upstream unitary class applies to the complete generated
carrier, with no omitted or inserted type.
\end{proof}

\begin{corollary}[Constituent-level co-realization of the momentum seed]
\label{p2b:cor:momentum-seed-constituent-corealization}
Relative to the G0-occurring transfer-operation root and the frozen
momentum-seed frontier, the complete nonredundant generated constituent family
co-realizes:
\[
 \PhysOcc(\mathscr O_{\rm tr})
 \Longrightarrow
 \operatorname{CoReal}\!\left(
 E_0,S_+,S_-,E_1,(\mathcal K_{\rm Kill})_{\mathbb C}
 \mid \mathscr O_{\rm tr}\right).
\]
Thus the two chiral spinor sectors and the two adjoint irreducible types inside
\((\mathcal K_{\rm Kill})_{\mathbb C}\cong(1,0)\oplus(0,1)\) are not
alternative momentum realizations from which one may be selected.  They are
nonredundant generated constituents of the same occurring momentum composite.
This constituent-level co-realization does not assert occurrence of an
arbitrary state, vector, measurement outcome, or labelled representative in
any constituent.
\end{corollary}

\begin{proof}
Theorem~\ref{p2:thm:momentum-host-H1} identifies exactly the complete generated
carrier family required by GP2 no-omission.  Paper~II-A
Theorem~\ref{p2:thm:generated-constituent-corealization} then applies to these
generated constituents of the already occurring transfer-operation composite.
No additional selector or occurrence root is introduced.
\end{proof}

\begin{corollary}[Common momentum type-support co-realization and conditional endpoint rigidity]
\label{p2b:cor:physical-momentum-carrier-corealization}
If the transfer-operation root occurs in the Paper-II-A G0 sense, the
round parent and the complete generated momentum constituent family
co-realize without a second momentum G0 witness:
\[
 \PhysOcc(\mathscr O_{\rm tr})\Longrightarrow
 \operatorname{CoReal}\!\left(
 S^3_A(\Rmom),\mathcal G^{\rm car}_{\Sigma_p}
 \mid\mathscr O_{\rm tr}\right),\qquad
 \mathcal G^{\rm car}_{\Sigma_p}\cong\mathcal H_1,
 \quad m_p=1.
\]
If, in addition, all hypotheses of
Theorem~\ref{p2b:thm:momentum-carrier-rigidity} are certified, this
co-realization identifies the full carrier of every GP3-irreducible
structured endpoint with \(\mathcal H_p\cong\mathcal H_1\) and its
corresponding exact represented operation envelope.  No measurement,
calibration, normalized state or independent momentum-occurrence claim
is part of either conclusion.
\end{corollary}
\begin{proof}
Paper~II-A's generated-role and generated-constituent co-realization
theorems transport the parent and its frozen generated seed objects from
the occurring transfer root.  Theorem~\ref{p2:thm:momentum-host-H1}
identifies their complete GP2 type-support, and
Corollary~\ref{p2:thm:m1-component} derives its level.  Physical
endpoint multiplicity-one uses the additional whole-structured
comparison hypotheses exactly as stated in
Theorem~\ref{p2b:thm:momentum-carrier-rigidity}; it is not inferred
from the common type-support alone.
\end{proof}

\begin{corollary}[Same-parent status of the retained smooth realization]
\label{p2b:cor:retained-smooth-realization-corealization}
The realization \((\mathcal L_R,J_R)\) is a natural derived presentation of
the common generated one-copy module over the already co-realized round
momentum parent.  It needs no second G0 occurrence witness for the common
generated object content.  Identifying it with the full underlying carrier
of every structured physical endpoint is conditional on the GP3
comparison certificate.  This conclusion actualizes neither an ambient
coordinate lift nor a physical kinetic operator; those are separate target
sorts.
\end{corollary}

\begin{proof}
Every operation used in Theorem~\ref{p2b:thm:canonical-retained-smooth-realization}
is an oriented-round-metric constructor on the already occurring parent:
normalization, exterior differentiation, the metric musical map, Hodge star,
and passage to the lowest spectral subspaces.  Its exact image is the
explicit one-copy generated type-support of
Theorem~\ref{p2b:thm:canonical-retained-smooth-realization}; the remaining
sector phases are unmarked equivalences.  The occurrence of the common
generated constituents follows from
Corollary~\ref{p2b:cor:physical-momentum-carrier-corealization}.  Paper~II-A's generated-role and
derived-constructor witness-transport rules therefore preserve occurrence in
this declared presentation sort.  No ambient coefficient or dynamical owner
appears in the constructor.
\end{proof}

\subparagraph{After-the-fact word classification.}
The already derived one-copy type-support satisfies
\(\mathcal H_1\cong\mathcal C^{\rm word}(1)\), so \(m_p=1\) is its
constructor-length descriptor and not a selector.  The general word-budget
identity and optional tower machinery are in
Paper~II-F, Section~\ref{p2b:app:optional-tower-machinery}.

\begin{theorem}[Intrinsic first-coordinate sector and exact orbit-retention criterion]
\label{p2:thm:intrinsic-coordinate-sector}
On the generated round \(S^3(\Rmom)\), the first nonconstant scalar
Laplace eigenspace is the intrinsically defined four-dimensional
\(\Spin(4)\)-irreducible object
\[
 E_1\cong V_{\frac12,\frac12},\qquad
 \lambda_1=\frac3{\Rmom^2}.
\]
If a finite reciprocal role has a frozen \emph{object-typed} coordinate
port with nonzero image of \(E_1\), and that port is faithful and
closed under the full inherited \(\Spin(4)\)-action and linear
object operations, then it must retain the entire one-copy image of
\(E_1\).  It cannot retain just one axis or a proper nonzero part of
that irreducible object.  A scalar readout without the coordinate-object
port, or a role inheriting only a proper symmetry subgroup, does not by
itself imply this full-sector retention.  This statement determines
irreducible \emph{type support}, not unmarked multiplicity in a larger
realization or an independently prepared coordinate basis.
\end{theorem}
\begin{proof}
For round \(S^3\), the scalar eigenvalues are
\(n(n+2)/\Rmom^2\) and the degree-one eigenspace is the real
coordinate module \(E_1^{\mathbb R}\cong\mathbb R^4\) (or its
complexification).  Its spectral projector is a function of the
intrinsic Laplacian and requires no frame.  The standard
\(SO(4)\)-module \(\mathbb R^4\) is irreducible: the orbit of each
nonzero vector is the full sphere of its radius and its linear span is
all of \(\mathbb R^4\).  Equivalently the complex representation is
\((\frac12,\frac12)\), irreducible under the full
\(SU(2)_L\times SU(2)_R\) action.

A nonzero faithful generated coordinate-object output and the frozen
full-group action therefore generate the whole image by group-orbit
span and the already constituted object-linear operations.  GP2 typed
no-omission retains that image in its coordinate-object sort.  This
argument is stronger than invoking
Lemma~\ref{p2:lem:equivariant-truncation-no-insertion}: that lemma
allows an arbitrary multiplicity projector in a non-multiplicity-free
ambient carrier, whereas the orbit argument establishes the
\emph{specific} inherited E1 image without choosing a copy or
imposing ambient multiplicity-freeness.  Without the declared object
port or full group action the orbit premise is missing, as claimed.
\end{proof}

\begin{proposition}[Frame-independent rank-four orbit certificate]
\label{p2b:prop:v281-e1-orbit-frame-certificate}
On the real first scalar-harmonic space \(E_1^{\mathbb R}\) of the
round \(S^3(\Rmom)\), let \(U(g)\) denote the orthogonal
\(\Spin(4)\) action and let \(0\ne f\in E_1^{\mathbb R}\).
With normalized Haar measure and the round \(L^2\) pairing,
\[
 \mathscr P_f:=\int_{\Spin(4)}
 |U(g)f\rangle\langle U(g)f|\,d\mu(g)
 =\frac{\|f\|_{L^2}^2}{4}\,I_{E_1^{\mathbb R}}.
\]
In particular \(\operatorname{rank}\mathscr P_f=4\), and the group
orbit of any nonzero generated first-coordinate object spans the entire
sector.  The constant \(1/4\) is fixed by dimension and trace, not by
an independently selected coordinate normalization.
\end{proposition}
\begin{proof}
By invariance of Haar measure \(\mathscr P_f\) commutes with the
full orthogonal \(SO(4)\) action.  It is self-adjoint and nonnegative.
The real standard representation is orthogonally irreducible, so every
invariant self-adjoint endomorphism is a real scalar multiple of the
identity.  Taking traces gives
\(\operatorname{tr}\mathscr P_f=\|f\|_{L^2}^2\) and hence the
factor \(1/4\).  Since \(f\neq0\), the operator has strictly
positive eigenvalues and rank four.  The Haar integral is a
mathematical orbit certificate; it is not a new finite physical
averaging or test constructor.
\end{proof}

\begin{remark}[Sharpness of full-symmetry inheritance]
\label{p2b:rem:v281-subgroup-coordinate-counterexample}
If only the subgroup \(SO(3)\subset SO(4)\) fixing one ambient
coordinate direction is constituted, the one-dimensional coordinate
line spanned by \(x\mapsto x_4\) is invariant.  Thus equivariance
under an unspecified proper subgroup, even with a nonzero coordinate
object, cannot force full \(E_1\) retention.  The group and typed-port
hypotheses above are necessary to the claimed orbit implication.
\end{remark}

\subsection{Optional mathematical towers and coordinate interfaces}

\begin{remark}[Integer coordinate-stable tower details]
\label{p2:rem:integer-coordinate-stable-deferred}
The Paper-I-A integer coordinate-stable subfamily is a later optional tower
interface and is not used to derive the momentum seed.  Its relation to the
half-integral component envelope and scalar-seeded integer boundaries is
deferred to Paper~II-F, Section~\ref{p2b:app:integer-coordinate-stable-details}.
\end{remark}

\begin{remark}[Optional lower-ideal classification and sharp countermodels]
\label{p2:rem:legacy-lower-ideal-pointer}
Within the optional \(\mathsf P_{\rm sli}\) grammar, finite supports are
complete total-weight lower ideals.  Under a weaker host signature,
the proper modules \(\mathcal E_5\), \(\mathcal E_7^{L/R}\), and
\(\mathcal E_9\) show sharply that coordinate content, faithfulness, chirality,
and left--right symmetry do not by themselves force \(\mathcal H_1\).  The
revised momentum theorem excludes these models for a different reason:
\((Z_\beta)_{\mathbb C}=(\mathcal K_{\rm Kill})_{\mathbb C}\), where
\(Z_\beta=\operatorname{span}_{\mathbb R}\operatorname{im}\beta\), inherits
the same-role section-object typing and, by the same-role no-splitting rule,
belongs to the momentum carrier; deleting it violates GP2 no-omission.  The full optional
lower-ideal classification and countermodel calculations are deferred to
Paper~II-F, Section~\ref{p2b:app:legacy-lower-ideal}.
\end{remark}

\label{p2:part:minimal-bookkeeping}

The mathematical interfaces below begin from the GP2-derived momentum type-support
\(\mathcal G^{\rm car}_{\Sigma_p}\cong\mathcal H_1\) and its
derived label \(m_p=1\).  It records
finite-state, spectral, and Clifford consequences of that seed and records
only the interfaces required by later-role and recovery analysis.  Detailed
optional word/lower-ideal and component-ceiling machinery is deferred to the
appendices.  The tower is not used to select the momentum carrier, and
\(\mathsf C_{\rm PI}\) is invoked only for the additional normalized
coordinate and protected Clifford data.

\section{Finite representation and operator interfaces of the carrier}
\label{paper2b:sec:organized-carrier}

\subsection{Generated-carrier consequences and downstream interfaces}
\label{p2:sec:uniqueness-m1}

The GP2-required momentum carrier type-support and its intrinsic level have
already been derived, without a GP3 endpoint-multiplicity premise:
\[
 \boxed{\mathcal G^{\rm car}_{\Sigma_p}\cong\mathcal H_1,\qquad m_p=1.}
\]
This section records the downstream consequences and the compact interfaces
consumed later in the series.  Detailed optional tower and ceiling analysis is
deferred to Paper~II-F, Section~\ref{p2b:app:optional-tower-machinery}.

\begingroup
\small
\setlength{\tabcolsep}{2.2pt}
\renewcommand{\arraystretch}{1.14}
\begin{longtable}{|
  L{0.18\linewidth}|
  L{0.37\linewidth}|
  L{0.18\linewidth}|
  L{0.19\linewidth}|}
\caption[Momentum constitution and selector-status interfaces]{Momentum constitution and selector-status interfaces.}
\label{p2:tab:truncation-minimal-bookkeeping}\\
\hline
\textbf{Item} & \textbf{Conclusion} & \textbf{Source} & \textbf{Boundary} \\
\hline
\endfirsthead
\hline
\textbf{Item} & \textbf{Conclusion} & \textbf{Source} & \textbf{Boundary} \\
\hline
\endhead
Momentum type-support / endpoint
& \(\mathcal G^{\rm car}_{\Sigma_p}\cong\mathcal H_1\), \(\dim_{\mathbb C}\mathcal H_1=15\)
& Theorems~\ref{p2:thm:momentum-host-H1}
  and~\ref{p2b:thm:momentum-carrier-rigidity}
& GP2 common type-support unconditionally relative to the II-A frontier; universal GP3 endpoint multiplicity-one requires the comparison certificate. \\
\hline
Derived level
& \(m_p=1\) from \(\ell_{\max}=2\)
& Corollary~\ref{p2:thm:m1-component}
& GP2 type-support descriptor independent of GP3; not a selector. \\
\hline
Optional component ceiling
& \(j_L+j_R\le s_\rho\) for a declared \(P_{\rm ceil}\)
& Theorem~\ref{p2:thm:truncation}
& Restricts allowed labels only; not used for the momentum seed. \\
\hline
Generic word tower
& \(\begin{gathered}
 \mathcal C^{\rm word}(m)\\[-2pt]
 {}=\bigoplus_{j_L+j_R\le m}V_{j_L,j_R}
 \end{gathered}\)
& Theorem~\ref{p2:thm:finite-spectral-exhaustivity}
& Mathematical/\allowbreak later-role grammar; not universal GP2 closure. \\
\hline
Invariant finite states
& commutant classification; normalized trace only under stronger hypotheses
& Theorem~\ref{p2:thm:no-insertion-finite-trace-weight}
& Symmetry alone leaves sector/multiplicity weights free. \\
\hline
Clifford package
& \(\dim_{\mathbb R}A_{\rm ev}=8\), traceless part dimension \(7\)
& Theorem~\ref{p2:thm:minimal-hilbert-structural-15}
& Specializes to \(\mathsf C_{\rm PI}\); does not follow from \(m_p=1\). \\
\hline
Finite supports
& octet/septet/trimer invariant weights
& Corollary~\ref{p2:cor:paper1-octet-descendant-weight-closure}
& Requires a constituted \(\mathsf{Supp}_{2T}\) output: direct core occurrence, or point resolution plus the II-A core-\allowbreak constructor/\allowbreak domain bridge. \\
\hline
\end{longtable}
\endgroup

\paragraph{Logical classification of the mathematical interfaces.}
The entries above are not one undifferentiated selector list.  They fall into
four logically distinct groups:
\begin{enumerate}[label=(\arabic*), leftmargin=2.2em]
\item \textbf{generated/co-realized momentum content:}
\(SU(2)\), the complete generated seed constituents and their
six-type \(\mathcal H_1\) model, and the derived label \(m_p=1\);
universal endpoint multiplicity-one remains proof-gated;
\item \textbf{conditional realization or optional branch data:}
the raw coordinate-completion cone outside an already derived ancestry
branch, the generic invariant-state simplex when a physical preparation is
claimed, and a specified/generated quotient spin or twist member only when a
nontrivial quotient branch is actually asserted;
\item \textbf{G0 occurrences and calibrations:}
actual preparations, point readouts, separately constituted finite supports,
numerical calibration, and other named experimental or dynamical tokens;
\item \textbf{mathematical/presentation infrastructure:}
Hopf representatives, the word tower, lower-ideal and ceiling arithmetic,
integer-boundary conventions, and Casimir normalization.
\end{enumerate}
The second group is branch-conditional rather than part of the active
through-Paper-VI-B structural frontier.  The third group is occurrence or
calibration data, not a structural selector.  Membership in neither group
reopens the derived momentum seed.

\subsubsection{Momentum level as a derived constructor-length label}
\label{p2:subsec:m1-kinematic}
The natural momentum unit \(\Punit=\hbar/\Rmom\) and any independently
declared component ceiling remain useful for later bounded-spectrum questions,
but they do not determine the initial momentum carrier.  The structural
statement is instead the GP2 constructor-relative result
\[
 \mathcal G^{\rm car}_{\Sigma_p}\cong\mathcal H_1,
 \qquad
 m_p=1,
\]
proved above.  Optional ceilings may subsequently be checked for compatibility
with this carrier.

\begin{remark}[Optional level-one tower checks deferred]
The optional consistency statements that recover \(\mathcal H_1\) after an
independent \(\mathsf P_{\rm sli}\) lower-ideal grammar and a level-one
ceiling have been imposed are supplementary checks, not part of the positive
momentum derivation.  They are recorded in
Paper~II-F, Section~\ref{p2b:app:optional-tower-machinery}.
\end{remark}

\begin{corollary}[Structural momentum-role handoff]
\label{p2:cor:neutral-level-one-handoff}
The constructor-relative momentum theorem supplies directly the generated
GP2 type-support and its descriptor
\[
 \boxed{\mathcal G^{\rm car}_{\Sigma_p}\cong\mathcal H_1,\qquad m_p=1.}
\]
Every later result requiring only this generated support and level may use
them without GP3 endpoint multiplicity-one.  The stronger handoff
\(\mathfrak F_p\cong\mathcal H_1\) for every structured physical
endpoint uses the whole-candidate comparison certificate in
Theorem~\ref{p2b:thm:momentum-carrier-rigidity}.  No second momentum witness is required.
This handoff selects neither a point of \(\mathfrak M_{\rm coord}\) nor a
point-reference readout token.  An occurring point-resolution token generates
only the generic finite-resolution support role
\(\mathsf{Supp}_{\varepsilon}(X)\).  The exceptional unmarked
\([2T]\) core-support orbit is co-realized without a second support witness
only after the additional candidate-independent core constructor and the
Paper-II-A exceptional-domain bridge are available, or after a direct
core-constructor occurrence, as stated in
Corollary~\ref{p2:cor:core-support-inheritance-no-duplicate-witness}.
\end{corollary}

\begin{remark}[Finite-family mathematics and physical selection boundary]
\label{p2:rem:paper1-truncation-data}
The geometric and representation atlases of Papers~I-A and I-B classify the mathematical family \(\mathcal H_m\), its dimensions, and
the low-degree representation identities.  Paper~II-A plus the host
closure theorem above supplies the six-type generated content
\(\mathcal H_1\) for the momentum-role signature and Paper~II-A's
generated-role theorem supplies its constituent co-realization.  Universal
endpoint multiplicity-one is subject to the separate GP3 comparison proof.  The raw normalized coordinate intertwiners and the particular protected
Euclidean Clifford vector used for the conditional calculation below belong to
the representative \(\mathsf C_{\rm PI}\); the carrier itself leaves the three
completion moduli unresolved.  This conditional presentation is compatible
with the later completion-independent normalized protected even-algebra result
recorded by Paper~II-A.
\end{remark}

\subsubsection{Invariant finite states and finite-orbit weights}
\label{p2:subsec:finite-carrier-trace-weight-closure}

Two notions must be kept separate.  A state on a matrix algebra is constrained
by the commutant of the represented symmetry; a probability weight invariant
under the constituted action on a finite transitive orbit is fixed by
transitivity.  Neither statement licenses replacing the constituted symmetry
by the full unitary group.

\subparagraph{Optional finite-orbit means.}
Uniform weights on separately actualized finite transitive supports and the
octet/septet/trimer consequences are deferred to
Paper~II-F, Section~\ref{p2b:app:optional-finite-support}.  The invariant matrix-state
classification below remains in the main chain.

\begin{theorem}[Classification of invariant positive normalized functionals]
\label{p2:thm:no-insertion-finite-trace-weight}
Let \(G\) be compact.  On a finite-dimensional complex unitary \(G\)-module,
write the isotypic decomposition as
\[
V=\bigoplus_{\lambda}V_\lambda\otimes M_\lambda,
\]
with pairwise inequivalent irreducible modules \(V_\lambda\).
A functional \(\tau:\End_{\mathbb C}(V)\to\mathbb C\) is positive,
normalized, and invariant under conjugation by \(G\) if and only if
\[
\tau(A)=\Tr_V(\rho A),
\qquad
\rho=\bigoplus_\lambda
I_{V_\lambda}\otimes B_\lambda,
\qquad
B_\lambda\ge0,
\]
with the single normalization condition
\[
\sum_\lambda
(\dim_{\mathbb C}V_\lambda)\Tr_{M_\lambda}(B_\lambda)=1.
\]
Consequently:
\begin{enumerate}[label=(\roman*), leftmargin=2em]
\item on a single irreducible carrier \(V_\lambda\), the invariant state is
the normalized trace
\[
\tau(A)=\frac1{\dim V_\lambda}\Tr(A);
\]
\item on a reducible multiplicity-free carrier, invariant states are
\[
\tau(A)=\sum_\lambda p_\lambda
\frac{\Tr(A_{\lambda\lambda})}{\dim V_\lambda},
\qquad
p_\lambda\ge0,\quad\sum_\lambda p_\lambda=1,
\]
so the sector weights \(p_\lambda\) remain free;
\item on an isotypic component with \(\dim M_\lambda>1\), the positive
operator \(B_\lambda\) may also distinguish multiplicity directions;
\item the normalized trace on all of \(\End(V)\) follows if one additionally
requires traciality, or invariance under the full unitary group \(U(V)\).
\end{enumerate}
Invariant states on a tensor product need not factor: on
\(V_{\frac12}\otimes V_{\frac12}\) under the \emph{diagonal} \(SU(2)\)
action the singlet state is invariant and is not a product state, its
reduced state on either factor having rank two.  (Under a factorwise
\(G_1\times G_2\) action on irreducibles the commutant is scalar and every
invariant state does factor.)  When the normalized
trace is explicitly selected on each factor, their product state has the
familiar multiplicative rank-one weights.
\end{theorem}

\begin{proof}
Finite-dimensional positivity gives a unique density operator
\(\rho\ge0\), \(\Tr\rho=1\), with
\[
\tau(A)=\Tr(\rho A).
\]
This is the finite-dimensional density-matrix representation of a state;
see \cite{BratteliRobinson1987,HornJohnson2013}.
Conjugation invariance is equivalent, by cyclicity and nondegeneracy of the
trace pairing, to
\[
\rho\,U(g)=U(g)\rho
\]
for every \(g\in G\).  The commutant theorem for the displayed isotypic
decomposition gives
\[
\End_G(V)
=
\bigoplus_\lambda
I_{V_\lambda}\otimes\End(M_\lambda),
\]
which proves the asserted form; positivity and normalization give the stated
conditions on \(B_\lambda\)
\cite{fulton1991representation,Hall2015lie}.  The irreducible and
multiplicity-free formulas are its special cases.  If \(\tau\) is tracial, then
\(\Tr(\rho[A,B])=0\) for all \(A,B\), so \(\rho\) is central in
\(\End(V)\) and hence scalar.  Full \(U(V)\)-invariance gives the same
conclusion from its scalar commutant.  Normalization then yields
\(\rho=I_V/\dim V\).  Finally,
\[
\frac{\Tr_{V\otimes W}(A\otimes B)}{\dim V\,\dim W}
=
\frac{\Tr_V(A)}{\dim V}\,
\frac{\Tr_W(B)}{\dim W},
\]
which proves the final conditional statement.
\end{proof}


\begin{proposition}[Converse limitation of normalized unmarked momentum moments]
\label{p2b:prop:normalized-multiplicity-indistinguishability}
For each positive integer \(k\), let
\(\mathcal H^{(k)}=\mathcal H_1\otimes\mathbb C^k\), with an unmarked
replicated momentum operator \(X_k=X\otimes I_k\).  Then for any
nonnegative integer \(n\),
\[
 \frac{\Tr_{\mathcal H^{(k)}}(X_k^n)}{15k}
 =\frac{\Tr_{\mathcal H_1}(X^n)}{15}.
\]
For the dimensionless component
\(X=\hat P_3/(\hbar/\Rmom)=L_3^+-L_3^-\) on the explicit
\(\mathcal H_1\) module, these give
\[
 \frac{\Tr_{\mathcal H_1}(X^2)}{15}=\frac7{15},\qquad
 \frac{\Tr_{\mathcal H_1}(X^4)}{15}=\frac5{12}.
\]
Thus \(m_p=1\) and normalized unmarked spectral moments are not
converse certificates of \(\dim\mathcal H_p=15\): they cannot exclude
uniform replicated implementations.  This does not make those replicas
GP3-admissible minima, nor constitute a measurement that distinguishes them.
\end{proposition}
\begin{proof}
The factorization
\(\Tr_{\mathcal H_1\otimes\mathbb C^k}(X^n\otimes I_k)
=k\Tr_{\mathcal H_1}(X^n)\)
proves the general identity.  For the explicit six types, sum the
weight powers from Proposition~\ref{p2:prop:component_spectrum}:
\(\Tr(X^2)=7\), \(\Tr(X^4)=25/4\).  Dividing by fifteen
gives the stated values.
\end{proof}

\begin{remark}[Support diagnostics are not a multiplicity selector]
\label{p2b:rem:casimir-difference-support-diagnostic}
The mathematical Casimir difference \(D=C_L-C_R\), where \(C_L,C_R\)
have eigenvalues \(j_L(j_L+1),j_R(j_R+1)\), has eigenvalues
\(+2\) and \(-2\) on the two adjoint three-dimensional summands,
respectively.  It can detect missing adjoint \emph{types} under a
separately licensed observable interpretation, but it does not detect
uniform implementation replication from its support alone and it
is not a new physical cutoff or state-selection rule.
\end{remark}

\begin{corollary}[Invariant preparations for the native operation algebras]
\label{p2b:cor:v233-native-invariant-states}
For the mixed scalar-position/enriched-momentum carrier
\(\mathcal P_m\otimes\mathcal H_1\), let the represented algebra be
generated by the factor actions and all native nonzero coordinate
edges. With \(d_m=\dim\mathcal P_m\), its full block algebra is
\(M_{11d_m}\oplus M_{4d_m}\).
The densities invariant under all of its unitary conjugations are
exactly
\[
 \rho_p=\frac{p}{11d_m}P_{\rm geo}
       +\frac{1-p}{4d_m}P_{\rm sp},\qquad 0\le p\le1.
\]
On the matched enriched product with Born exchange, the algebra is
\(M_{121}\oplus M_{88}\oplus M_{16}\), and these densities are
\[
 \rho_{a,b,c}=\frac a{121}P_{gg}
       +\frac b{88}(P_{gs}+P_{sg})+\frac c{16}P_{ss},
 \qquad a,b,c\ge0,\quad a+b+c=1.
\]
The normalized full traces have \(p=11/15\) and
\((a,b,c)=(121,88,16)/225\), respectively; the stated operation
invariances do not select those points.
\end{corollary}
\begin{proof}
The connected coordinate-graph blocks calculated in Paper~I-C,
\ref{p1b:der:g03:s11:sub2}, \ref{p1b:der:g03:s11:sub3} and
\ref{p1b:der:g03:s17}, give the algebras. Explicitly, sandwiching
a nonzero edge between the group matrix units gives every rectangular
map between its endpoints; products give the full matrix algebra
on each connected component. The native component dimensions are
eleven and four. Tensor products give \(121,44,44,16\);
exchange joins the two middle components.
Invariance under all algebra unitaries is membership in its commutant.
Each full block has scalar commutant, so positivity and trace one
give precisely the displayed simplex.
\end{proof}

The free geometric/protected population is distinguishable without selecting
a preferred preparation.  Paper~II-C,
Corollary~\ref{p2c:cor:v2211-tagged-intrinsic-populations}, constructs a
central-sector refinement of the intrinsic localization POVM using the
already generated block projectors.  It preserves the exact coordinate
moments and reads the two populations.  Its untagged predecessor is blind to
that population difference.  This is an available mathematical readout;
an actual tagged measurement still requires its occurrence certificate and
does not select the population before the measurement.

\begin{remark}[A specified loop and a specified ground-state preparation]
Corollary~\ref{p2a:cor:v233-loop-coefficients} fixes a canonical
closed-word readout without preparing any of these invariant states.
For the precise \(K_{\rm LC}\), a zero-energy preparation is uniquely
the scalar projector and the normalized heat family converges to it
at the rate in Corollary~\ref{p1:cor:v233-geometric-ground-state}.
That branch retains the physical operator owner and preparation
condition: the identical principal symbol and coordinate class can
have a different ground space
(Appendix~\ref{p1:app:v233-square-owner}).
\end{remark}


\begin{remark}[Preparation status of invariant-state specialization]
\label{p2b:rem:invariant-state-selector-status}
For the multiplicity-free momentum carrier \(\mathcal H_1\),
\(\Spin(4)\)-invariance alone leaves the simplex
\[
 p_\lambda\ge0,\qquad \sum_\lambda p_\lambda=1
\]
of sector probabilities.  This proves that momentum constitution alone does
not actualize a generic physical preparation.  On a physically constituted zero-energy branch for the precise
\(K_{\rm LC}\) of Corollary~\ref{p1:cor:v233-geometric-ground-state},
and on a branch whose specified ergodic dynamics uniquely fixes the
central weights, the state follows from those preparation conditions.
The symbol alone does not select that branch.  Outside those branches, an actual choice of
\((p_\lambda)\) belongs to the single preparation/measurement occurrence
package of G0.  It is not an additional structural root.

If a later theorem uses the normalized trace on all of
\(\End(\mathcal H_1)\), it must cite the corresponding canonical closed-word, tracial,
dynamical, or occurrence antecedent. The closed-word case is
Corollary~\ref{p2a:cor:v233-loop-coefficients}; it is a readout,
not an assertion that a maximally mixed state was prepared.  Sector-conditioned normalized traces do not by
themselves determine the global preparation weights.
\end{remark}

\begin{remark}[Scope of the generated tests on a finite carrier]
\label{p2:rem:generated-test-scope}
The theorem above classifies invariant positive normalized functionals on the
full algebra \(\End(V)\).  For the block embedding
\(\iota_\lambda(A)\), equal to
\(A\otimes\mathbf1_{M_\lambda}\) on the \(\lambda\)-block and zero on all
other blocks, put
\[
 P_\lambda:=\iota_\lambda(I_{V_\lambda}),
 \qquad
 p_\lambda:=\tau(P_\lambda)
 =(\dim V_\lambda)\Tr_{M_\lambda}(B_\lambda).
\]
Then the exact block formula is
\[
 \tau\!\left(\iota_\lambda(A)\right)
 =p_\lambda\frac{\Tr_{V_\lambda}(A)}{\dim V_\lambda}.
\]
Only after conditioning on a sector of positive weight does the sector weight
cancel:
\[
 \frac{\tau(\iota_\lambda(A))}{\tau(P_\lambda)}
 =\frac{\Tr_{V_\lambda}(A)}{\dim V_\lambda},
 \qquad p_\lambda>0.
\]

The unmarked finite-role doctrine of
Proposition~\ref{p2:prop:gp3-multiplicity-bound} uses the observable algebra
\[
 \mathcal A_{\rm unmarked}
 =\bigoplus_\lambda
 \End(V_\lambda)\otimes\mathbf1_{M_\lambda}.
\]
Restrictions of invariant states to this algebra retain the free sector
weights \(p_\lambda\) but cannot distinguish directions inside
\(M_\lambda\).  A state on the full \(\End(V)\) can distinguish such
directions only through observables acting nontrivially on \(M_\lambda\); the
unmarked hypothesis excludes those observables rather than silently replacing
the sector weight by one.  Finally, \(\tau(\mathbf1)=1\) for every state, so
normalization alone does not return the carrier dimension.  This fixes the
scope of ``generated test'' used in that proposition.
\end{remark}

\begin{remark}[Full basis neutrality is an additional hypothesis]
\label{p2:rem:no-insertion-basis-neutrality}
The absence of a preferred basis does not enlarge a constituted symmetry
group to \(U(V)\).  Full-unitary invariance is a sufficient extra hypothesis
for the normalized trace; it is not a consequence of \(G\)-equivariance on a
reducible carrier.
\end{remark}

\subparagraph{Conditional protected-completion outcome.}
The unconditional GP2 main-text result is the six-type generated support
\(\mathcal G^{\rm car}_{\Sigma_p}\cong\mathcal H_1\), with
\(m_p=1\) and one-copy complex dimension
\(15\).  Under the optional coordinate-completion specialization
\(\mathsf C_{\rm PI}\), the protected real even Clifford algebra has dimension
\(8\), its matrix-traceless part has dimension \(7\), and the resulting
coordinate algebra is noncommutative on the protected block.  These are
conditional completion consequences, not part of the positive momentum-seed
derivation.  Their complete statements, scope guards, and proofs are deferred
to Paper~II-F, Section~\ref{p2b:app:conditional-clifford-package}.
\subparagraph{Deferred optional geometric developments.}
Quotient restrictions, total-weight shell filtration, and Hopf
phase--direction coupling are valid conditional or supplementary structures
but are not used to derive the momentum seed.  Their full statements and
proofs are in Paper~II-F, Section~\ref{p2b:app:optional-finite-support}.

\subsubsection{Momentum-generator convention}

The left- and right-invariant differential generators are retained as auxiliary
spectral tools.  The kinematic momentum triple is
\[
\hat P_i=\frac{\hbar}{\Rmom}(L_i^+-L_i^-),
\]
the transvection complement of the diagonal rotation algebra.  All later
operator and commutator statements use this convention unless a left- or
right-invariant generator is named explicitly.


\section{Recovery maps and continuum comparison}
\label{paper2b:sec:organized-recovery}

\subsection{\texorpdfstring{Continuum Limit and Relation to Classical $S^3$}{Continuum Limit and Relation to Classical S³}}

\subsubsection{\texorpdfstring{The Limit $m \to \infty$}{The Limit m}}

\begin{definition}[Retained Peter--Weyl symbol spaces]
\label{p2:def:retained-pw-symbols}
For each level $m$, define
\[
\mathcal P_m
:=
\operatorname{span}\Bigl\{D^j_{ab}: j\in \tfrac12\ZZ_{\ge0},\; 2j\le m,\; -j\le a,b\le j\Bigr\}
\subset C^\infty(S^3),
\]
where $D^j_{ab}$ are the matrix coefficients of the spin-$j$ irreducible
representation of $SU(2)\cong S^3$.  This is exactly the retained scalar
diagonal coefficient space
\(\mathcal H_m^{\rm sc}\); it is not the full \(\mathcal H_m\),
\(\mathcal A_m\), or \(\End(\mathcal H_m)\).  Its precise finite-operator
relation is the section/retraction proved next.
\end{definition}

\begin{proposition}[Exact scalar Toeplitz section and symbol retraction]
\label{p2:prop:scalar-toeplitz-symbol-retraction}
Use the scalar diagonal identification
\[
 \mathcal H_m^{\rm sc}=\mathcal P_m
 \subset L^2(S^3),
\]
and let \(\Pi_m^{\rm sc}:L^2(S^3)\to\mathcal H_m^{\rm sc}\) be
the orthogonal projector.  Let
\(\iota_m^{\rm sc}:\mathcal H_m^{\rm sc}\hookrightarrow\mathcal H_m\)
be the diagonal-sector inclusion and
\(r_m^{\rm sc}:\mathcal H_m\to\mathcal H_m^{\rm sc}\) its orthogonal
projection.  For \(f\in\mathcal P_m\), define the scalar-block Toeplitz
section
\[
 Q_m^{\rm sc}(f)
 :=
 \iota_m^{\rm sc}\Pi_m^{\rm sc}M_f r_m^{\rm sc}
 \in\End(\mathcal H_m),
\]
where \(M_f\) is scalar multiplication on \(L^2(S^3)\).  For
\(A\in\End(\mathcal H_m)\), define the scalar symbol evaluated at the
constant unit
\[
 \sigma_m^{\rm sc}(A)
 :=
 r_m^{\rm sc}A\iota_m^{\rm sc}(1)
 \in\mathcal P_m.
\]
Here \(1\) is the algebraic constant unit.  It has norm one for
probability-normalized Haar measure, whereas for Riemannian volume its norm
is \(\sqrt{V_R}\), with
\(V_R=\operatorname{Vol}(S^3(\Rmom))\).  The written map uses \(1\) in
either convention.  Using the normalized constant \(e_0=V_R^{-1/2}\)
under Riemannian volume instead would require a compensating factor
\(\sqrt{V_R}\) in the symbol map.
Then both maps are linear and \(\Spin(4)\)-equivariant, and
\[
 \sigma_m^{\rm sc}\!\circ Q_m^{\rm sc}
 =
 \operatorname{id}_{\mathcal P_m}.
\]
Consequently \(Q_m^{\rm sc}\) is injective and
\(\sigma_m^{\rm sc}\) is surjective.  This is an exact bridge between the
retained scalar coefficient space and a scalar-supported linear operator
space in
\(\End(\mathcal H_m)\).  It is not a multiplicative embedding, it does not
identify \(\mathcal P_m\) with \(\mathcal A_m\) or
\(\End(\mathcal H_m)\), and it supplies no scalar symbols for the
non-diagonal bundle sectors.
\end{proposition}

\begin{proof}
The invariant constant function \(1\) belongs to \(\mathcal P_m\), and
\(\Pi_m^{\rm sc}\) fixes every element of \(\mathcal P_m\).  Therefore, for
each \(f\in\mathcal P_m\),
\[
\begin{aligned}
 \sigma_m^{\rm sc}\!\left(Q_m^{\rm sc}(f)\right)
 &=
 r_m^{\rm sc}\iota_m^{\rm sc}
 \Pi_m^{\rm sc}M_f
 r_m^{\rm sc}\iota_m^{\rm sc}(1)\\
 &=
 \Pi_m^{\rm sc}(f\cdot1)
 =f.
\end{aligned}
\]
This proves the retraction identity, injectivity, and surjectivity.  The
left--right action preserves the scalar diagonal sector and its orthogonal
complement, fixes \(1\), commutes with \(\Pi_m^{\rm sc}\), and satisfies
\(U_gM_fU_g^{-1}=M_{U_gf}\).  Hence both displayed maps intertwine the
\(\Spin(4)\)-actions.  The remaining boundary statements follow directly
from the types of the maps: no injectivity of the retraction and no
product-preservation property has been imposed.
\end{proof}

\begin{theorem}[Peter--Weyl parent-symbol density from the finite tower]
\label{p2:thm:continuum_limit}
Let $\mathcal P_m$ be the retained Peter--Weyl coefficient space of Definition~\ref{p2:def:retained-pw-symbols}. Then
\[
\mathcal P_{\mathrm{poly}}:=\bigcup_{m\ge0}\mathcal P_m
\]
is the Peter--Weyl polynomial algebra on $SU(2)$ and is uniformly dense in
$C(S^3)$.  Equivalently, for every \(f\in C(S^3)\),
\[
 \lim_{m\to\infty}
 \inf_{p\in\mathcal P_m}\|f-p\|_\infty=0.
\]
The exact maps in
Proposition~\ref{p2:prop:scalar-toeplitz-symbol-retraction} type the
finite-level scalar coefficient bridge.  The limit is a statement about the
parent presentation and the directed family of finite symbols; no realized
stage has \(m=\infty\).  The conclusion is coefficient-space density only.  It
does not identify finite matrix algebras as literal subalgebras of
\(C(S^3)\), and it is neither an operator-norm nor a dynamical convergence
theorem.
\end{theorem}

\begin{proof}
For every spin \(j\in\frac12\ZZ_{\ge0}\), choosing any integer \(m\ge 2j\)
places all matrix coefficients \(D^j_{ab}\) in \(\mathcal P_m\).  Hence
\[
\bigcup_m\mathcal P_m
=
\operatorname{span}\{D^j_{ab}:j\in\tfrac12\ZZ_{\ge0},\,-j\le a,b\le j\},
\]
which is precisely the Peter--Weyl polynomial algebra of the compact group
\(SU(2)\).  The Peter--Weyl uniform-density theorem
\cite{Peter1927,helgason1984groups} states that this algebra is dense in
\(C(SU(2))\) in the uniform norm.  After the identification
\(SU(2)\cong S^3\), this gives
\[
\overline{\mathcal P_{\mathrm{poly}}}^{\|\cdot\|_\infty}=C(S^3).
\]
Because the spaces are nested, this closure statement is equivalent to the
displayed limit of best uniform-approximation errors.  The section and
retraction of
Proposition~\ref{p2:prop:scalar-toeplitz-symbol-retraction} do not change the
type of this result: it is not an operator-norm convergence theorem for
\(\mathcal A_m\) or \(\End(\mathcal H_m)\).
\end{proof}

\begin{remark}[Mode of convergence]
\label{p2:rem:continuum-mode}
The theorem establishes density of the retained Peter--Weyl symbols in the
classical commutative algebra $C(S^3)$.  This parent-symbol density statement concerns the scalar Peter--Weyl symbol content and is typed as a presentation recovery, not an ontic continuum.  The additional non-diagonal
$(j_L,j_R)$ sectors of the full finite module package require their separate
vector/spinor/bundle interpretation and should not be identified with scalar
functions on $S^3$.  Their share is the larger one: since
$\dim\mathcal P_m=(m+1)(m+2)(2m+3)/6$ and
$\dim\mathcal H_m=\binom{2m+4}{4}$,
\[
 \frac{\dim\mathcal P_m}{\dim\mathcal H_m}=\frac{1}{2m+1},
\]
so the density statement addresses a fraction of the carrier that falls to
zero as the level grows.  More refined semiclassical convergence statements require
additional product, norm, and asymptotic estimates; the exact finite-level
section/retraction above does not supply them by itself.
\end{remark}


\section{Finite products and exact retained support}
\label{paper2b:sec:organized-finite}

\subsection{Quantitative product-word retention at the frozen seed}
\label{p2b:sec:v2210-lit-product-retention}
The generated $\mathcal H_1$ remains the frozen linear/bracket momentum
carrier.  The following result extends the multiplication-defect identity
already proved in Paper~I-C to a quantitative test for a consumed finite
word.  It does not license a pointwise multiplication constructor or a new
physical momentum level.

\begin{theorem}[Exact ordered-word defect and its bound]
\label{p2b:thm:v2210-lit-word-defect}
Let $P=P^\dagger=P^2$, $Q=I-P$, and let $A_1,\ldots,A_n$ be bounded
operators on a supplied common parent, $n\ge2$.  On $P\mathcal H$, put
\[
 W_P=(PA_1P)\cdots(PA_nP),\qquad W=P A_1\cdots A_n P.
\]
Then exactly
\[
 W-W_P=\sum_{j=1}^{n-1}
 (PA_1P)\cdots(PA_{j-1}P)\,PA_jQ A_{j+1}\cdots A_nP,
\]
where the prefix is the identity on $P\mathcal H$ when $j=1$.  Consequently
\[
 \|W-W_P\|\le\sum_{j=1}^{n-1}
 \left(\prod_{i<j}\|A_i\|\right)
 \|PA_jQ\|\,\|QA_{j+1}\cdots A_nP\|.
\]
For a supplied isometry $Z$ with $P=ZZ^\dagger$, define
$\pi_Z(A)=Z^\dagger AZ$.  Its positive square defect and mixed defect obey
\[
 D_Z(A):=\pi_Z(A^\dagger A)-\pi_Z(A)^\dagger\pi_Z(A)
 =(QAZ)^\dagger(QAZ),
 \qquad \|D_Z(A)\|=\|QAZ\|^2,
\]
\[
 \|\pi_Z(AB)-\pi_Z(A)\pi_Z(B)\|
 \le\|QA^\dagger Z\|\,\|QBZ\|.
\]
\end{theorem}
\begin{proof}
Insert $I=P+Q$ at the first internal interface of $W$, and repeat only on
the term whose retained prefix still ends in $P$.  This partitions the
difference by the first omitted intermediate image, giving the exact sum.
Submultiplicativity gives the bound.  The square and mixed identities
follow by inserting $ZZ^\dagger=P$ between the compressed factors, as in
Paper~I-C, Lemma~\ref{p1b:lem:v2167-compression-multiplication}.
For the square, $T^\dagger T$ has norm $\|T\|^2$; for the mixed defect
use $Z^\dagger AQBZ=(QA^\dagger Z)^\dagger(QBZ)$.
\end{proof}

\begin{corollary}[A consumed comparison needs only its visible defect]
\label{p2b:cor:v2210-lit-visible-word-defect}
For normalized source vectors $u,v\in P\mathcal H$ or a density operator
$\sigma$ on that range,
\[
 |u^\dagger(W-W_P)v|\le\|W-W_P\|,
 \qquad |\operatorname{Tr}\sigma(W-W_P)|\le\|W-W_P\|.
\]
Exact reduction of the consumed parent algebra makes all defects zero.
An observable-specific zero contraction can suffice for that comparison
without proving reduction of every mathematically available parent word.
If source derivatives are consumed, the same test applies separately to
the actual differentiated word, including derivatives of $P$ or $Z$ when
the retained chart moves.
\end{corollary}
\begin{proof}
The first inequality is the operator norm estimate, and the second follows
from positivity and unit trace.  Reduction annihilates $PAQ$ and $QAP$.
The final assertion is application of the same identities to the fully
differentiated supplied word, not deletion of chart-derivative terms.
\end{proof}
The $SU(2)$ Fourier/Duflo construction
\cite{LitGEOFreidelMajid2008Fourier} provides a useful explicit comparison
between group composition, star multiplication and ordinary symbols, but
its $2+1$-dimensional Euclidean model and ambient function algebra are not
an identification of this finite carrier.  Any such comparison uses its
licensed map and the defect above.  Full pointwise scalar products still
have the additional harmonic outputs proved in
Section~\ref{p2b:sec:v2201-derived-d01}; no Fourier name removes those
outputs or changes $m_p=1$.


\section{Physical scope and downstream handoff}
\label{paper2b:sec:organized-scope-exports}

\subsection{Geometry-first consequences and physical scope}
\label{p2b:sec:v2201-geometry-first-consequences}
\subparagraph{Geometry ancestry.}
Round Hopf harmonics $\to$ frozen momentum seeds and linear/bracket carrier $\to$ explicitly licensed product outputs.
\subparagraph{Hypotheses and occurrence boundary.}
These are consequences on the domains stated below.
A mathematical certificate or an executed restriction retains every
physical source, state, clock and readout hypothesis of its inputs.

\subsubsection{First multiplicative extension of the frozen carrier}
\label{p2b:sec:v2201-derived-d01}

\begin{proposition}[Degree-two output under licensed full pointwise multiplication]
\label{p2b:prop:derived-d01-pointwise-extension}
The scalar constant and degree-one harmonics occur in the retained carrier.
If a downstream constructor licenses their \emph{full pointwise product},
it must retain the nine-dimensional degree-two scalar harmonic output.
For ambient scalar coordinates on $S^3(R)$,
\[
 x_ax_b=\frac{R^2}{4}\delta_{ab}
 +\left(x_ax_b-\frac{R^2}{4}\delta_{ab}\right),\qquad
 \dim E_2=9.
\]
The second term is the symmetric traceless quadratic sector. More generally,
products of at most $r$ degree-one functions span
\[
 \bigoplus_{\ell=0}^r E_\ell,
 \qquad \dim=\sum_{\ell=0}^r(\ell+1)^2
 =\frac{(r+1)(r+2)(2r+3)}6.
\]
Here $E_\ell$ is the degree-$\ell$ scalar harmonic eigenspace.
It is not the full auxiliary module $\mathcal H_\ell$; in particular
$\dim E_2=9$ whereas $\dim\mathcal H_2=\binom84=70$.
\end{proposition}
\begin{proof}
Restrict homogeneous polynomials to the sphere and remove their
Euclidean traces. Each polynomial decomposes into harmonic polynomials times
powers of $|x|^2$, which become constants on the sphere. Conversely the
coordinate monomials generate every such polynomial. The symmetric
quadratic tensors have dimension ten, with one trace.
\end{proof}
Thus linear/bracket closure of the 15-dimensional one-copy module is not closure
under all pointwise products. Projected operator composition and full
function multiplication are different words; a finite matrix algebra can
close under the former without acquiring the latter's harmonic outputs.
This is a capacity result, not a license to enlarge the physical carrier. The exact omitted-image error for each actually licensed retained word is now given by Theorem~\ref{p2b:thm:v2210-lit-word-defect}; its source-visible contraction remains the application calculation.
\par\noindent\textit{Inherited inputs:} \ref{p2:thm:momentum-host-H1}, \ref{p2b:thm:momentum-carrier-rigidity}, \ref{p3f:thm:v180-incidence-harmonic-resolution}.

\begin{theorem}[No finite covariant multiplicative sphere realization]
\label{p2b:thm:finite-covariant-classical-algebra-obstruction}
Let a group $G$ act transitively by homeomorphisms on an infinite compact Hausdorff space $X$,
and let $U$ be a unitary representation on a nonzero finite-dimensional
complex Hilbert space $H$.  There is no unital $G$-equivariant
$*$-homomorphism
\[
 \pi:C(X)\longrightarrow\operatorname{End}(H),\qquad
 \pi(f\circ g^{-1})=U_g\pi(f)U_g^*.
\]
In particular, the round $\Spin(4)$ action on $S^3$ admits no such finite
realization.  Every finite unital covariant positive linear map from its
full continuous function algebra must therefore fail multiplicativity on
some pair of functions.
\end{theorem}

\begin{proof}
The finite-dimensional commutative $C^*$-algebra $\pi(C(X))$ is simultaneously
diagonalizable.  Its distinct characters are evaluations at finitely many
points $x_1,\ldots,x_N\in X$, so
\[
 \pi(f)=\sum_{j=1}^N f(x_j)P_j,
 \qquad P_j\ne0,\quad \sum_jP_j=I.
\]
The kernel consists precisely of functions vanishing on the nonempty finite
support $F=\{x_1,\ldots,x_N\}$.  Equivariance preserves this kernel and hence
preserves $F$.  Transitivity forces every nonempty invariant subset to equal
$X$, contradicting its infinitude.  A unital positive map on $C(X)$ respects
complex conjugation; if it were multiplicative everywhere, it would be the
forbidden $*$-homomorphism.
\end{proof}

\begin{remark}[Structural defect and the frozen carrier]
\label{p2b:rem:finite-covariant-defect-scope}
The obstruction explains why the product outputs in D01 and the word defect
in Theorem~\ref{p2b:thm:v2210-lit-word-defect} cannot all disappear in a finite
round-covariant function realization.  It does not prohibit exactness on a
specified finite test space, finite sampling without the full continuous
group action, or closure under matrix multiplication.  Nor does it select a
particular POVM map or enlarge the frozen momentum carrier.  Paper II-C's
optimal coordinate POVMs realize the positive-map alternative; their
remaining family freedom must be retained by a downstream prediction that
consumes classical quadratic readouts.
\end{remark}


\subsection{Finite coordinate-word closure on the frozen momentum seed}
\label{p2b:sec:v226-coordinate-words}
\begin{corollary}[Degree six is an operation bound, not a momentum level]
\label{p2b:cor:v226-coordinate-word-degree}
On the generated intrinsic class define
$\mathcal W_d=\operatorname{span}_{\mathbb C}\{I,Y_{a_1}\cdots Y_{a_s}:s\le d\}$.
Then $\dim\mathcal W_5=136$ and
$\mathcal W_6=\mathcal A_D$ has dimension 137. The intrinsic coordinate-word calculation here acts on the explicit
one-copy module $\mathcal H_1$ of dimension fifteen; universal
identification with every GP3 physical endpoint retains the additional
comparison hypothesis of Theorem~\ref{p2b:thm:momentum-carrier-rigidity}.
The label $d$ here is not the carrier tensor-word level $m$.
\end{corollary}
\begin{proof}
The degree-six identity and explicit degree-five obstruction are
Theorem~\ref{p1:thm:v226-chirality}; the exact rank certificate gives the
lower dimension 136. These are products of represented endomorphisms on
one carrier, not an insertion of a new representation sector. The frozen
seed frontier and its no-insertion theorem therefore retain their domain.
\end{proof}
The 137-outcome native measurement is a finite effect construction on that
same algebra (Theorem~\ref{p1:thm:v226-tomography}). Its cardinality is not
a site count, particle species count or new occupation ceiling. Exact
moment-preserving localization still has its separate radius and finite
support constraints, including the universal chirality-information bound
of Theorem~\ref{p2c:thm:v226-chirality-radius}.


\section{Finite-algebra capacity and native multiplicity limits}
\label{p2b:sec:v352-capacity-topology}
Here the finite algebra is first taken as an \emph{already constituted
abstract test algebra}.  Reconstructing its minimum module must not
be reversed into an independent proof that GP2 necessarily generates
that abstract algebra from bare G0.

\begin{theorem}[Faithful module and commutant characterization]
\label{p2b:thm:v352-minimal-faithful}
For $\mathcal A=\bigoplus_{i=1}^{r}M_{d_i}(\mathbb C)$, every
finite-dimensional faithful unital $*$-representation is of the form
\[
 \pi(\mathcal A)=\bigoplus_{i=1}^{r}
     \bigl(M_{d_i}(\mathbb C)\otimes I_{m_i}\bigr),\quad m_i\geq1,
 \qquad \dim_{\mathbb C}\mathcal H=\sum_i d_i m_i.
\]
The minimum is $\sum_i d_i$, achieved precisely when every $m_i=1$.
The minimum faithful module is unique up to a unitary intertwiner
preserving the labeled central blocks.  Moreover
\[
 \pi(\mathcal A)'=\bigoplus_i
  I_{d_i}\otimes M_{m_i}(\mathbb C),\qquad
 \pi(\mathcal A)'=Z(\pi(\mathcal A))
 \quad\Longleftrightarrow\quad m_i=1\ \text{for all }i.
\]
\end{theorem}
\begin{proof}
A full complex matrix algebra has, up to equivalence, one irreducible
unital representation, its defining module.  Restrict a faithful
representation to every central simple summand to obtain the displayed
multiplicities, which must all be positive.  Compute the commutant
block by block using matrix units.  The inequality and equality
criterion follow immediately.
\end{proof}

\paragraph{Constructive matrix-unit verification of Theorem~\ref{p2b:thm:v352-minimal-faithful}.}
For each central summand choose matrix units $e_{ab}^{(i)}$ and put
$P_a^{(i)}=\pi(e_{aa}^{(i)})$.  These are pairwise orthogonal
projections.  The $\pi(e_{ab}^{(i)})$ identify their ranges unitarily,
so each block has a common positive rank $m_i$ (faithfulness makes
$m_i\geq1$).  Choosing an orthonormal basis in the range of
$P_1^{(i)}$ and transporting it by $\pi(e_{a1}^{(i)})$ constructs
$\mathcal H_i\cong\mathbb C^{d_i}\otimes\mathbb C^{m_i}$ explicitly.
The matrices commuting with all these units are precisely
$I_{d_i}\otimes M_{m_i}(\mathbb C)$, and central projections forbid
inter-block commutant terms.  Hence $\dim\mathcal H=\sum_i d_i m_i$
and $\pi(\mathcal A)'=Z(\pi(\mathcal A))$ exactly when all $m_i=1$.
This proof does not invoke the classification of irreducible representations
as an external starting result.  It nevertheless presupposes the same
already constituted block algebra; it provides no new GP2 generation,
physical preparation, or universal GP3 endpoint certificate.

\begin{theorem}[Single-shot distinguishability and pure-state Euler invariant]
\label{p2b:thm:v352-capacity-euler}
Let $C_1(\mathcal A)$ be the maximum number of algebraic states
perfectly discriminated by \emph{one} POVM with effects in
$\mathcal A$, and let $\mathcal P(\mathcal A)$ denote the abstract
pure-state space with its weak-$*$ topology.  Then
\[
 \boxed{C_1(\mathcal A)=\sum_i d_i
 =\chi(\mathcal P(\mathcal A))
 =\dim_{\mathbb C}\mathcal H_{\min}},\qquad
 \mathcal P(\mathcal A)\simeq\bigsqcup_i\mathbb{CP}^{d_i-1}.
\]
Both quantities are independent of all faithful-representation
multiplicities $m_i$ as long as the \emph{same abstract effects}
constitute the test algebra.
\end{theorem}
\begin{proof}
A rank-one orthogonal basis of each summand, assembled across all
summands, gives $\sum_i d_i$ states and a single perfectly
discriminating POVM.  Conversely, for such a POVM
$\{e_j\}$ and perfectly discriminated density functionals
$\rho_j$, positivity implies that $e_j$ equals the identity on
the support of $\rho_j$ and vanishes on every other support.
The nonzero supports are pairwise orthogonal; their total rank is at
most $\sum_i d_i$.  A pure state of a finite direct sum belongs
to exactly one simple summand, whose pure states are rank-one
projectors $\mathbb{CP}^{d_i-1}$.  Finally
$\chi(\mathbb{CP}^{n})=n+1$ and the Euler characteristic is
additive over finite disjoint unions.  Apply the previous theorem.
\end{proof}

\begin{corollary}[Invariant fifteen-state capacity across the three algebra stages]
\label{p2b:cor:v352-three-stage-capacity}
For the already generated abstract stages
\[
 \mathcal A_{\rm kin}=\mathbb C\oplus M_2\oplus M_2\oplus
 M_4\oplus M_3\oplus M_3,\quad
 \mathcal A_D=M_{11}\oplus M_4,\quad
 \mathcal A_{\rm full}=M_{15},
\]
\[
 \boxed{C_1(\mathcal A_{\rm kin})=C_1(\mathcal A_D)
  =C_1(\mathcal A_{\rm full})=15},\qquad
 \chi(\mathcal P(\mathcal A_j))=15.
\]
The algebra dimensions are $43,137,225$, the minimal-module
commutant dimensions are $6,2,1$, and the numbers of connected
pure-state components are $6,2,1$, respectively.  The maximum
single-shot classical alphabet carries $\log_2(15)$ bits of
\emph{distinguishability capacity}; this is not a physical bit rate.
The third algebra and its state/effect role remain conditional on
licensing the full field/event test domain.
\end{corollary}
\begin{proof}
Apply the preceding two theorems to the displayed block sizes.
The component count is the number of central simple summands, not
the Euler characteristic.
\end{proof}

\begin{proposition}[Normalized-copy no-go for universal multiplicity one]
\label{p2b:prop:v352-normalized-amplification}
For any $k\geq1$ and a finite $*$-polynomial $p$ in canonical
algebra elements $a_1,\ldots,a_n$,
\[
 \boxed{\frac{\operatorname{Tr}_{15k}
   p(a_1\otimes I_k,\ldots,a_n\otimes I_k)}{15k}
 =\frac{\operatorname{Tr}_{15}p(a_1,\ldots,a_n)}{15}.}
\]
The same algebraic effect statistics of a state are obtained by
partial trace over the unmarked multiplicity factor.  Thus normalized
trace words and these multiplicity-blind protocols cannot force
$k=1$.  A full GP3 whole-candidate comparison, independently
licensed multiplicity-resolving role, or absolute capacity witness
for an \emph{enlarged} algebra is still required to assert a
universally one-copy physical Hilbert carrier.
\end{proposition}
\begin{proof}
Polynomial evaluation commutes with tensoring by $I_k$ and
$\operatorname{Tr}(T\otimes I_k)=k\operatorname{Tr}(T)$.
For a density $\rho$ and effect $a\otimes I_k$, the partial-trace
identity is $\operatorname{Tr}[\rho(a\otimes I_k)]
=\operatorname{Tr}[(\operatorname{Tr}_k\rho)a]$.
\end{proof}

\begin{remark}[Scope of the capacity closure]
\label{p2b:rem:v352-capacity-occurrence}
The algebraic maximum presumes the role allows every mathematical
state and effect used in its definition; an actually performed
preparation, instrument and outcome require their original G0/GP2
certificates.  The new statement fixes the \emph{15-state capacity
of the constituted observable algebra} without claiming that all
physical carriers or all experimental instruments have Hilbert
space dimension fifteen.  The constitutional axiom GP3 is unchanged.
\end{remark}


\subsection{Process-level native multiplicity descent with finite ancillas}
\label{p2b:sec:v355-native-multitime}
The following theorems strengthen the preceding normalized-trace
obstruction. Here \emph{native} means that every readout and instrument
branch is a completely positive map of the specified abstract constituted
test algebra; no dynamics in the otherwise invisible multiplicity commutant
is admitted. In particular, the universal GP3 same-role comparison is \emph{not}
proved by the process-equivalence statements.


\begin{theorem}[Finite native-instrument histories descend under amplification]
\label{p2b:thm:v353-multitime-amplification}
Let $\mathcal A$ be a finite-dimensional unital $C^*$-algebra and
$\pi_k(a)=\pi(a)\otimes I_k$ an unmarked common amplification of a
faithful finite-dimensional unital $*$-representation
$\pi:\mathcal A\to B(\mathcal H)$, with
$\pi(1_{\mathcal A})=I_{\mathcal H}$. For any finite sequence of completely
positive Heisenberg instrument branches
$\mathcal J_{s_j,x_j}:\mathcal A\to\mathcal A$ with
$\sum_x\mathcal J_{s,x}(1)=1$, the joint probability of a history
in an initial density matrix $\rho$ is
\[
p(x_1,\ldots,x_n\mid s_1,\ldots,s_n)
=\omega_\rho\!\left(
\mathcal J_{s_1,x_1}\circ\cdots\circ
\mathcal J_{s_n,x_n}(1)\right),\quad
\omega_\rho(a)=\operatorname{Tr}[\rho\,\pi_k(a)].
\]
Every such history distribution agrees with that of a state on the
minimum faithful module of $\mathcal A$, independently of $k$;
the claim extends to classical-adaptive choices of settings.
This is a theorem about the abstract native instrument algebra, not about
all possible amplified-system apparatus interactions.
\end{theorem}
\begin{proof}
The specified composition takes values in $\mathcal A$, so all history
probabilities depend only on the restricted positive normalized functional
$\omega_\rho$. Every finite-dimensional $C^*$-algebra state has a density
operator in its minimum faithful representation, proving equivalence for
all histories. Adaptivity only chooses a branch-dependent CP word.
\end{proof}


\begin{theorem}[Ancilla-stable native equivalence of all faithful multiplicities]
\label{p2b:thm:v354-ancilla-stable-descent}
Let $\mathcal A=\bigoplus_{\lambda}M_{d_\lambda}(\mathbb C)$.
Every finite-dimensional faithful unital $*$-representation has the form
$\pi_{\mathbf m}(a)=\bigoplus_\lambda(a_\lambda\otimes I_{m_\lambda})$
with $m_\lambda\ge1$. For any finite reference Hilbert space $\mathcal K$
and any density on the amplified representation tensored with
$\mathcal K$, sectorwise partial trace of the unmarked multiplicity
factors yields a canonical density reproducing all tests of
$\mathcal A\otimes B(\mathcal K)$. This equivalence persists through
all finite adaptive histories of CP instruments defined on that
abstract joint algebra. Therefore no such native protocol, even with
entangled finite ancillas, can infer $\mathbf m$.
\end{theorem}
\begin{proof}
Decompose the state by central sectors and trace out each
$\mathbb C^{m_\lambda}$. The resulting canonical block density gives
identical expectations for every abstract joint observable.
Each $\pi_{\mathbf m}\otimes\operatorname{id}_{M_n}$ is a
complete-order isomorphism onto its image. Composing CP Heisenberg
instrument maps in the abstract algebra preserves the equality
on each outcome-history branch. No operation outside the native algebra
is included.
\end{proof}


\begin{remark}[GP3 and physical-operation boundary]
\label{p2b:rem:v355-native-operation-firewall}
The canonical minimum faithful representation determines the complete
operator-algebraic native test capacity, but neither the actual microscopic
Hilbert-space multiplicities nor whether a multiplicity-sensitive operation
occurs. These theorems do not remove GP3 as a primitive principle or replace
its universal whole-candidate comparison certificate. Finite entangled
ancillas change none of the abstract-native statistics, even though an
independently licensed commutant-sensitive channel would change the test
algebra and invalidate that no-go's premises.
\end{remark}

\section{Inverse diagnostics and conditional action interfaces}
\label{p2b:sec:v264-reverse-certificates}
\noindent
This section uses mathematical inputs from Papers~I-A, I-B, II-C, II-D
and III-F to derive converses and obstructions on their stated domains.
None of these inputs is used to select the momentum
parent, the frozen constructor frontier, or the level $m_p=1$.
A canonical mathematical operator is not, without its own stage-safe
consumer, an occurring physical event, action, process, or test.
Detailed supporting calculations and counterexamples belong to the
technical verification record, not to the assumptions of these results.

\begin{theorem}[Reverse rigidity of the native $15/137$ data]
\label{p2b:thm:v264-algebra-rigidity}
If $\mathcal A\subseteq M_{15}(\CC)$ is a \emph{unital complex
finite-dimensional $C^*$-subalgebra} with
$\dim_{\CC}\mathcal A=137$, then its represented form is, up to a
unitary change of basis,
\[
 \boxed{\mathcal A=M_{11}(\CC)\oplus M_4(\CC)}
\]
with each simple block represented once.  Neither the original
coordinate-graph constructor nor its occurring action/readout ownership
follows from these two numbers alone.
\end{theorem}
\begin{proof}
Write the faithful unital representation as
$\mathcal A\cong\bigoplus_i M_{n_i}(\CC)$ on
$\bigoplus_i(\CC^{n_i}\otimes\CC^{m_i})$, where $m_i\ge1$.
Then $\sum_i n_i^2=137$ and $\sum_i n_im_i=15$.
If every $n_i\le9$, then
$137=\sum_i n_i^2\le9\sum_i n_i\le135$, impossible.
Also $n_i\le11$ because $12^2>137$. If a largest block has
$n_i=10$, the remaining blocks have total defining dimension at most
five and therefore total squared dimension at most $5^2=25$, whereas
$137-100=37$ is required. Hence a largest block is $11$.
The remaining blocks have total defining dimension at most four and
squared dimension $137-121=16$, which forces precisely one block of
size four. These two blocks exhaust the $15$ representation dimensions,
so both multiplicities are one.
\end{proof}

\begin{proposition}[Inverse quadratic fingerprint within the coordinate atlas]
\label{p2b:prop:v264-quadratic-inversion}
Fix the normalized Hermitian equivariant coordinate-completion atlas
of Paper~I-A with $R_{\rm mom}>0$ and nonnegative reduced magnitudes
$(r_{1,0},r_{0,1},\varrho_{\rm sp})$. Denote the scalar, vector, left
bivector, right bivector, and spinor eigenvalues of
$Q=\sum_{a=1}^{4}X_a^2$ by $(q_0,q_V,q_L,q_R,q_S)$.
Within this declared atlas, their exact inverse is
\[
\begin{gathered}
 \boxed{\begin{gathered}
  R_{\rm mom}=\sqrt{q_0},\qquad r_{1,0}=\sqrt{q_L/2},\\
  r_{0,1}=\sqrt{q_R/2},\qquad \varrho_{\rm sp}=\tfrac12\sqrt{q_S}
 \end{gathered}},\\[5pt]
 \boxed{q_V=\tfrac14q_0+\tfrac34(q_L+q_R)}.
\end{gathered}
\]
Thus nonnegative sector data satisfying the last identity reconstruct
all reduced magnitudes modulo the atlas's phase equivalence. Outside
the stipulated atlas, neither the same fingerprint nor the same abstract
algebra reconstructs the coordinate operator or its action provenance.
\end{proposition}
\begin{proof}
The fixed scalar/vector edge and outgoing-path trace identities
of Paper~I-A give
$q_0=R_{\rm mom}^2$, $q_L=2r_{1,0}^2$,
$q_R=2r_{0,1}^2$, $q_S=4\varrho_{\rm sp}^2$ and
$q_V=R_{\rm mom}^2/4+\tfrac32(r_{1,0}^2+r_{0,1}^2)$.
Solving gives the claimed necessary and sufficient algebraic relations
inside that atlas. Existence and phase classification are inherited
from the specified atlas, not inferred for arbitrary operator tuples.
\end{proof}

\begin{corollary}[Explicit central grading as a coordinate word]
\label{p2b:cor:v264-central-polynomial}
On the \emph{already generated} intrinsic fixed-radius class of
Paper~II-C, put $q=Q_D/R_{\rm mom}^2$. Then
\[
 q=P_0+\frac{11}{12}P_V+\frac49P_B+\frac34P_S,
\quad
\boxed{P_S=\frac{864}{11}(q-I)
 (q-\tfrac{11}{12}I)(q-\tfrac49I),\quad
 \Gamma_c=I-2P_S.}
\]
Consequently $P_S$ and $\Gamma_c$ are polynomials of coordinate-word
degree at most six; no new operator \emph{choice} is needed on this branch.
A performed measurement of $P_S$ still needs an actual readout/test token.
\end{corollary}
\begin{proof}
The fingerprint is Paper~II-C's intrinsic coordinate theorem.
Its four distinct eigenvalues permit Lagrange interpolation at
$3/4$, with denominator
$(3/4-1)(3/4-11/12)(3/4-4/9)=11/864$.
The grading acts as $+1$ on the rank-$11$ complement and $-1$ on
its rank-$4$ spinor block. Each $Q_D$ uses two coordinate factors.
\end{proof}

\begin{theorem}[Full spinorial representation content of the odd bridge]
\label{p2b:thm:v264-odd-bridge-types}
For the derived carrier, write
\[
\begin{split}
\mathcal H_g&=(0,0)\oplus(\tfrac12,\tfrac12)
  \oplus(1,0)\oplus(0,1),\\
\mathcal H_s&=(\tfrac12,0)\oplus(0,\tfrac12).
\end{split}
\]
Then as $\Spin(4)$ modules,
\[
\boxed{\begin{aligned}
\operatorname{Hom}(\mathcal H_s,\mathcal H_g)
 \cong{}&3(\tfrac12,0)\oplus3(0,\tfrac12)\\
 &\oplus2(1,\tfrac12)\oplus2(\tfrac12,1)\\
 &\oplus(\tfrac32,0)\oplus(0,\tfrac32).
\end{aligned}}
\]
It has complex dimension $44$ and no trivial or ordinary-vector
$(\tfrac12,\tfrac12)$ summand. The self-adjoint off-block operators
form an $88$-real-dimensional space. In particular the irreducible
\emph{type} of a spinorial input leaves three possible channels per
chirality and does not select an action coefficient.
\end{theorem}
\begin{proof}
Take the tensor product of $\mathcal H_g$ with the dual of each
half-spinor; all $SU(2)$ finite irreducibles are self-dual. The
Clebsch--Gordan rule
$V_j\otimes V_{1/2}=V_{j-1/2}\oplus V_{j+1/2}$
(with the negative term omitted) gives the displayed multiplicities.
The dimension check is $3(2)+3(2)+2(6)+2(6)+4+4=44$.
The Hermitian odd block is uniquely specified by one $11\times4$
complex matrix and its adjoint, thus has $2(44)=88$ real parameters.
\end{proof}

\begin{theorem}[Central-even input obstruction and its sharpness]
\label{p2b:thm:v264-central-even-no-go}
Let $z=(-I,-I)\in\Spin(4)$ and let $V$ be any $\Spin(4)$-space on
which $z$ acts trivially. Every $\Spin(4)$-equivariant map
\[
 F:V\longrightarrow
 \mathcal O:=\operatorname{Hom}(\mathcal H_s,\mathcal H_g)
 \oplus\operatorname{Hom}(\mathcal H_g,\mathcal H_s)
\]
is zero, even if $F$ is nonlinear. Thus wholly central-even
coordinate, tensor, curvature, or bosonic-jet inputs cannot alone
generate a central-odd \emph{one-system} bridge by an equivariant rule.
The obstruction is sharp: for any actual spinorial datum
$0\ne\chi\in\mathcal H_s$ and normalized invariant scalar
$e_0\in\mathcal H_g$, the family
\[
 B_\chi=|e_0\rangle\langle\chi|+|\chi\rangle\langle e_0|
\]
is nonzero, self-adjoint, central-odd and $\Spin(4)$-covariant.
This is mathematical availability only, not physical occurrence.
\end{theorem}
\begin{proof}
The central element acts as $+I$ on $\mathcal H_g$ and $-I$ on
$\mathcal H_s$, hence as $-I$ on $\mathcal O$. Equivariance gives
$F(v)=F(zv)=z\cdot F(v)=-F(v)$. For $B_\chi$, the two terms are
adjoint and exchange the blocks; $\Spin(4)$ fixes $e_0$ and sends
$\chi$ to $\rho(g)\chi$, so the claimed covariance follows.
\end{proof}

\begin{theorem}[No parallel central-odd section for the canonical Born holonomy]
\label{p2b:thm:v264-no-parallel-odd}
On the associated $\mathcal H_1$ representation bundle for the
\emph{already generated correlated fibre-spin connection} of
Paper~III-F, restrict to its established full $SU(2)_{\rm diag}$
holonomy. Then
\[
\mathcal H_g\!\downarrow_{SU(2)_{\rm diag}}
 =2V_0\oplus3V_1,
\qquad
\mathcal H_s\!\downarrow_{SU(2)_{\rm diag}}=2V_{1/2},
\]
\[
 \operatorname{Hom}(\mathcal H_s,\mathcal H_g)
 \!\downarrow_{SU(2)_{\rm diag}}
 =10V_{1/2}\oplus6V_{3/2}.
\]
On any connected domain carrying this holonomy, every parallel odd
section vanishes. This does not rule out nonparallel sections or
create a physical $\mathcal H_1$ interaction bundle from the
fibre connection alone.
\end{theorem}
\begin{proof}
Under the diagonal embedding $(j_L,j_R)\downarrow SU(2)$ the vector
$(\tfrac12,\tfrac12)$ is $V_0\oplus V_1$ and each bivector triplet
is $V_1$. Thus $\mathcal H_g=2V_0\oplus3V_1$;
the chiral doublets are each $V_{1/2}$. Tensoring with two copies
of $V_{1/2}$ gives $2(2V_{1/2}+3(V_{1/2}\oplus V_{3/2}))
=10V_{1/2}\oplus6V_{3/2}$. No trivial summand exists.
A parallel section corresponds to a holonomy-fixed vector, hence is zero.
\end{proof}

\begin{proposition}[Chiral action oddness is not central-oddness]
\label{p2b:prop:v264-chiral-central-distinction}
The rooted nonzero $W_L\leftrightarrow W_R$ Dirac coefficient of
Paper~V-A is odd for its \emph{internal chiral} grading but even
under diagonal-center conjugation, since
$\rho_W(z)=-I_W$ on both halves. Moreover its printed action domain
is $\operatorname{End}(W\otimes Q\otimes G)$, not the mixed block
$\operatorname{Hom}(\mathcal H_s,\mathcal H_g)$.
Therefore its rooted nonzero action incidence by itself does not
discharge the $\mathsf{FENV}$ central-odd bridge or supply a two-carrier
exchange Hessian. This is a type distinction, not a contradiction of
Paper~V-A's nonzero chiral source theorem.
\end{proposition}
\begin{proof}
For every $D_W\in\operatorname{End}(W)$,
$(-I_W)D_W(-I_W)=D_W$; the internal chiral grading need not commute
with $D_W$. The declared source and target spaces of the two
operator owners are distinct, so a nonzero coefficient in the first
is no incidence certificate for the second.
\end{proof}

\begin{theorem}[Basis-free even paired-odd exchange]
\label{p2b:thm:v264-canonical-paired-exchange}
On the \emph{mathematical} two-copy carrier
$\mathcal H_1\otimes\mathcal H_1$, let $F(u\otimes v)=v\otimes u$
be the tensor flip, $P_g,P_s$ the central rank-$11$ and rank-$4$
projectors, and put
\[
 Q_{\times}=P_g\otimes P_s+P_s\otimes P_g,
 \qquad K=FQ_{\times}.
\]
Then
\[
 K^*=K,\quad K^2=Q_{\times},\quad
 \operatorname{rank}K=88,\quad
 [K,\rho(g)\otimes\rho(g)]=0,
\]
\[
 [K,\Gamma_c\otimes\Gamma_c]=0,
 \qquad [K,P_g\otimes I]\ne0.
\]
Thus a symmetry-invariant \emph{total-even} operator can change an
individual carrier's $11/4$ population. It is not in the native
$\mathcal A_{\rm ev}\otimes\mathcal A_{\rm ev}$ test algebra, and
the mathematical tensor construction alone does not actualize a
second carrier, its identification, or an interaction/process token.
The exact conditional population evolution is proved at its
physical-owner boundary in Paper~IV-A.
\end{theorem}
\begin{proof}
The flip commutes with the symmetric crossed projector $Q_{\times}$;
both are self-adjoint and $F^2=I$. Therefore $K^*=K$,
$K^2=Q_{\times}$ and its rank is twice $11\cdot4=44$.
The diagonal symmetry and total central parity commute with both
factors. On a vector $u_g\otimes v_s$, $K$ exchanges it with
$v_s\otimes u_g$, so it does not commute with the first system's
$P_g\otimes I$. Every native even tensor-product test does commute
with $P_g\otimes I$, hence $K$ is not one of them.
\end{proof}

\begin{remark}[Non-converses preserved by this reverse audit]
\label{p2b:rem:v264-converse-firewall}
The derived carrier implies $m_p=1$, not conversely: for example
$\mathcal E_9=(0,0)\oplus(\tfrac12,\tfrac12)\oplus
(\tfrac12,0)\oplus(0,\tfrac12)$ also has maximal half-level one.
Likewise $\mathcal H_1$ can be produced abstractly by different
mathematical tensor-prefix/lower-ideal grammars and hence does not
identify its \emph{physical} bracket ancestry. The $15/137$ algebra
and exact quadratic fingerprint are independent reconstruction tests,
not G0 provenance certificates. Nor does a covariant mathematical
exchange imply an action-owned $\mathsf{PROC}$, $\mathsf{DYN}$,
$\mathsf{FENV}$ or $\mathsf{ENV}$ occurrence.
\end{remark}


\section{Exceptional finite-support comparison and nonidentification}
\label{p2b:sec:v318-petersen-momentum-firewall}

\begin{proposition}[No dimension-based carrier identification]
\label{p2b:prop:v318-petersen-vs-h1}
On the conditional embedded $2I$ support of Paper~II-A, the real
Petersen-edge module $W_P\cong\mathbb R^{15}$ has an $A_5$ action and
a finite decomposition
\[
 W_P\cong\mathbf4\oplus\mathbf5\oplus\mathbf3\oplus\mathbf3'.
\]
It is not the independently generated fifteen-complex-dimensional
round-host momentum carrier
\[
 \mathcal H_1\cong\mathbf1\oplus S\oplus\mathfrak c_{\mathbb C}
 \quad(15=1+4+10)
\]
and it is not the distinct fifteen-dimensional Bose--Mesner block
inside $\mathbb C^{2T}$.  The central binary element $-1\in2I$ acts
trivially on $W_P$; the signed Hopf-axis and quaternionic-fibre
lifts must be retained separately if spinorial central parity is to
be detected.  Neither matching dimension nor restriction of a
continuous representation to a finite subgroup supplies a canonical
intertwining map between these three carriers.
\end{proposition}
\begin{proof}
The three spaces are defined on different source objects and with
different actions.  Paper~I-B gives the Petersen representation, and
the round-parent generator/gradient/field-bracket construction here
gives $\mathcal H_1$ independently of any finite subgroup selection.
The $A_5$ permutation action factors through the quotient by the
binary center.  Without a separately owned action-intertwiner, no
identification is licensed by GP2's exact/faithful requirement.
\end{proof}

\begin{remark}[No changed momentum frontier]
\label{p2b:rem:v318-petersen-momentum-scope}
The intrinsic result $\mathcal H_p\cong\mathcal H_1$ and the derived
label $m_p=1$, the coordinate-completion moduli and all existing
normalized momentum scales remain as before.  The Petersen packet is
an extra finite comparison structure on a separately constituted
exceptional support role, not a new level-one momentum constructor,
component ceiling or physical Hilbert-space multiplicity.
\end{remark}


\section{Conclusions and export boundary}
\subsection{Conclusion}
The inverse mathematical certificates of Section~\ref{p2b:sec:v264-reverse-certificates} sharpen the already generated carrier and its observability boundary without altering the frozen GP2 frontier or discharging the separate action/physical-occurrence obligations.

Paper~II-B completes the constructor-relative momentum constitution required by
Paper~II-A.  The single momentum parent is
\(SU(2)\cong S^3(\Rmom)\); its connected isometry completion is \(SO(4)\),
with derived universal spin lift \(\Spin(4)\).  The finite carrier is not
selected by a numerical cutoff or by minimizing a word-stage label.

The round host supplies the intrinsic lowest spinor objects
\(S_+\cong V_{\frac12,0}\) and \(S_-\cong V_{0,\frac12}\), together with
\(E_0\), \(E_1\cong V_{\frac12,\frac12}\), and the canonical bracket linear image
\[
 \begin{aligned}
  Z_\beta
  &:=\operatorname{span}_{\mathbb R}\operatorname{im}
  \bigl(\beta:E_1^{\mathbb R}\times E_1^{\mathbb R}
      \to\Gamma(TS^3)_{\rm sec,obj}\bigr)\\
  &=\operatorname{im}\bar\beta=\mathcal K_{\rm Kill},\\
  (\mathcal K_{\rm Kill})_{\mathbb C}
  &\cong(1,0)\oplus(0,1).
 \end{aligned}
\]
The literal bilinear set image is only the decomposable locus; the equality is
for its generated real linear hull.  Paper~II-A finite generated-image
inheritance applies to this finite-dimensional section-object subspace, and
its same-role no-splitting lemma places the inherited complexified object copy
in the momentum carrier because the frozen seed frontier generates no
autonomous role split.
Therefore GP2 exactness and host faithfulness fix the carrier type-support
\[
 E_0\oplus S_+\oplus S_-\oplus E_1
 \oplus(\mathcal K_{\rm Kill})_{\mathbb C}
 =\mathcal H_1.
\]
Equivalently, the six irreducible summands have dimensions
\(1,2,2,4,3,3\), so
\[
 \mathcal H_1\cong
 V_{0,0}\oplus V_{\frac12,0}\oplus V_{0,\frac12}
 \oplus V_{\frac12,\frac12}\oplus V_{1,0}\oplus V_{0,1},
 \qquad
 \dim_{\mathbb C}\mathcal H_1=15.
\]
This records the sectorwise dimension check inside the conclusion; it is not
a second theorem independent of the host derivation.
The canonical smooth-model theorem of Paper~I-B supplies the retained
presentation
\[
 \mathcal L_R
 =\operatorname{span}\{V_R^{-1/2}\}
 \oplus\operatorname{span}\!\left\{
   \frac{2}{\sqrt{3V_R}}dx_a\right\}
 \oplus\operatorname{span}\!\left\{
   \sqrt{\frac2{V_R}}dx_a\wedge dx_b\right\}
 \oplus S_+\oplus S_-,
\]
The smooth-model theorem identifies the explicit one-copy GP2 type-support
with that upstream model by \(J_R\), without invoking universal GP3
rigidity.  It is a natural same-parent presentation of the generated
constituents and needs no second occurrence witness for those constituents.  It does not select a
coordinate completion, ambient lift, or kinetic operator.
These five generated sectors co-realize as constituents of the same occurring
momentum composite; they are not competing branches.  No-omission excludes the
former \(\mathcal E_9\)-type countermodels, while constructor-relative
no-insertion prevents arbitrary tensor products or coordinate polynomial
products from generating higher carrier sectors.  GP3 removes unmarked implementation/multiplicity redundancy where the
whole-structured comparison certificate is proved.  Only under the additional
hypotheses of Theorem~\ref{p2b:thm:momentum-carrier-rigidity} does it follow
that
\[
 \forall\mathfrak M\in\mathcal M_p^{\rm GP3},\qquad
 \mathcal H_{\mathfrak M}\cong\mathcal H_1,
\]
without any assertion that \(|\mathcal M_p^{\rm GP3}|=1\).

Independently of the GP3 comparison, the representation-theoretic word
length is read off from the GP2 type-support.  Since the terminal sectors of \(\mathcal H_1\) have
\(\ell=2\), the structural momentum-role label is
\[
 \boxed{m_p=1}.
\]
Thus the strongest verified type-support chronology is
\[
 \boxed{\mathcal G^{\rm car}_{\Sigma_p}\cong\mathcal H_1\Longrightarrow m_p=1,}
\]
not the reverse.  Equivalently, the GP2-generated momentum support obeys
\[
 \boxed{
 \mathcal G^{\rm car}_{\Sigma_p}\cong\mathcal H_1
 \quad\text{independently of}\quad
 \mathsf P_{\rm sli},\ \mathrm{TWR},\ \mathrm{SPEX},\ \mathrm{FIT},\
 \mathsf P_{\rm lvl},\ P_{\rm ceil},\ \mathsf C_{\rm PI},}
\]
where \(\mathsf C_{\rm PI}\) is relevant only to later completion-specific
coordinate statements.

The generic family \(\mathcal H_m\), optional lower-ideal and component-ceiling
grammars, shell filtration, invariant-state classification, and Peter--Weyl
parent-symbol recovery remain available mathematics and later-role
infrastructure.  They no longer explain the initial momentum seed.
The coordinate-completion family and the point-reference role remain
logically separate: neither a completion point nor a point-reference readout
token follows from \(\mathcal H_1\) or the integer equation \(m_p=1\).  Raw
ambient coordinate actions remain parameterized by
\(\mathfrak M_{\rm coord}(\mathcal H_1)/\!\sim\cong\mathbb R_{\ge0}^3\),
although current Paper~II-C later proves the normalized protected even
\(\Cl(0,3)\) algebra completion-independent on the nondegenerate protected
stratum.  An occurring point-resolution readout generates only
\(\mathsf{Supp}_{\varepsilon}(X)\); the exceptional \([2T]\) support
co-realizes only with the additional II-A core constructor/domain bridge or a
direct core-constructor occurrence.  Likewise the carrier label fixes no
numerical \(\Rmom\), no equality \(\rho_{\rm prot}=\Rmom\), no physical
kinetic Hamiltonian, and no measurement or calibration token.

G0 supplies the common semantic boundary.  Its original transfer-operation
occurrence co-realizes the parent and the generated seed constituent family,
without a second G0 momentum witness.  Universality of the multiplicity-free
physical endpoint still requires the GP3 structured comparison certificate.  A later named test, process, trajectory,
bridge, readout, or calibration that is not already a generated constituent
requires its own stage-safe, result-independent token witness.  Such a token
does not reopen the carrier theorem.


\printbibliography
\end{document}
